\documentclass[11pt]{article}
\usepackage[T1]{fontenc}
\usepackage{lmodern,microtype,amsmath,amssymb,amsthm,booktabs,array}
\usepackage[margin=1in]{geometry}
\usepackage[hidelinks]{hyperref}
\hypersetup{pdftitle={Sharp subsidies for truthful production and network cooperation},pdfauthor={Maxwell Flanagan}}
\newtheorem{theorem}{Theorem}[section]
\newtheorem{lemma}[theorem]{Lemma}
\newtheorem{proposition}[theorem]{Proposition}
\newtheorem{corollary}[theorem]{Corollary}
\theoremstyle{definition}\newtheorem{definition}[theorem]{Definition}
\theoremstyle{remark}\newtheorem{remark}[theorem]{Remark}
\newcommand{\R}{\mathbb R}
\newcommand{\E}{\mathbb E}
\newcommand{\Prob}{\mathbb P}
\newcommand{\core}{\mathcal C}
\newcommand{\pos}[1]{(#1)_+}
\newcommand{\ind}{\mathbf 1}
\title{Online Resource 1: Indexed supplement\\to Sharp subsidies for truthful production\\and network cooperation}
\author{Maxwell Flanagan}
\date{17 October 2026\\[2pt]\small Economic Theory candidate; not submitted}
\begin{document}
\maketitle
\noindent\textbf{Author information.} The author has no affiliation.
Correspondence details require confirmation before submission.

\section*{Scope and provenance}
This supplement preserves the mathematical sections of the 7 October
consolidated research manuscript and adds the audited 16 October owner-count,
dimension and restricted-architecture results. Sections~\ref{general:section}
and~\ref{arch:section} are new; the earlier square-root-free upper estimates
are retained as valid historical intermediate bounds where explicitly marked. The focused main paper contains the complete
general upper and simple-DAG lower proofs. Sections~\ref{prod:section}
and~\ref{dag:section} repeat those foundations to keep every subsequent
reference local and to make this file independently readable. This repetition
is not an additional contribution.

The original source is retained unchanged as
\texttt{research\_record/manuscript\_2026-10-07.tex}. Its SHA-256 is
\begin{center}\small\texttt{4786126d3a4473bf04b06384ad7151a83576}\\
\texttt{321c4734718f26f29ba6ff13e67f}\end{center}
The section index below maps mathematical content to that record. Existing
proof labels are preserved. The source-applicability paragraph for the strictly
concave projection diagnostic and the Br\^anzei--Solymosi--Tijs bibliography
are updated using newly inspected originals. Two audited corrections are
also explicit: the uncapped-rule example now specifies the owner bundles
and uses a two-owner Shapley alternative valid on the full bundle domain;
the cut-budget inversion specifies $0<\varepsilon<1$. The sharp complexity upper and lower statements are unchanged. The new
general bounds improve the owner-count and concave-dimension estimates;
the new matching barriers restrict a specified mechanism architecture. Online Resource 2 records fresh checks for the
central theorem, projection diagnostic, and the 16 October owner/dimension
checks; the October consolidation ledger
records the scoped prior audits of unchanged supplementary results.

\paragraph{Index.}
Sections~\ref{sec:model}--\ref{sec:escape} develop the weighted multi-glove
model, sharp hierarchy, proportional rules and restricted nonlinear escape.
Section~\ref{general:section} gives the new square-root upper bounds and
Section~\ref{arch:section} proves the restricted mechanism barriers.
Section~\ref{ext:sec} gives parameter lower bounds, general
concave homogeneous production and optional original-core protection.
Section~\ref{alg:section} gives finite sampling and all-profile certification,
with potentially exponential preprocessing. Section~\ref{sec:remaining}
states remaining questions. The appendices retain finite-pool constructions,
network and noise mechanisms, envelope obstructions, projection diagnostics,
and exact failures of alternative rules. The full contents provides page
locators. These are supplementary results within one research program,
not a list of separately validated papers.

\paragraph{Review boundary.}
Mathematical arguments and exact finite checks received internal model audits.
The checks are not formal verification, and they do not prove asymptotic
claims. No human specialist endorsement or comprehensive priority clearance
is claimed. The author must confirm responsibility and declarations before
submission. This local preparation does not publish or submit any material.
\tableofcontents
\clearpage
\section{General production: a sharp complexity bound}\label{prod:section}
There are $n$ fixed owners and $d$ fixed resource types. Owner $i$ reports
$b_i\in\R_+^d$; write $c=\sum_i b_i$. Fix a nonempty finite family
$V=\{w^1,\ldots,w^K\}\subseteq\R_+^d$, with a report-independent priority,
and define
\[
 f(c)=\min_{w\in V}w\cdot c,\qquad
 v_b(S)=f\left(\sum_{i\in S}b_i\right),\qquad T=f(c).
\]
A rule is eligible if its payments are nonnegative, cover every coalition,
and are componentwise nondecreasing in each owner's whole bundle, with all
other reports fixed. These conditions hold on all of $(\R_+^d)^n$.
They express truthfulness for binding underreports and payment-only utility;
there are no operating costs, overreport, false-name, population-change or
coordinated-report guarantees. A factor $C$ bounds \emph{total payment} by
$CT$; external subsidy is at most $(C-1)T$. The family $V$ remains fixed
across reports. Its size may be exponential in a compact description. All logarithms are natural.

\begin{theorem}[General sharp-order upper bound]\label{prod:upper}
For every integer $j\ge2$, there exists an eligible deterministic rule with factor
\begin{equation}\label{prod:constant}
 C_j(K)=j+\frac{\log K+2\log j+1}{\alpha_j},\qquad
 \alpha_j=\log\frac{j^3}{j^2+j-1}>0.
\end{equation}
With $L=\log K$ and
$j=\max\{2,\lceil L/[\log(2+L)]^2\rceil\}$, this factor is
$(1+o(1))\log K/\log\log K$, uniformly over owner counts, dimensions,
technologies and reports. Zero-value profiles receive zero; for $K=1$ the
constructed rule is efficient.
\end{theorem}

\paragraph{The mechanism.}
For $q\in(0,1)$ let independent owner bits satisfy $\Prob(B_i=0)=q$.
Choose the first minimizer of the \emph{uncapped} objective,
\[
 W(b,B)=\operatorname*{first\,argmin}_{w\in V}
               \sum_i B_i(w\cdot b_i),\qquad
 a_i(w,b)=\min\{w\cdot b_i,T\}.
\]
Only the payment contribution $a_i$ is capped. Put
\begin{align}
 z_i^q&=\E a_i(W,b),&
 p_i^q&=\E[a_i(W,b)\mid B_i=0],\label{prod:p}\\
 m_i^q&=\E_{B_{-i}}\left[f\left(b_i+\sum_{h\ne i}B_hb_h\right)
                         -f\left(\sum_{h\ne i}B_hb_h\right)\right],
 &r_i^q&=qp_i^q+(1-q)m_i^q.\label{prod:hybrid}
\end{align}
The final payment is the deterministic finite expectation
\begin{equation}\label{prod:payment}
 M_i(b)=j\int_0^{1/j}r_i^q(b)\,dq.
\end{equation}
These quantities are finite polynomials in $q$ at a fixed profile; no
continuity or generic-position assumption about the reports is required.

\begin{proof}[Proof of Theorem~\ref{prod:upper}]
Conditional on $B_i=0$, every score and the selected form are independent of
$b_i$. For that fixed nonnegative form, both $w\cdot b_i$ and $T$ are
nondecreasing in $b_i$, so their minimum is nondecreasing. In $m_i^q$ the
off term is independent of $b_i$ and the on term is nondecreasing. The
nonnegative fixed mixtures preserve this property and nonnegativity.

For any $w,S$, if some $i\in S$ has $w\cdot b_i\ge T$, then
$\sum_{i\in S}a_i(w,b)\ge T\ge v_b(S)$. Otherwise no term in $S$ is capped,
and the sum is $w\cdot\sum_{i\in S}b_i\ge v_b(S)$. Thus every capped form
covers coalitions. Fix $B_{-i}$ and let $w_1$ be the on minimizer. The off
optimum is at most the off objective of this same form, whence
\begin{equation}\label{prod:endpoint}
 f\left(b_i+\sum_{h\ne i}B_hb_h\right)-f\left(\sum_{h\ne i}B_hb_h\right)
 \ge w_1\cdot b_i=a_i(w_1,b).
\end{equation}
The equality holds because the contribution is a nonnegative summand of
the selected on value, which is at most $T$. Averaging proves $r_i^q\ge z_i^q$,
and hence coverage for $r^q$ and $M$.

For the budget define $g(q)=\E f(\sum_iB_ib_i)$, $Z(q)=\sum_i z_i^q$,
and $R(q)=\sum_i r_i^q$. Capping leaves every open contribution unchanged:
\[
 \sum_iB_i a_i(W,b)=f\left(\sum_iB_ib_i\right).
\]
Consequently $q\sum_i p_i^q=Z-g$. Differentiating the finite mask polynomial
gives $\sum_i m_i^q=-g'$, so $R=Z-g-(1-q)g'$. Since $g(0)=T$ and $g\ge0$,
for $0<\eta<1$ integration by parts gives the exact identity
\begin{equation}\label{prod:integral}
 \int_0^\eta R\,dq=\int_0^\eta Z\,dq+T-(1-\eta)g(\eta)
                         -2\int_0^\eta g\,dq
 \le\int_0^\eta Z\,dq+T.
\end{equation}
Both the sign and the factor two in this identity matter.

Suppose $T>0$. For each fixed form let
$d_i(w)=a_i(w,b)/T\in[0,1]$ and $\rho_w=\sum_i d_i(w)$.
The selected form satisfies $\sum_iB_id_i(W)\le1$. For $\lambda>0$ put
$\alpha=-\log(q+(1-q)e^{-\lambda})$. The convex function
$u\mapsto\log(q+(1-q)e^{-\lambda u})$ lies below its endpoint chord
$-\alpha u$ on $[0,1]$. Independence and Markov's inequality give
\[
 \Prob\left(\sum_iB_id_i(w)\le1\right)
 \le e^\lambda\prod_i(q+(1-q)e^{-\lambda d_i(w)})
 \le e^{\lambda-\alpha\rho_w}.
\]
A union over the fixed family and integration of the resulting tail yield
\begin{equation}\label{prod:chernoff}
 \Prob(\rho_W\ge x)\le\min\{1,K e^{\lambda-\alpha x}\},\qquad
 \frac{Z(q)}T\le\frac{\log K+\lambda+1}{\alpha}.
\end{equation}
Take $\eta=1/j$ and $\lambda=2\log j$. For $q\le1/j$,
$q+(1-q)e^{-\lambda}\le(j^2+j-1)/j^3$, so
\eqref{prod:integral}--\eqref{prod:chernoff} give \eqref{prod:constant}.
For the displayed choice of $j$, as $L\to\infty$,
$j\sim L/(\log L)^2$, $\alpha_j\sim\log L$, and
$j=o(L/\log L)$, proving the asserted asymptotic.

If $T=0$, all masked production values and capped contributions vanish.
If $K=1$, no contribution is capped, $p_i^q=m_i^q=w\cdot b_i$, and
$M_i=w\cdot b_i$. Finally each finite-mask term is
$q^{n-1-|S|}(1-q)^{|S|}[q a_i(W(b,1_S),b)+(1-q)(v_b(S\cup\{i\})-v_b(S))]$.
Its integral is a finite sum with rational coefficients when $j$ is integer.
This gives exact rational payments for rational input, without an efficiency
claim about enumeration of masks or forms.
\end{proof}

\begin{corollary}[Production LPs and networks]\label{prod:lp}
The theorem applies to
$f(c)=\max\{\pi\cdot x:x\ge0,Ax\le c\}$ for fixed $A\ge0$, with every
activity column nonzero and fixed positive profits $\pi$. Here $V$ can be
the vertices of $\{w\ge0:A^Tw\ge\pi\}$. For single-commodity directed flow,
one may instead use the characteristic vectors of the inclusion-minimal
structural source--sink edge cuts.
\end{corollary}
\begin{proof}
Each activity is bounded by a positive entry in its column, so every feasible
production set is compact at finite $c$. The dual is nonempty, pointed and
polyhedral. Strong duality and attainment at a vertex give the claimed finite
min-linear representation, including boundary capacities. Network max-flow/
min-cut gives the cut representation directly. Redundant nonminimal cuts may
be removed once and for all because capacities are nonnegative. All owners
and forms remain fixed across profiles. A positive-profit zero-resource
activity is excluded because it would make production unbounded.
\end{proof}
Dual core allocations are established ingredients~\cite{Owen,SametZemel}.
General dual vertices may be fractional; their count is not an ordinary
network-cut count. Constraints introducing capacity-independent dual offsets
are outside the homogeneous model here.

\begin{proposition}[Why the payment cap is essential]\label{prod:uncapped}
Removing the cap from the own-excluded contribution in
\eqref{prod:p} makes the uniform-$q$ hybrid unbounded even with two owners
and two supporting forms.
\end{proposition}
\begin{proof}
Take unit-profit activity columns $(1,1)$ and $(M,0)$, $M>1$, and two
fixed owners with bundles $b_1=(M,0)$ and $b_2=(0,1)$, so $c=(M,1)$. The two dual vertices are $(1,0)$ and $(1/M,1-1/M)$, and
$T=2-1/M$. If owner 1 is off and owner 2 on, $(1,0)$ uniquely minimizes the
raw score, giving uncapped contribution $M$. Thus $r_1^q\ge q(1-q)M$;
with $j=2$ its average is at least $M/6$, an unbounded ratio to $T$.
This obstructs only the uncapped rule. On the full two-owner bundle
domain, the Shapley payments
$P_i(b)=[f(b_i)+T-f(b_{-i})]/2$ are efficient and own-monotone.
They are nonnegative by monotonicity of $f$, and cover singleton coalitions
because $f(b_1+b_2)\ge f(b_1)+f(b_2)$. This last inequality follows by
minimizing nonnegative linear forms. Thus an efficient eligible alternative
exists without restricting each owner's report to a single resource.
\end{proof}

\section{A matching lower bound on simple networks}\label{dag:section}
A network has one distinct fixed strategic owner per edge, and coalition
worth is its maximum source--sink flow using only its owned edges. Structural
cuts are inclusion-minimal edge sets disconnecting source from sink; their
number depends on the fixed graph, not its reported capacities. Define
$\mathcal B_{\rm dag}(K)$ as the supremum, over simple finite DAGs having at
most $K$ structural cuts, of the best all-profile eligible total-payment
factor on that graph. The analogous quantity for finite nonnegative
min-linear technologies with at most $K$ forms is $\mathcal B_{\rm prod}(K)$;
owner counts and dimensions are unrestricted finite integers.

\begin{theorem}[Sharp first-order worst-case factors]\label{dag:sharp}
As $K\to\infty$,
\[
 \mathcal B_{\rm dag}(K)=(1+o(1))\frac{\log K}{\log\log K},\qquad
 \mathcal B_{\rm prod}(K)=(1+o(1))\frac{\log K}{\log\log K}.
\]
The lower holds along an explicit sequence of simple DAGs at their exact cut
counts, and for the worst case with \emph{at most} $K$ cuts. It does not
assert an exact-count padding theorem for every integer $K$.
\end{theorem}

\begin{proof}
The upper follows from Theorem~\ref{prod:upper}. We give the simple-DAG lower,
which also supplies the production lower. Fix integers $d_0\ge1$, $m>d_0$,
$R\ge1$, and $k\ge2$. (Here $d_0$ is hierarchy height, not resource dimension.)
One module with terminals $a,b$ has vertices $L_1,\ldots,L_R$,
$u_1,\ldots,u_m$, $v_1,\ldots,v_m$, $W_1,\ldots,W_R$ and edges
\[
 aL_s,\quad L_su_j,\quad u_jv_j,\quad v_jW_t,\quad W_tb
 \quad(1\le s,t\le R,\ 1\le j\le m).
\]
Take $k$ copies in series by identifying consecutive terminals. Each middle
edge $u_jv_j$ has report in $\{1,m\}$; all auxiliary edges have report
$H=m^2$. Every edge, including every auxiliary edge, is a strategic original
owner whose full report domain is nonnegative. Only the finite profiles
used in this proof keep auxiliary reports fixed.

A type-$h$ module has $h$ high middle edges, $0\le h\le d_0$. For any subset
$J$ of middle edges, choose one left fan and one right fan, and include the
source/final fan edges plus the incident edges supporting exactly $J$.
This actual coalition has module flow $\sum_{j\in J}c_j$: the middle cut
bounds it above, and the selected branches carry it because their total is
at most $m^2$. In particular the full graph has
\begin{equation}\label{dag:value}
 T=m+(m-1)\min_{1\le l\le k}H_l.
\end{equation}
Choosing the two fans independently uniformly, any particular auxiliary
edge is included with probability at most $1/R$.

Symmetrize an arbitrary eligible factor-$B$ rule over module permutations
and branch permutations, relabeling each branch's incident auxiliary edges.
These are permutations of the fixed owner set preserving every coalition's
flow at all profiles: module flows combine by their minimum, so even module
permutations preserve coalition values. Averaging preserves all required
inequalities. At a histogram $\nu=(\nu_0,\ldots,\nu_{d_0})$ of module types,
write $x_h(\nu),z_h(\nu)$ for low and high variable-owner payments, and
\[
 Y(\nu)=\sum_h\nu_h[(m-h)x_h(\nu)+hz_h(\nu)],\qquad
 Q(\nu)=\text{total auxiliary-owner payment}.
\]
No auxiliary-payment symmetry is assumed; $Y+Q\le BT$.

First choose one low variable in each module, supported by random fans.
Averaged coalition coverage and nonnegativity give
$\sum_h\nu_hx_h+Q/R\ge1$. Next put $q_0=m-d_0+1$ and choose $q_0$ low
variables in each type below $d_0$, and one high variable in each type
$d_0$. These exist because $m>d_0$. The resulting coalition has worth at
least $q_0$, giving
$q_0\sum_{h<d_0}\nu_hx_h+\nu_{d_0}z_{d_0}+Q/R\ge q_0$.
Multiply the first inequality by $m-d_0$ and the second by $d_0$. With
$A=(d_0+1)m-d_0^2$,
\begin{equation}\label{dag:coalitions}
 A\le A\sum_{h<d_0}\nu_hx_h+
 \nu_{d_0}[(m-d_0)x_{d_0}+d_0z_{d_0}]+(m/R)Q.
\end{equation}

Fix $\alpha>0$, set $\lambda=A/\alpha$,
$Z=\sum_{h=0}^{d_0}\lambda^h/h!$, and $p_h=\lambda^h/(h!Z)$.
Average profiles with independent module types of probabilities $p_h$;
let $\mu(\nu)$ be their multinomial histogram weights. Raising one actual
low variable owner from $1$ to $m$ gives
$x_h(\nu)\le z_{h+1}(\nu')$ at
$\nu'=\nu-e_h+e_{h+1}$. Every other report, including every auxiliary report,
is unchanged. The exact identity
\[
 \mu(\nu)\nu_h=(p_h/p_{h+1})\mu(\nu')\nu'_{h+1},\qquad
 A p_h/p_{h+1}=\alpha(h+1)
\]
transports only this changing owner's payment. Averaging
\eqref{dag:coalitions} therefore yields
\begin{align}
 A&\le\alpha\E\sum_{h=1}^{d_0}h\nu_hz_h+
 \E\nu_{d_0}[(m-d_0)x_{d_0}+d_0z_{d_0}]+(m/R)\E Q\notag\\
 &\le(1+\alpha)\E Y+(m/R)\E Q
 \le\max\{1+\alpha,m/R\}B\E T.\label{dag:transport}
\end{align}
The last step uses the disjoint budgets $Y,Q$; their coefficients are not
added. By the tail-sum formula and $p_0=1/Z$,
\begin{equation}\label{dag:mean}
 \E T=m+(m-1)\sum_{r=1}^{d_0}\left(\sum_{h=r}^{d_0}p_h\right)^k
 \le m[1+d_0 e^{-k/Z}].
\end{equation}

Let $C_{m,R}$ be the number of structural cuts of one module. Any minimal
cut is the outgoing cut of its reachable source side after deletion, so
$1\le C_{m,R}\le2^{2R+2m}$. A minimal cut in a series composition lies
entirely in one module: otherwise each module retains a path, or an edge
removed elsewhere is redundant. Hence the full cut count is exactly
$K=kC_{m,R}$.
For $d_0\ge2$ choose
\[
 g=\lceil\sqrt{\log(d_0+1)}\rceil,\quad
 m=d_0g,\quad R=m,\quad \alpha=1/g,\quad
 \ell=\lceil\log(d_0+1)\rceil,\quad k=\lceil2\ell Z\rceil.
\]
The separate multiplier $\ell$, rather than $g$, ensures the tail error
vanishes. Equations~\eqref{dag:transport}--\eqref{dag:mean} give
\[
 B\ge\frac{d_0+1-d_0/g}{(1+1/g)(1+d_0e^{-2\ell})}=(1-o(1))d_0.
\]
Here $\lambda=(d_0+1)d_0g^2-d_0^2g>d_0$, so
$\lambda^{d_0}/d_0!\le Z\le(d_0+1)\lambda^{d_0}/d_0!$.
Stirling's estimate implies
$\log k=(1+o(1))d_0\log d_0$; since $m=o(d_0\log d_0)$,
\[
 \log k\le\log K\le\log k+4m\log2=(1+o(1))d_0\log d_0.
\]
This proves the lower along the sequence $K_{d_0}$. For any fixed
$0<\varepsilon<1$ and large external cut budget $K$, take
$d_0=\lfloor(1-\varepsilon)\log K/\log\log K\rfloor$.
Then $K_{d_0}\le K$ eventually, yielding lower coefficient $1-\varepsilon$.
Letting $\varepsilon$ decrease to zero proves the at-most-$K$ assertion.
\end{proof}

\begin{proposition}[A subdivision that destroys the lower bound]\label{dag:subdivision}
Fully subdividing every edge of any finite directed network into two edges
owned separately produces a network with an efficient eligible rule.
\end{proposition}
\begin{proof}
For pair reports $(a_e,b_e)$ set $d_e=\min(a_e,b_e)$. Choose the original
network's fixed-priority minimum-cut core allocation $Q(d)$, which pays
$d_e$ on the selected cut and zero elsewhere. Pay $Q_e(d)$ to the higher
reporting member of pair $e$, split equally at equality, and pay the lower
member zero. Holding the mate's report fixed, the payment is zero below it,
half the fixed $Q_e$ at equality and the fixed $Q_e$ above it: in the last
two cases all virtual capacities are fixed. Thus it is own-monotone, even
if $Q$ is not. For a coalition $S$, let $S^*$ contain the original edges
whose two members both lie in $S$. Its value is exactly $v_d(S^*)$, and
$p(S)\ge Q(S^*)\ge v_d(S^*)$. Total payment is $T$. This shows why the naive subdivision reduction fails and motivates the
explicit auxiliary-budget accounting in the shared-fan construction.
\end{proof}

\section{Model and two reductions}\label{sec:model}
Let $N=N_1\sqcup\cdots\sqcup N_k$ be a fixed finite partition with
$k\ge2$ and nonempty pools. Reports are $c\in\R_+^N$. Write
\[
 C_j=\sum_{i\in N_j}c_i,\qquad
 v_c(S)=\min_j\sum_{i\in S\cap N_j}c_i,\qquad T=\min_j C_j.
\]
A payment rule $y(c)\in\R_+^N$ \emph{covers coalitions} if
$y(c)(S)\ge v_c(S)$ for every $S\subseteq N$. It is \emph{own-report
monotone} if $y_i(t,c_{-i})$ is nondecreasing in $t$. Under binding
underreports and payment utility, this is exactly the required no-profitable-
underreport condition. We impose no overreport, group-report, split, merge,
or population-change guarantee. On a restricted domain, all comparisons
whose endpoints belong to that domain are required.

The original efficient core and coordinate benchmarks are
\[
 \core(c)=\{x:x(N)=T,\ x(S)\ge v_c(S)\ \forall S\subseteq N\},
 \qquad H_i(c)=\max_{x\in\core(c)}x_i.
\]
Paying capacities in one minimum-total pool and zero elsewhere gives a core
allocation. Singleton constraints imply $x_i\ge0$, so this nonempty core is
compact. Our optional fairness requirement is $y_i(c)\ge H_i(c)/k$.
A factor $B$ means $y(c)(N)\le BT$ at every profile in the stated domain.
At $T=0$ we set payments to zero.

Let $\mathfrak P_k$ be the set of finite ownership partitions with $k$
nonempty pools, and $\mathcal Y(P)$ the nonnegative covering own-report
monotone rules on all of $\R_+^N$ for partition $P$. Define the floor-free
unrestricted minimax factor
\begin{equation}\label{eq:Bk}
 B_k=\sup_{P\in\mathfrak P_k}\ \inf_{y\in\mathcal Y(P)}\quad
       \sup_{c:T(c)>0}\frac{y(c)(N)}{T(c)}.
\end{equation}
Define $B_k^{\mathrm{floor}}$ by the same formula with the infimum
restricted to rules also satisfying $y_i(c)\ge H_i(c)/k$ at every profile.
Rules for different ownerships need not obey any consistency requirement.
The supremum permits both pool sizes and the number of pools to grow in
asymptotic statements. All logarithms in this paper are natural.

\begin{lemma}[Marginal floor]\label{lem:floor}
Put $M_i=T-v_c(N\setminus\{i\})$. Then $H_i\le M_i$, $M_i$ is own-report
monotone, and $\sum_iM_i\le kT$. If a nonnegative covering own-monotone base
has factor $B$, replacing $y_i$ by $\max\{y_i,M_i/k\}$ gives the original
floor with factor at most $B+1$.
\end{lemma}
\begin{proof}
Core efficiency and the complementary coalition constraint give
$x_i=T-x(N\setminus\{i\})\le M_i$. If $i\in N_j$, set
$A_i=C_j-c_i$ and $D_j=\min_{\ell\ne j}C_\ell$. Then
\[
 M_i=\min\{c_i,(D_j-A_i)_+\}=(c_i-(C_j-T))_+.
\]
The first form proves own-monotonicity. For $s=C_j-T\ge0$, if $p\ge1$
terms $(c_i-s)_+$ in pool $j$ are positive, their sum is at most
$C_j-ps\le C_j-s=T$; if none is positive it is zero. Sum over pools.
The maximum of two own-monotone functions is own-monotone, and its additional
cost is at most $\sum_iM_i/k\le T$.
\end{proof}
Consequently the existence of a universal constant factor with this floor
is equivalent to its existence without the floor, although the best constants
need not coincide. In the notation of \eqref{eq:Bk},
$B_k\le B_k^{\mathrm{floor}}\le B_k+1$.

\begin{lemma}[Capping]\label{lem:cap}
Let $d_i=\min\{c_i,T\}$ and $D_j=\sum_{i\in N_j}d_i$. Then
\[
 v_c(S)=\min_j\sum_{i\in S\cap N_j}d_i,\qquad \min_jD_j=T.
\]
\end{lemma}
\begin{proof}
For any owner set $A$, $\min\{c(A),T\}=\min\{d(A),T\}$: if it contains
an owner above $T$, both truncated sums are $T$; otherwise the sums agree.
Coalition values are at most $T$. A minimum-total original pool has no
capacity above $T$, and has capped total $T$. Taking minima proves the claim.
\end{proof}

\section{An all-profile upper bound}\label{sec:banzhaf}
For $T>0$, draw independent $U_i\sim\operatorname{Unif}[0,1]$ as an
integration device and define
\[
 X_j=\sum_{i\in N_j}d_iU_i,\qquad
 p_j=\Prob(X_j<X_\ell\text{ for all }\ell\ne j),\qquad z_i=d_ip_j
 \quad(i\in N_j).
\]
Every pool has a positive capped capacity. Its sum $X_j$ is continuous, and
the pool sums are independent. Thus ties have probability zero and
$\sum_jp_j=1$. The mechanism pays deterministic expectations, not random
realizations of these amounts. Let $y_i=\max\{z_i,M_i/k\}$.

\begin{theorem}\label{thm:upper}
The base $z$ covers every coalition and is own-report monotone on all
nonnegative real profiles. For $R_k>e$ defined by
$R_k\log(R_k/e)=\log k$,
\[
 \sum_i z_i\le\left(R_k+\frac1{\log R_k}\right)T,
 \qquad
 \sum_i y_i\le\left(1+R_k+\frac1{\log R_k}\right)T.
\]
Both factors are $(1+o(1))\log k/\log\log k$ as $k\to\infty$,
uniformly over all finite owner counts. The second rule meets $H_i/k$.
\end{theorem}
\begin{proof}
By Lemma~\ref{lem:cap}, for every $S$,
\[
 z(S)=\sum_jp_jd(S\cap N_j)\ge\min_jd(S\cap N_j)=v_c(S).
\]
For own-monotonicity, fix $i\in N_j$, write its report as $x$, let
$a=\sum_{h\in N_j\setminus\{i\}}c_h$ and
$b=\min_{\ell\ne j}C_\ell$. If $b=0$ all payments vanish. Otherwise
$T(x)=\min\{x+a,b\}$ and $d_i(x)=\min\{x,b\}$. Conditional on $U_{-i}$,
put
\[
 R=\sum_{h\in N_j\setminus\{i\}}U_h\min\{c_h,T(x)\},\qquad
 Y(x)=\min_{\ell\ne j}\sum_{h\in N_\ell}U_h\min\{c_h,T(x)\}.
\]
Here $R$ is independent of $x$: if $a\le b$, every other own-pool capacity
is at most $a\le T(x)$; if $a>b$, $T(x)=b$ throughout. Also $Y(x)$ is
nondecreasing. Integrating only $U_i$ gives
\[
 d_i\Prob(U_id_i<Y-R\mid U_{-i})=\min\{d_i,(Y-R)_+\}.
\]
At $d_i=0$ both sides are zero. The expression is nondecreasing in $x$;
its expectation is $z_i$. Lemma~\ref{lem:floor} proves the floor assertions.

For the budget, normalize $a_i=d_i/T\in[0,1]$, put $\rho_j=D_j/T$, and
let $W$ be the winning pool. Then $z(N)/T=\E\rho_W$. A short pool has
$X_s/T\le1$, so
$p_j\le\Prob(\sum_{i\in N_j}a_iU_i\le1)$.
For $t>0$ and $a\in[0,1]$, concavity of $u\mapsto u^a$ gives
\[
 \E e^{-taU}\le(\E e^{-tU})^a
 =\left(\frac{1-e^{-t}}t\right)^a\le t^{-a}.
\]
Independence and exponential Markov imply
$p_j\le e^t t^{-\rho_j}$. Set $t=\rho_j$. Because $(e/s)^s$ decreases
for $s\ge e$,
\[
 \Prob(\rho_W>s)\le\min\{1,k(e/s)^s\},\qquad s\ge e.
\]
At $s=R_k$ the second expression is one; its logarithmic derivative is
$-\log s\le-\log R_k$ thereafter. Tail integration therefore yields
\[
 \E\rho_W\le R_k+\int_{R_k}^{\infty}
 e^{-(s-R_k)\log R_k}\,ds=R_k+\frac1{\log R_k}.
\]
Finally, $R_k\log(R_k/e)=\log k$ implies
$R_k\sim\log k/\log\log k$. Add the floor cost from
Lemma~\ref{lem:floor}.
\end{proof}

\paragraph{Identification with an established index.}
For a fixed profile define the continuous monotone game
\[
 f_d(u)=\frac1T\min_j\sum_{i\in N_j}d_iu_i,\qquad u\in[0,1]^N.
\]
It has endpoint values zero and one. Kurz's absolute continuous Banzhaf
index integrates $f_d(1,u_{-i})-f_d(0,u_{-i})$ over the other coordinates.
With $R,Y$ as above,
\[
 T\bigl[f_d(1,u_{-i})-f_d(0,u_{-i})\bigr]
 =\min\{R+d_i,Y\}-\min\{R,Y\}=\min\{d_i,(Y-R)_+\}.
\]
Hence $z_i=T\operatorname{BZI}_i(f_d)$. This pointwise identity does not
by itself prove monotonicity across reports, since $f_d$ changes with $c$;
the preceding conditional integration argument supplies that step.
Exact evaluation can require exponentially many subset-sum knots. No
polynomial-time or arbitrary rounded-payment implementation is claimed.

\section{The unrestricted budget grows sublogarithmically}\label{sec:lower}
\begin{lemma}[A finite hierarchy certificate]\label{lem:hierarchy}
Fix integers $d\ge1$ and $m>d$. Choose integers $n_0\ge2$ and
$n_h\ge1$ for $1\le h\le d$, with $\sum_{h=0}^d n_h=k$.
Every eligible all-profile rule on $k$ pools of $m$ owners has factor at least
\begin{equation}\label{eq:hierarchy}
 \frac{A}{m\left[1+A\displaystyle\sum_{h=0}^{d-1}
          \frac{n_h}{(h+1)(n_{h+1}+1)}\right]},
 \qquad A=(d+1)m-d^2.
\end{equation}
\end{lemma}
\begin{proof}
If the rule has infinite factor there is nothing to prove; otherwise write
its factor as $B$. Average it over the finite group of pool permutations and
independent owner permutations inside pools. Relabeling preserves the game,
and the averaging preserves nonnegativity, coalition coverage, budgets and
every own-report comparison. A lower bound for this invariant rule therefore
bounds the original rule. No anonymity axiom or strategic cloning is used.

All reports below are either $1$ or $m$. A type-$h$ pool has $h$ high-report
owners and $m-h$ low-report owners. At a base profile $S$ let $n_h$ pools
have type $h$, and write $x_h,z_h$ for the invariant payments of low and
high owners at that profile (there is no $z_0$). We have $T=m$.
For each $h<d$, raise one original low owner's report in a type-$h$ pool
from $1$ to $m$, giving a profile $S_h$. Its counts are $n_h-1$ and
$n_{h+1}+1$ at these two types. An all-low pool remains, since $n_0\ge2$,
so its grand value is still $m$.
All $(h+1)(n_{h+1}+1)$ high owners in its type-$(h+1)$ pools receive the
same payment. Their nonnegative total is at most $Bm$. Monotonicity of
the changing owner alone gives the simultaneous base-profile bounds
\begin{equation}\label{eq:exports}
 x_h\le \frac{Bm}{(h+1)(n_{h+1}+1)},\qquad 0\le h<d.
\end{equation}
No unchanged owner's payment is transported between profiles.

Two actual coalitions at $S$ complete the argument. The first chooses one
low owner from every pool; this is possible because $m>d$ and has worth one.
The second chooses $m-d+1$ low owners in every type-$h$ pool with $h<d$,
and one high owner in every type-$d$ pool. Its worth is $m-d+1$.
Their coverage inequalities are
\begin{align}
 \sum_{h=0}^d n_hx_h&\ge1,\label{eq:hierarchy-c0}\\
 (m-d+1)\sum_{h<d}n_hx_h+n_dz_d&\ge m-d+1.
 \label{eq:hierarchy-c1}
\end{align}
Multiply these by $m-d$ and $d$, respectively. The type-$d$ coefficients
then equal that type's full budget coefficients. Thus
\begin{align*}
 A&\le A\sum_{h<d}n_hx_h+n_d[(m-d)x_d+dz_d]\\
  &\le A\sum_{h<d}n_hx_h+Bm\\
  &\le Bm\left[1+A\sum_{h<d}
            \frac{n_h}{(h+1)(n_{h+1}+1)}\right].
\end{align*}
The middle step drops only other nonnegative base payments, and the last
uses \eqref{eq:exports}. Rearranging proves \eqref{eq:hierarchy}.
\end{proof}

\begin{theorem}[Single-base hierarchy bound]\label{thm:half-order}
For the factor in \eqref{eq:Bk},
\[
 (1/2-o(1))\frac{\log k}{\log\log k}
 \le B_k\le(1+o(1))\frac{\log k}{\log\log k}.
\]
In particular no absolute factor works uniformly in both pool and owner
counts. The same bounds hold to first order with the optional $H_i/k$ floor.
\end{theorem}
\begin{proof}
For $d\ge2$ set
\[
 q_d=\lceil\log(d+1)\rceil,\quad m=dq_d,\quad
 A=(d+1)m-d^2,\quad L=\lceil A(1+\log d)^2\rceil,
 \quad n_h=2L^h.
\]
Then $m>d$ and $n_h/(n_{h+1}+1)<1/L$. Since
$\sum_{h<d}1/(h+1)\le1+\log d$, the bracket in
\eqref{eq:hierarchy} is at most $1+1/(1+\log d)$. Moreover
\[
 A/m=d+1-d/q_d=(1-o(1))d,\qquad
 \log L=(2+o(1))\log d.
\]
For $k_d=2\sum_{h=0}^d L^h$ this gives
$\log k_d=(2+o(1))d\log d$, and the lower factor is $(1-o(1))d$.
For any $k\ge k_d$, add $k-k_d$ type-$d$ pools. Only the final ratio in
\eqref{eq:hierarchy} changes, and it decreases. Thus the same lower bound
holds on the padded ownership, without assuming that $B_k$ is monotone in $k$.
For fixed $0<\epsilon<1/2$, choose
$d=\lfloor(1/2-\epsilon)\log k/\log\log k\rfloor$.
Then $k_d\le k$ for all sufficiently large $k$. Apply the padded bound,
take the lower limit and then let $\epsilon\downarrow0$.
Theorem~\ref{thm:upper} gives the upper bound for every ownership in
\eqref{eq:Bk}. Lemma~\ref{lem:floor} adds at most one to the upper factor;
the lower already holds without the floor.
\end{proof}

\paragraph{An elementary divergence sequence and its dual.}
For $d\ge1$, take $m=2d$, $A=d(d+2)$, $L=2Ad$ and $n_h=2L^h$.
The bracket in \eqref{eq:hierarchy} is strictly below $3/2$, so
$B>(d+2)/3$. In the convention that every row has left side at least
right side, a nonnegative certificate uses the two coalition rows with
weights $m-d,d$, the base budget with weight one, the changing-owner
monotonicity row at type $h$ with weight $An_h$, and the budget at $S_h$
with weight $An_h/[(h+1)(n_{h+1}+1)]$. All remaining payment coefficients
are nonpositive and are canceled by nonnegativity rows. This yields
\eqref{eq:hierarchy} using exact rational coefficients.

\subsection{Averaging profiles closes the leading-constant gap}
The preceding hierarchy charges a separate destination budget for each
exporting type. We now average all histograms with carefully chosen weights
so the same budget controls every incoming type simultaneously.

\begin{lemma}[Weighted-profile hierarchy]\label{lem:weighted-hierarchy}
Fix integers $d\ge1$, $m>d$, $k\ge2$, and a real $\alpha>0$. Set
\[
 A=(d+1)m-d^2,\quad \lambda=A/\alpha,\quad
 Z=\sum_{h=0}^d\frac{\lambda^h}{h!},\quad
 p_h=\frac{\lambda^h}{h!Z}.
\]
Every eligible all-profile rule on $k$ pools of $m$ fixed owners has factor
\begin{equation}\label{eq:weighted-bound}
 B\ge\frac{A}{(1+\alpha)\mathcal T},\qquad
 \mathcal T=m+(m-1)\sum_{r=1}^d
             \left(\sum_{h=r}^d p_h\right)^k
 \le m\bigl[1+d e^{-kp_0}\bigr].
\end{equation}
\end{lemma}
\begin{proof}
Use finite symmetrization and reports in $\{1,m\}$ as before. We now allow
every histogram $n=(n_0,\ldots,n_d)$ with $\sum_hn_h=k$, including those
without a type-zero pool. Denote its low and high payments by $x_h(n)$ and
$z_h(n)$; terms for absent types are assigned zero. Since $m>d$, every
pool has a low owner. The two coalitions from Lemma~\ref{lem:hierarchy}
still give
\begin{equation}\label{eq:weighted-cover}
 A\le A\sum_{h<d}n_hx_h(n)
          +n_d[(m-d)x_d(n)+dz_d(n)].
\end{equation}
If every pool has type $d$, the second coalition is worth $m$, which is
at least the $m-d+1$ used in this inequality.

Average report profiles by drawing their $k$ pool types independently
from $p$. This is only a proof device; the payment rule stays deterministic.
The multinomial weight of histogram $n$ is
\[
 \mu(n)=k!\prod_{h=0}^d\frac{p_h^{n_h}}{n_h!}.
\]
For $h<d$ and $n_h>0$, raise one original low report in a type-$h$ pool.
The destination $n'=n-e_h+e_{h+1}$ satisfies
$x_h(n)\le z_{h+1}(n')$ by own-report monotonicity. Direct cancellation
in the multinomial weights yields
\begin{equation}\label{eq:transport}
 \mu(n)n_h=\frac{p_h}{p_{h+1}}\mu(n')n'_{h+1}.
\end{equation}
The map is a bijection between histograms with $n_h>0$ and those with
$n'_{h+1}>0$. Summing the changing-owner inequalities therefore gives
\[
 \E[n_hx_h]\le\frac{p_h}{p_{h+1}}\E[n_{h+1}z_{h+1}].
\]
No constancy of $T$ or retention of unchanged owners' payments is required.
Since $A p_h/p_{h+1}=\alpha(h+1)$, averaging
\eqref{eq:weighted-cover} implies
\begin{align*}
 A&\le \alpha\E\left[\sum_{r=1}^d r n_rz_r\right]
       +\E\left[n_d((m-d)x_d+dz_d)\right]\\
  &\le(1+\alpha)\E[y(N)]\le(1+\alpha)B\E[T].
\end{align*}
Both bracketed sums are parts of the same nonnegative payment budget.

A type-$h$ pool has total $m+(m-1)h$. If $H_1,\ldots,H_k$ are the
independent types, then $T=m+(m-1)\min_j H_j$. The tail-sum identity gives
$\E T=\mathcal T$. The probability of no type-zero pool is $(1-p_0)^k$;
otherwise $T=m$, while always $T\le m(1+d)$. Hence
$\mathcal T\le m[1+d(1-p_0)^k]\le m[1+d e^{-kp_0}]$.
This proves \eqref{eq:weighted-bound}.
\end{proof}

\begin{theorem}[Sharp unrestricted leading constant]\label{thm:order}
Under the model of Section~\ref{sec:model},
\[
 B_k=(1+o(1))\frac{\log k}{\log\log k},\qquad
 B_k^{\mathrm{floor}}=(1+o(1))\frac{\log k}{\log\log k}.
\]
Both statements allow arbitrary finite ownership sizes and all nonnegative
real report profiles. No proportionality or anonymity is imposed.
\end{theorem}
\begin{proof}
For $d\ge2$, choose $q=\lceil\log(d+1)\rceil$, $m=dq$,
$\alpha=1/q$, and $\lambda=Aq$ in Lemma~\ref{lem:weighted-hierarchy}.
Put $k_d=\lceil2qZ\rceil$. For every integer $k\ge k_d$, $kp_0\ge2q$,
so the lemma gives
\begin{equation}\label{eq:sharp-finite}
 B\ge\frac{d+1-d/q}{(1+1/q)(1+d e^{-2q})}=(1-o(1))d.
\end{equation}
Here $\lambda=(d+1)dq^2-d^2q>d$ and
$\log\lambda=(2+o(1))\log d$. The summands $\lambda^h/h!$ increase
for $0\le h\le d$, whence
\[
 \frac{\lambda^d}{d!}\le Z\le(d+1)\frac{\lambda^d}{d!}.
\]
Using $\log(d!)=d\log d-d+O(\log d)$, we obtain
\[
 \log Z=d\log\lambda-\log(d!)+O(\log d)
       =(1+o(1))d\log d,
 \qquad \log k_d=(1+o(1))d\log d.
\]
For fixed $0<\epsilon<1$, choose
$d=\lfloor(1-\epsilon)\log k/\log\log k\rfloor$.
Then $k\ge k_d$ eventually, so \eqref{eq:sharp-finite} applies on that
exact $k$-pool ownership. Thus
$\liminf_{k\to\infty}B_k\log\log k/\log k\ge1-\epsilon$.
Let $\epsilon\downarrow0$ and combine with Theorem~\ref{thm:upper}.
The floor inequalities $B_k\le B_k^{\mathrm{floor}}\le B_k+1$
give the second assertion.
\end{proof}

\paragraph{A finite exact dual.}
At every histogram use its two coalition rows with weights
$\mu(n)(m-d)$ and $\mu(n)d$, its budget row with weight
$\mu(n)(1+\alpha)$, and its type-$h$ changing-owner row with weight
$A\mu(n)n_h$. Identity~\eqref{eq:transport} makes every residual owner
coefficient nonpositive; type-$d$ high coefficients cancel exactly.
For rational $\alpha$, all weights are rational. The actual coalition
constants sum to $A+d(d-1)p_d^k$, at least $A$, because the all-type-$d$
profile has the larger second-coalition worth. The budget coefficient is
exactly $(1+\alpha)\mathcal T$. The separate \texttt{verify.py} row checker verifies certificates generated
by \texttt{poisson\_certificate.py} in small instances, reconstructing each
row without presupposing the transport identity.

\paragraph{Scope for max-flow games.}
Give each pool a layer of parallel directed edges between consecutive
vertices, one edge per owner. The max flow of every coalition is exactly
the minimum of its layer totals. Hence the lower bound also precludes an
absolute truthful coalition-covering budget factor on the class of all
edge-owned max-flow games. Our upper bound is proved only for these layered
games; no corresponding upper bound for arbitrary networks follows.

\subsection{A smaller explicit finite family}
The following earlier bound is weaker asymptotically but uses only two
linked profiles and can be useful at specified finite parameters.

The next theorem uses neither proportionality nor a fairness floor.
\begin{theorem}\label{thm:unrestricted}
Fix $k\ge3$ and $r\ge2$ owners per pool, and set $t=2r-2$. Every
nonnegative covering own-monotone rule with factor $1+\beta$ on all real
profiles satisfies
\[
 \beta\ge\frac{(k-1)(t-1)}{t(k-1+2t)}.
\]
Therefore any total-payment constant uniform over both pool and owner counts
must be at least two.
\end{theorem}
\begin{proof}
Average the rule over all pool permutations and all within-pool owner
permutations. Relabeling and averaging preserve the guarantees, so a bound
for the resulting invariant rule bounds the original rule. This is a proof
device on fixed owners, not permission for an owner to split or merge.
Put $q=r-2$, so $t=2q+2$. Use two pool types
\[
 L=(1,1,\underbrace{2,\ldots,2}_{q}),\qquad
 H=(1,\underbrace{2,\ldots,2}_{q},t).
\]
Their totals are $t$ and $2t-1$, respectively. All capacities are positive.
At profile $X$ with one $L$ and $k-1$ copies of $H$, let the payments be
$a,a,b,\ldots,b$ in $L$ and $u,v,\ldots,v,w$ in each $H$. Write
$P=2a+qb$, $Q=u+qv+w$. The budget is $P+(k-1)Q\le t(1+\beta)$.
Actual coalitions give
\begin{align*}
 a+(k-1)u&\ge1, &2a+(k-1)v&\ge2,\\
 a+qb+(k-1)(u+qv)&\ge t-1, &P+(k-1)w&\ge t.
\end{align*}
The inequality involving $v$ is used only if $q>0$. The choices in the $L$ pool are,
respectively, one unit owner, both unit owners, one unit plus all 2-owners,
and all owners. Each $H$ contributes its unit owner, one 2-owner, the unit
plus all 2-owners, or its $t$-owner. Combine these inequalities with weights
$1,q,t-1,1$. The left side is
$t[P+(k-1)Q]-(t-1)(k-1)w$ and the right side is $t^2$. Hence
\begin{equation}\label{eq:X}
 (t-1)(k-1)w\le t^2\beta.
\end{equation}
At profile $Y$ with two $L$ pools and $k-2$ copies of $H$, denote analogous
payments by $A,B,U,V,W$. Its budget is
$2(2A+qB)+(k-2)(U+qV+W)\le t(1+\beta)$. Coalitions selecting one unit,
both units, or all owners from each $L$ yield
\[
 2A+(k-2)U\ge1,\quad4A+(k-2)V\ge2,\quad
 2(2A+qB)+(k-2)W\ge t.
\]
Again omit the middle inequality if $q=0$. Their combination with weights
$1,q,1$ is the budget sum plus $2(t-1)A\ge2t-1$, so
\begin{equation}\label{eq:Y}
 2(t-1)A\ge t-1-t\beta.
\end{equation}
Raising one unit report in an $L$ pool of $Y$ from $1$ to $t$ gives $X$
up to relabeling. Own-monotonicity gives $w\ge A$. Equations
\eqref{eq:X}--\eqref{eq:Y} therefore imply
\[
 (k-1)(t-1-t\beta)\le2(t-1)(k-1)A
 \le2(t-1)(k-1)w\le2t^2\beta,
\]
which rearranges to the claim. No sign assumption on $t-1-t\beta$ is needed.
For $r=n+1$, $t=2n$, $k=8n^2+1$, the displayed lower bound is
$(2n-1)/(2n+1)\to 1$.
\end{proof}
This two-profile family alone only gives a necessary constant. The hierarchy
in Theorem~\ref{thm:order} establishes the stronger growing lower bound.

\section{Sharp asymptotics for proportional payments}\label{sec:prop}
Consider final base payments of the form
\[
 z_i(c)=d_i(c)\pi_j(c),\qquad i\in N_j,
\]
where $\pi_j$ is an arbitrary nonnegative function of the whole labeled
profile and fixed ownership. There is no normalization, anonymity,
continuity, locality, or independence assumption on these prices.

\begin{theorem}\label{thm:prop}
On $k\ge2$ pools with $r\ge2$ owners each, every covering own-monotone
capped-proportional rule with factor $B$ satisfies
\begin{equation}\label{eq:prop}
 B\ge\frac{kr}{k+(r-1)r^r}.
\end{equation}
The same lower bound holds for raw-capacity-proportional payments $c_i\pi_j$.
\end{theorem}
\begin{proof}
Let $H$ be the all-unit profile. Its value is $r$, and its budget implies
$\sum_j\pi_j(H)\le B$. Let $L_j$ have all owners of pool $j$ at $1/r$
and all other owners at $1$. It has value one. The coalition containing all
of pool $j$ and one owner of each other pool has worth one, so
\[
 \pi_j(L_j)+\sum_{\ell\ne j}\pi_\ell(L_j)\ge1,
 \qquad \pi_j(L_j)+r\sum_{\ell\ne j}\pi_\ell(L_j)\le B.
\]
Multiply the first by $r$ and subtract the second:
\begin{equation}\label{eq:price-low}
 B+(r-1)\pi_j(L_j)\ge r.
\end{equation}
Raise each of pool $j$'s reports successively from $1/r$ to $1$, obtaining
$P_{j,0}=L_j,\ldots,P_{j,r}=H$. Every capacity is at most one and every
profile value is at least one, so capping changes no report. At step $h$,
own-monotonicity of the changing owner gives
\[
 \pi_j(P_{j,h+1})\ge\pi_j(P_{j,h})/r,
 \qquad\text{hence}\qquad r^r\pi_j(H)\ge\pi_j(L_j).
\]
This transports a \emph{common pool price} between changes made by different
owners; that identity is precisely where proportionality is used. Substitute
in \eqref{eq:price-low}, sum over $j$, and use the budget at $H$:
\[
 kr\le kB+(r-1)r^r\sum_j\pi_j(H)
 \le[k+(r-1)r^r]B.
\]
The uncapped version has the same proof because the witness capacities are
never capped.
\end{proof}

\paragraph{A nonnegative inequality certificate.}
Write all constraints as left side at least right side. For each $j$, use
$r$ times its $L_j$ coverage row, once its budget row
$B-\pi_j-r\sum_{\ell\ne j}\pi_\ell\ge0$, and for each $h=0,\ldots,r-1$
use $(r-1)r^{h+1}$ times the actual monotonicity row
$\pi_j(P_{j,h+1})-\pi_j(P_{j,h})/r\ge0$. Add once in total
$(r-1)r^r$ times $B-\sum_j\pi_j(H)\ge0$. Prices cancel, leaving
\eqref{eq:prop}. All weights are nonnegative integers.

\begin{corollary}\label{cor:minimax}
Let $A_k$ be the infimum of worst-case total-payment factors among families
of floor-free capped-proportional rules, defined for every finite fixed
ownership with $k$ pools and all nonnegative real capacities, covering all
coalitions and satisfying own-monotonicity. Explicitly, if $\mathfrak P_k$
contains all finite ownership partitions with $k$ nonempty pools and
$\mathcal F=(z^P)_{P\in\mathfrak P_k}$ ranges over such rule families, then
\[
 A_k=\inf_{\mathcal F}\ \sup_{P\in\mathfrak P_k}\ \sup_{c:T(c)>0}
 \frac{z^P(c)(N)}{T(c)}.
\]
Rules for different fixed ownerships need not satisfy cross-ownership
consistency. We have
\[
 A_k=(1+o(1))\frac{\log k}{\log\log k}.
\]
\end{corollary}
\begin{proof}
For fixed $\epsilon>0$, let $L=\log k$ and
$r=\lfloor(1-\epsilon)L/\log L\rfloor$. For all sufficiently large $k$,
$r\ge2$ and
$\log((r-1)r^r)\le(1-\epsilon)L+o(L)$.
Thus $(r-1)r^r/k\to0$, and Theorem~\ref{thm:prop} gives
$A_k\ge(1-o(1))r$. Take the lower limit and then $\epsilon\downarrow0$.
The base of Theorem~\ref{thm:upper} belongs to this class and gives the
matching upper limit.
\end{proof}
The floor-restoring maximum need not remain proportional. Accordingly the
corollary is floor-free; it is not a claim about the minimax among final
proportional rules also meeting the floor, nor about unrestricted rules.

\section{A nonlinear escape on the linked profiles}\label{sec:escape}
Fix equal pool size $r$. Let $\mathcal D_1$ contain \emph{every labeled}
profile with reports in $\{1/r,1\}$ and at most one pool containing a low
report. The family includes all orders of the report changes in
Theorem~\ref{thm:prop}.

\begin{theorem}\label{thm:escape}
On $\mathcal D_1$ there are nonnegative covering payments that are own-monotone
on every retained report edge and have factor at most $2$. The original-core
floor can be added with factor at most $3$.
\end{theorem}
\begin{proof}
At $H$, pay $1/k$ to each owner. Otherwise let the altered pool have $h$
unit owners, $0\le h<r$, and define
\[
 T_h=h+(r-h)/r,\quad \delta_h=(r-h)/r,\quad z_h=T_h/r,
\]
\begin{equation}\label{eq:abc}
 a_h=\frac{(\delta_h-z_h)_+}{r-h},\qquad
 b_h=1-z_h,\qquad e_h=\frac{z_h}{k-1}.
\end{equation}
Pay its low owners $a_h$, its unit owners $b_h$, and every outside owner
$e_h$. For a coalition, let $b$ be its unit-owner count in the altered pool,
$s$ its selected low capacity, and $n$ the smallest selected-owner count
outside that pool. Its worth is $\min\{b+s,n\}$, and its payment is at least
\[
 b(1-z_h)+s(1-z_h/\delta_h)_++nz_h.
\]
If $n\le b$, the first and last terms already give at least $n$. Otherwise
$n\ge b+1$ because $n,b$ are integers. It suffices that
$z_h+s(1-z_h/\delta_h)_+\ge s$, which follows from $s\le\delta_h$ in
both cases $z_h\le\delta_h$ and $z_h\ge\delta_h$. At $H$, the average
selected-owner count bounds its minimum, proving coverage there too.

The total off $H$ is
$F_h=h(1-z_h)+(\delta_h-z_h)_++rz_h$.
If $z_h\le\delta_h$, this equals $T_h+(r-h-1)z_h\le2T_h$.
Otherwise it equals $h+(r-h)z_h\le T_h+T_h$.

Every retained increase raises a low owner to one. For $0\le s\le r-2$,
direct substitution in \eqref{eq:abc} gives
\begin{equation}\label{eq:ratio}
 \frac{a_s}{b_{s+1}}=
 \frac1{r-1}\left(1-\frac{s^2}{(r-s)(r-s-1)}\right)_+\le1,
\end{equation}
with $b_{s+1}>0$. Also $a_{r-1}=0$. Thus $a_s\le b_{s+1}$ for internal
edges and $a_{r-1}\le1/k$ on the last-low boundary. This exhausts all labeled
retained edges. Apply Lemma~\ref{lem:floor}.
\end{proof}
The preceding certificate cannot show unrestricted divergence on these same
profiles, even after adding the stated fairness floor. This is a separation
\emph{on the witness family}; it is not a separation between two global
minimax asymptotics. Appendix~\ref{app:binary} extends the factor-two base to
two altered binary pools while preserving this boundary.

\section{Continuous reports with a bounded number of altered pools}
Here equal pool size $r\ge2$ and the fixed public reference capacity $1$ are
assumptions. A pool is altered if at least one report differs from $1$.
There is no data-dependent choice of reference or ownership.

\begin{theorem}\label{thm:continuous}
For reports in $[0,1]$ and at most $m$ altered pools, the following rule is
nonnegative, covering and own-monotone, with factor at most $2m+4$:
\[
 b=\frac1k\min\left\{1,\frac{2T}{r-1}\right\},\quad
 E_i=\sum_{h\in N_j\setminus\{i\}}(1-c_h),\quad
 y_i=b+\frac{2E_i}{r-1}\frac{Tc_i}{C_j}\quad(i\in N_j).
\]
Set all payments to zero at $T=0$. On $[1/r,1]$, the factor improves to
$2m+3$. The original-core floor adds at most one in either domain.
\end{theorem}
\begin{proof}
When $T>0$, $C_j\ge T$ so all denominators are positive. Fixing other
reports makes $E_i$ constant and nonnegative. If their own-pool sum is $A_i$,
both $T(t)$ and $t/(A_i+t)$ are nonnegative and nondecreasing wherever the
ratio is defined. Their product is nondecreasing; so is $b(T(t))$. At a
zero-value endpoint all payments are zero. This proves own-monotonicity on
the full cube, including all boundary transitions. If the changing owner
was its pool's last non-unit owner, $E_i=0$ throughout that change.

If $T\ge(r-1)/2$, then $b=1/k$ and, writing
$n_j=|S\cap N_j|$, the uniform payments give
$|S|/k\ge\min_j n_j\ge v_c(S)$ because capacities are at most one.
If $0<T\le(r-1)/2$, choose a pool with $C_j=T$. For each of its owners,
$E_i=r-1-(T-c_i)\ge(r-1)/2$, so its additional payment is at least $c_i$.
Payments in this one pool alone cover every coalition. These two regimes
include their common boundary; $T=0$ is immediate.

Because $0\le E_i\le r-1$, the additional payments in any pool sum to at
most $2T$. They vanish in unaltered pools. The uniform part costs
$r\min\{1,2T/(r-1)\}\le4T$. This proves the first bound.
On $[1/r,1]$ we have $T\ge1$: for $r=2$ the uniform part is at most
$2\le2T$; for $r\ge3$ it is at most $2rT/(r-1)\le3T$.
Lemma~\ref{lem:floor} proves the floor claim.
\end{proof}
For two altered pools this gives base $7T$ and fair $8T$ on $[1/r,1]$.
It revises the earlier binary payments: at $k=r=2$, on
$((1/2,1/2),(1,1))$, the new base is $(1,1,1/2,1/2)$, whereas
Theorem~\ref{thm:escape} gives $(1/4,1/4,1/2,1/2)$.
Preservation of the old boundary on continuous reports is not proved
impossible, and none of these constants is asserted to be minimax sharp.

The formula is continuous when $T>0$; at $T=0$ its full-cube bound with
$m=k$ forces every coordinate to tend to zero. Thus a single direct
continuous rule works on every finite report grid. No inference from grid
feasibility or exchange of limiting operations is required.

\paragraph{An exact limitation of this formula.}
Take $r=3$, any $k\ge2$, and $c_i=1/3$ everywhere. Then $T=1$, $E_i=4/3$,
$b=1/k$, and $y_i=1/k+4/9$. Total payment is exactly $3+4k/3$.
The floor is inactive since $M_i/k=1/(3k)<b$. Therefore this particular
rule has genuinely growing cost when all pools are altered. This says nothing
about the best possible unrestricted rule.


\section{Bounds in owner count and resource dimension}\label{general:section}
The number of supporting forms can be much larger than the economic description.
We now obtain bounds directly in the numbers of fixed owners and resources.
Here the production technology may be any finite-valued monotone concave
function $f:\R_+^d\to\R_+$ homogeneous of degree one; a finite min-linear
representation is not required. The report space and eligibility conditions
remain unchanged. In particular, there is no additional fairness floor.

\begin{theorem}[General owner and dimension bounds]\label{general:sqrt}
For every fixed $n,d\ge1$ and every such technology, there exists an eligible
deterministic rule on all nonnegative real report profiles with
\begin{equation}\label{general:combined}
 \sum_i M_i(b)\le
 \min\left\{d,\ n\sqrt{\frac{2}{n+1}},\ 1+2\sqrt d\right\}T.
\end{equation}
The individual bounds are uniform in the other parameter. They also hold when
$f$ has downward discontinuities at the boundary. More precisely, for every
fixed $0<a<1$ the supporting-price/endpoint rule averaged over $q\in[0,a]$
has both bounds
\begin{align}
 B_n(a)&=\frac{na}{2}+
 \frac{n[1-(1-a)^{n+1}]}{(n+1)a},\label{general:owner}\\
 D_d(a)&=\frac{1+2d[-a-\log(1-a)]}{a}.\label{general:dimension}
\end{align}
The owner expression remains valid at $a=1$ by polynomial integration.
No polynomial-time evaluation or preprocessing claim is made.
\end{theorem}

\subsection{Supporting prices on boundary faces}
For $x\ne0$, restrict $f$ to its coordinate face
$F_x=\{z\ge0:z_e=0\text{ whenever }x_e=0\}$.
At the relative-interior point $x$, finite concavity gives a supergradient $u$.
Monotonicity forces each active coordinate of $u$ to be nonnegative, and
homogeneity applied above and below $x$ gives $u\cdot x=f(x)$.
For a boundary point $z\in F_x$, apply the supporting inequality at
$(1-t)x+tz$, $0<t<1$, and the concavity lower bound on that segment to obtain
$u\cdot z\ge f(z)$. Thus boundary continuity is unnecessary.
Choose the unique minimum-norm vector in the nonempty closed convex set of
such nonnegative supports on $F_x$. Extend its linear form to
\[
 \ell_x(z)=\begin{cases}u\cdot z,&z\in F_x,\\+\infty,&z\notin F_x.
 \end{cases}
\]
Set $\ell_0(0)=0$ and $\ell_0(z)=+\infty$ for $z\ne0$.
This is one fixed selection depending only on $x$. For finite min-linear
technologies the original fixed-priority global minimizing form is an
alternative; the bounds below apply to either choice.

For any fixed mask $B\subseteq N$, put $x_B=\sum_{i\in B}b_i$ and
$A_i(B)=\min\{\ell_{x_B}(b_i),T\}$. This vector covers every coalition:
if any member's raw price reaches $T$, its capped payment alone suffices;
otherwise all relevant prices are finite on the support face and their sum
majorizes the coalition's worth. For $i\in B$,
\[
 f(x_B)-f(x_B-b_i)\ge \ell_{x_B}(b_i),
\]
because the support touches at $x_B$ and majorizes $f(x_B-b_i)$.
The marginal is at most $T$. Replace open owners' prices by these marginals
and keep $A_i(B)$ for closed owners. Coverage is preserved. A closed owner's
selected price excludes its own report, while its raw price and cap are
nondecreasing in that report. An open owner's marginal has a fixed off term
and a nondecreasing on term. This proves eligibility for every fixed mask
and every report-independent mixture of masks.

With independent missing probability $q$, retain the notation $p_i^q,m_i^q,r_i^q$
from \eqref{prod:p}--\eqref{prod:hybrid}, now using $\ell_{x_B}$.
At a fixed profile every capped mask payment is finite, so the expectations
are finite polynomials in $q$. Define $M_i^a=a^{-1}\int_0^a r_i^q\,dq$.
The endpoints have measure zero, including when $a=1$.
Writing $g(q)=\E f(x_B)$ and $C(q)=q\sum_i p_i^q$, differentiation of the
finite mask polynomial gives
\begin{equation}\label{general:endpoint}
 \sum_i r_i^q=C(q)-(1-q)g'(q),\qquad
 \int_0^a-(1-q)g'(q)\,dq
 =T-(1-a)g(a)-\int_0^a g(q)\,dq.
\end{equation}
No derivative of $f$ with respect to reports occurs here.

\subsection{Owner-count bound}
Every closed payment is at most $T$, whence $C(q)\le nqT$.
The all-open mask alone gives $g(q)\ge(1-q)^nT$.
Substituting both estimates in \eqref{general:endpoint}, integrating and
dividing by $a$ yields
\[
 \sum_i M_i^a\le\left(\frac{na}{2}+
 \frac{n[1-(1-a)^{n+1}]}{(n+1)a}\right)T.
\]
Dropping the bracketed power term gives
$B_n(a)\le na/2+n/[(n+1)a]$. The fixed choice
$a=\sqrt{2/(n+1)}\le1$ equalizes these terms and proves the owner bound
in \eqref{general:combined}. For $n=1$ the exact expression is $B_1(a)=1$.
The former $O(n/\log n)$ upper bound recorded in the supplement is therefore
superseded by $(\sqrt2+o(1))\sqrt n$.

\subsection{Resource-dimension bound}
The following scalar identity controls contributions from closed owners
without counting owners or supporting forms.
\begin{lemma}[Bernoulli reciprocal bound]\label{general:reciprocal}
Let $u_i\ge0$, independently retain each $i$ with probability $p=1-q$, and
write $S_B=\sum_{i\in B}u_i$. Define $\min\{u_i/S_B,1\}$ to be zero when
$u_i=0$ and one when $u_i>0=S_B$. Then
\[
 \E\sum_{i\notin B}\min\{u_i/S_B,1\}\le\frac{2q}{1-q}.
\]
\end{lemma}
\begin{proof}
With zero-over-zero interpreted as zero,
$\min\{u/S,1\}\le2u/(u+S)$. Conditioning on the own bit gives
\[
 \E\sum_{i\notin B}\frac{u_i}{u_i+S_B}
 =\frac qp\E\sum_{i\in B}\frac{u_i}{S_B}
 =\frac qp\Prob(S_B>0)\le\frac qp.
\]
In the first equality both conditioned ratios are
$u_i/(u_i+\sum_{j\ne i}B_ju_j)$; the other bits have the same law.
\end{proof}

Suppose $T>0$. For every mask and owner,
\begin{equation}\label{general:coordinate}
 \frac{A_i(B)}T\le
 \sum_{e=1}^d\min\left\{\frac{b_{ie}}{(x_B)_e},1\right\},
\end{equation}
with the lemma's conventions. If $b_i$ has a positive coordinate outside
$F_{x_B}$, the right side is at least one. Otherwise the selected nonnegative
support satisfies $\sum_eu_e(x_B)_e=f(x_B)\le T$, so each active term obeys
$u_eb_{ie}/T\le b_{ie}/(x_B)_e$. Apply
$\min\{\sum_e z_e,1\}\le\sum_e\min\{z_e,1\}$.
The case $x_B=b_i=0$ is immediate. The same argument applies to a finite
global minimizing form, since off-face prices remain capped by $T$.

Sum \eqref{general:coordinate} over closed owners and use
Lemma~\ref{general:reciprocal} separately in each resource coordinate:
$C(q)\le 2dqT/(1-q)$. The endpoint integral in
\eqref{general:endpoint} is at most $T$. Since
$\int_0^a q/(1-q)\,dq=-a-\log(1-a)$, this proves
\eqref{general:dimension}. Also
\[
 -a-\log(1-a)=\sum_{k\ge2}\frac{a^k}{k}
 \le\frac{a^2}{2(1-a)},\qquad
 D_d(a)\le\frac1a+\frac{da}{1-a}.
\]
Choosing $a=1/(1+\sqrt d)$ gives $1+2\sqrt d$.
At $T=0$, monotonicity makes all coalition values, caps and endpoint
marginals vanish, completing these cases.

\subsection{Baseline, selection and remaining gap}
The grand marginal $U_i=f(c)-f(c-b_i)$ is eligible: it is own-monotone and
dominates the support allocation at $c$. Put
$\alpha_i=\max_{e:c_e>0}b_{ie}/c_e$ when $T>0$.
Then $c-b_i\ge(1-\alpha_i)c$, so $U_i\le\alpha_iT$ and
$\sum_i\alpha_i\le d$. This proves the $dT$ baseline.
Choosing an entire rule using the fixed public $n,d$ before reports proves
\eqref{general:combined}; no pointwise minimum of payment vectors is used.
For a finite $K$ representation, the complexity bound in
Theorem~\ref{prod:upper} is another available choice.

The scalar leading constant two is already sharp as $q\downarrow0$:
two equal positive weights give expected closed ratio exactly $2q$.
This does not prove optimality of the dimension constant for the full mechanism.
The unrestricted owner and dimension rates remain open between the earlier
logarithmic-over-logarithmic lower bounds and the square-root uppers above.
The lower bounds and their parameter translations are proved in Online
Resource 1, theorems titled ``Owner-count range'' and ``Concave dimension gap''.
Section~\ref{arch:section} gives matching square-root growth obstructions only for
specified fixed-mask, common-support and endpoint architectures, including
a maximum-coalition-marginal cap variant. Those restrictions are essential;
they are not lower bounds for arbitrary eligible payment rules.

Square-root payment bounds also arise in static cost-of-stability theory:
Bachrach et al.~\cite[Theorem 3.5]{Bachrach2009} and their published
version~\cite[Theorem 7]{Bachrach} give a total-funding factor $\sqrt n$ for
superadditive games without the cross-report incentive requirement. Every profile in our
model already admits an efficient support allocation, so its static cost of
stability is zero. The new bounds concern one rule linking those profiles.


\section{Owner-count and resource-dimension extensions}\label{ext:sec}
The hybrid of Section~\ref{prod:section} extends to every finite monotone,
concave, degree-one homogeneous production function, including discontinuities
at boundary faces.  Its budget can also count the selected forms rather than
the entire technology. This section retains the earlier $n/\log n$
owner-count estimate as an intermediate bound, proves the parameter
translations giving $\log n/\log\log n$ and $\log d/\log\log d$ lower bounds,
and records a dimension fallback sharp without concavity. The stronger
square-root upper bounds in Theorem~\ref{general:sqrt} are used in the
updated ranges below. The unrestricted optimal rates remain open.

Throughout, $f:\R_+^d\to\R$ is finite, coordinatewise nondecreasing and
positively homogeneous of degree one; concavity is assumed where stated.
Homogeneity forces $f(0)=0$, hence monotonicity gives $f\ge0$.
There are $n$ fixed owners $N=[n]$, with arbitrary bundles $b_i\ge0$;
$c=\sum_i b_i$, $T=f(c)$, $v_b(S)=f(\sum_{i\in S}b_i)$, and
$x_B=\sum_iB_i b_i$.  An eligible rule is nonnegative, coalition-covering,
and componentwise nondecreasing in each owner's whole report at every real
profile, as in Section~\ref{prod:section}.  All domains and owners remain fixed.
Use $\alpha_j,C_j(H)$ from \eqref{prod:constant}, with $H\ge1$ in place of $K$.

\subsection{Reusable estimates and the selected range}
\begin{lemma}[Fixed-vector Chernoff bound]\label{ext:lem:chernoff}
Let independent bits have $\Prob(B_i=0)=q\in(0,1)$ and put
$\alpha=-\log(q+(1-q)e^{-\lambda})$ for $\lambda>0$.
For every fixed $\delta\in[0,1]^n$,
\[
 \Prob\Bigl(\sum_iB_i\delta_i\le1\Bigr)
 \le e^{\lambda-\alpha\sum_i\delta_i}.
\]
If $\mathcal L\subseteq[0,1]^n$ has at most $H$ elements and a mask-dependent
choice $\delta^B\in\mathcal L$ always satisfies
$\sum_iB_i\delta_i^B\le1$, then
$\E\sum_i\delta_i^B\le(\log H+\lambda+1)/\alpha$.
For $q\le1/j$, $\lambda=2\log j$, one has $\alpha\ge\alpha_j$.
\end{lemma}
\begin{proof}
The fixed-vector exponential calculation in \eqref{prod:chernoff} uses only
$\delta_i\in[0,1]$ and independence, so applies verbatim.
The event $\sum_i\delta_i^B\ge x$ implies that some vector in the fixed
list with sum at least $x$ has sampled score at most one.
A union bound gives $\min\{1,He^{\lambda-\alpha x}\}$; its integral is
$(\log H+\lambda+1)/\alpha$.  Finally
$q+(1-q)/j^2\le(j^2+j-1)/j^3$ when $q\le1/j$.
The list must be fixed before sampling, but may depend on the full profile.
\end{proof}

\begin{lemma}[Budget from capped mask contributions]\label{ext:lem:budget}
Fix a profile and, for every mask, a vector $a(B)\in[0,T]^n$ satisfying
\begin{equation}\label{ext:eq:open-sum}
 \sum_iB_i a_i(B)=f(x_B).
\end{equation}
Define $z_i^q=\E a_i(B)$, $p_i^q=\E[a_i(B)\mid B_i=0]$, and use the endpoint
marginal $m_i^q$ and hybrid $r_i^q=qp_i^q+(1-q)m_i^q$ of
\eqref{prod:hybrid}.  With $g=\E f(x_B)$, $Z=\sum_i z_i^q$, $R=\sum_i r_i^q$,
all these functions are finite mask polynomials and
\[
 q\sum_i p_i^q=Z-g,\qquad \sum_i m_i^q=-g',\qquad R=Z-g-(1-q)g'.
\]
Thus $\int_0^{b_0}R\le\int_0^{b_0}Z+T$ for $0<b_0<1$.
If the vectors $a(B)$ take at most $H$ values, then for $0<q\le1/j$,
$Z(q)\le T(\log H+2\log j+1)/\alpha_j$, and
$j\int_0^{1/j}R(q)\,dq\le C_j(H)T$.
\end{lemma}
\begin{proof}
Condition on the own bit and sum \eqref{ext:eq:open-sum} to obtain the first
identity.  Differentiating the finite multiaffine expectation on its diagonal
gives the second.  The exact integration identity \eqref{prod:integral}
then applies, since $g(0)=T$ and $g\ge0$.
If $T>0$, the fixed list $\{a(B)/T\}$ lies in $[0,1]^n$, and its selected
score is $f(x_B)/T\le1$.  Apply Lemma~\ref{ext:lem:chernoff}, then integrate.
For $T=0$, monotonicity makes every masked value and contribution zero.
\end{proof}

\begin{proposition}[Selected-range refinement]\label{ext:prop:range}
For a finite min-linear family $V$ of size $K$, the unchanged hybrid rule
of Section~\ref{prod:section} satisfies
\[
 \sum_iM_i(b)\le\min\{n,C_j(\min\{K,2^n\})\}T.
\]
One may choose $j$ in advance using $\min\{K,2^n\}$; zero-report owners
are not removed and the selected range is not used to retune after reports.
\end{proposition}
\begin{proof}
At a fixed profile, $\{W(b,B):B\in\{0,1\}^n\}$ has at most
$\min\{K,2^n\}$ elements.  Its capped vectors satisfy
\eqref{ext:eq:open-sum}, because each nonnegative open contribution is a
summand of the selected value $f(x_B)\le T$ and hence is uncapped.
Apply Lemma~\ref{ext:lem:budget}.  Also $a_i,m_i^q\in[0,T]$, so
$r_i^q,M_i\le T$ and the total is at most $nT$.
\end{proof}

\subsection{General concave homogeneous production}
\begin{lemma}[Relative supporting forms]\label{ext:lem:support}
Assume $f$ is concave.  For $x\ge0$ let $J=\operatorname{supp}(x)$ and
$F_J=\{y\ge0:\operatorname{supp}(y)\subseteq J\}$.
If $J\ne\emptyset$, the set of $u\in\R_+^J$ satisfying
\[
 u\cdot x_J=f(x),\qquad f(y)\le u\cdot y_J\quad(y\in F_J)
\]
is nonempty, compact and convex.  Its unique minimum-norm element $u(x)$
therefore defines a choice depending only on $f,x$.
Set $\ell_x(y)=u(x)\cdot y_J$ on $F_J$ and $+\infty$ outside it.
For $x=0$, instead set $\ell_0(0)=0$, $\ell_0(y)=+\infty$ for $y\ne0$.
Then $\ell_x$ is componentwise nondecreasing, majorizes $f$ everywhere,
and touches it at $x$.  On its finite face it is an ordinary nonnegative
linear form.  No boundary continuity is required.
\end{lemma}
\begin{proof}
On the relative interior of $F_J$, finiteness and concavity supply a finite
supergradient $u$ at $x$~\cite{BoydSubgradients}:
$f(y)\le f(x)+u\cdot(y_J-x_J)$.
Putting $y=\lambda x$ with both $\lambda>1$ and $0<\lambda<1$ yields
$u\cdot x_J=f(x)$.  Monotonicity applied to $y=x+te_e$ gives $u_e\ge0$.
For any boundary point $y\in F_J$ and $0<\epsilon<1$, the point
$y_\epsilon=(1-\epsilon)y+\epsilon x$ is in the relative interior.
Concavity and the supporting inequality give
\[
 (1-\epsilon)f(y)+\epsilon f(x)\le f(y_\epsilon)
 \le(1-\epsilon)u\cdot y_J+\epsilon f(x).
\]
Cancel and divide by $1-\epsilon$ to obtain the boundary inequality.
The admissible set is an intersection of closed half-spaces and the touching
hyperplane.  It is bounded by $0\le u_e\le f(x)/x_e$, so is compact;
strict convexity of squared norm gives the unique choice.
The extension majorizes $f$ trivially outside its face.  Increasing an argument
within the face increases its nonnegative linear value; leaving the face
changes that value to infinity.  The zero-face convention has the same
properties.  In particular, no subtraction of infinite quantities is needed.
\end{proof}

\begin{definition}[Relative-support hybrid]\label{ext:def:general}
For each profile and mask choose the form $\ell_{x_B}$ above and set
$a_i(B)=\min\{\ell_{x_B}(b_i),T\}$.
Define $z_i^q,p_i^q,m_i^q,r_i^q$ as in Lemma~\ref{ext:lem:budget} and
$M_i=j\int_0^{1/j}r_i^q\,dq$, with $j\ge2$ fixed independently of reports.
The capped coefficients are finite, so these are finite sums and exact
integrals of polynomials, defining deterministic payments.
\end{definition}

\begin{theorem}[General concave homogeneous production]\label{ext:thm:general}
For every finite, monotone, concave, degree-one homogeneous $f$, the rule of
Definition~\ref{ext:def:general} is eligible at every real bundle profile and
\[
 0\le M_i(b)\le T,\qquad
 \sum_iM_i(b)\le\min\{n,C_j(2^n)\}T.
\]
A zero bundle receives zero, and every zero-value profile receives zero.
\end{theorem}
\begin{proof}
Both capped contributions and endpoint marginals belong to $[0,T]$, so the
same holds for their mixtures.  If $b_i=0$, then $\ell(0)=0$ and its on/off
aggregates coincide.  If $T=0$, monotonicity makes every masked value zero.

For incentives, fix $b_{-i},B_{-i}$ with $B_i=0$.
The aggregate $x_B$ and its selected form are independent of the entire own
bundle $b_i$.  Both $\ell_{x_B}(b_i)$ and $T$ increase weakly with $b_i$,
including finite-to-infinite transitions; their capped minimum is therefore
nondecreasing.  The off term of $m_i^q$ is fixed and the on term increases.
Averaging with report-independent nonnegative weights proves own-bundle
monotonicity on the whole domain.

For coverage, fix a selected form $\ell$ and coalition $S$.
If $\ell(b_i)\ge T$ for some $i\in S$, that capped term alone covers
$v_b(S)\le T$.  Otherwise all coalition terms are finite, uncapped and on
the same face, so their sum is $\ell(\sum_{i\in S}b_i)\ge v_b(S)$.
Thus every vector $a(B)$, and hence $z^q$, covers coalitions.
For a fixed $B_{-i}$ select $\ell=\ell_x$ at the on aggregate
$x=b_i+\sum_{h\ne i}B_hb_h$.  The bundles $b_i,x-b_i$ lie on its finite face.
Touching, majorization and finite additivity give
\begin{equation}\label{ext:eq:endpoint}
 f(x)-f(x-b_i)\ge\ell(x)-\ell(x-b_i)=\ell(b_i),\qquad
 0\le\ell(b_i)\le f(x)\le T.
\end{equation}
The middle term is the uncapped on contribution.  Averaging implies
$m_i^q\ge\E[a_i(B)\mid B_i=1]$, hence $r_i^q\ge z_i^q$ and coverage of $M$.

For the budget, every open bundle lies on the finite face of $\ell_{x_B}$.
Its nonnegative contributions sum to $\ell_{x_B}(x_B)=f(x_B)\le T$,
so \eqref{ext:eq:open-sum} holds.  At the fixed profile there are at most
$2^n$ capped vectors.  Lemma~\ref{ext:lem:budget} gives the $C_j(2^n)$ bound;
$M_i\le T$ gives the other bound.  All endpoint arithmetic above is finite.
\end{proof}

\begin{corollary}[Owner-count upper bound]\label{ext:cor:owner-upper}
Let $L_n=n\log2$ and $j_n=\max\{2,\lceil L_n/[\log(2+L_n)]^2\rceil\}$.
The rule with this report-independent parameter satisfies
\[
 \sum_iM_i(b)\le C_{j_n}(2^n)T
   =(\log2+o(1))\frac n{\log n}T,
\]
uniformly in dimension, technology and profiles.  For $n=1$ it is efficient.
\end{corollary}
\begin{proof}
Here $j_n\sim L_n/(\log L_n)^2$, $\alpha_{j_n}\sim\log n$, and
$j_n=o(n/\log n)$; substitute into $C_j$.
For $n=1$ and $b_1\ne0$, the off form is $\ell_0$, so its capped
contribution is $T$, and the endpoint marginal is also $T$.
Thus $r_1^q=M_1=T$; the zero bundle is immediate.
\end{proof}
\begin{remark}\label{ext:rem:direct}
The direct construction covers downward boundary discontinuities, for example
$f(x_1,x_2)=x_1$ when $x_2>0$ and zero otherwise.
It avoids limits of approximating technologies.  The canonical relative
support is a mathematical definition, not an efficient oracle; no regularity
or measurability of the resulting all-profile payment map is asserted.
\end{remark}

\subsection{Owner-count and dimension lower bounds}
\begin{definition}[Total-payment factors]\label{ext:def:factors}
Let $\mathcal Y_n(f)$ denote eligible rules on all of $(\R_+^d)^n$ and put
\[
 \mathfrak b_n(f)=\inf_{y\in\mathcal Y_n(f)}\sup_{b:T(b)>0}\frac{y(b)(N)}{T(b)}.
\]
The inner supremum is zero if $f\equiv0$.  These are total-payment factors,
not subsidy quantities.  Let $\mathfrak M_d$ be finite nonnegative min-linear
technologies on $\R_+^d$, and $\mathfrak F_d$ all finite monotone concave
homogeneous technologies there.  Define
\[
 B_n^{\rm own}=\sup_{d\ge1,f\in\mathfrak M_d}\mathfrak b_n(f),\quad
 B_n^{\rm conc}=\sup_{d\ge1,f\in\mathfrak F_d}\mathfrak b_n(f),\quad
 B_d^{\rm dim}=\sup_{n\ge1,f\in\mathfrak M_d}\mathfrak b_n(f).
\]
Owners, resources and technology are fixed inside each infimum; different
problems need not share a rule.  The first two suprema allow dimension to
grow; the last fixes dimension and permits arbitrarily many owners.
\end{definition}

\begin{theorem}[Owner-count range]\label{ext:thm:owner-range}
As $n\to\infty$ along all integers,
\[
 (1-o(1))\frac{\log n}{\log\log n}\le B_n^{\rm own}
 \le B_n^{\rm conc}\le n\sqrt{\frac{2}{n+1}}
       =(\sqrt2+o(1))\sqrt n.
\]
The lower bound is a parameter translation of the weighted hierarchy,
not a new construction; it permits dimension to grow with $n$.
\end{theorem}
\begin{proof}
The updated upper follows from Theorem~\ref{general:sqrt} and
$\mathfrak M_d\subseteq\mathfrak F_d$.
For the lower, use $f_k(x)=\min_{1\le e\le k}x_e$.
Restrict $km$ owners to axis-supported bundles, with $m$ owners per axis.
An eligible all-bundle rule restricts to a valid scalar-report multi-glove
rule, and its restricted total-payment factor cannot be larger.
Thus Lemma~\ref{lem:weighted-hierarchy} applies to $\mathfrak b_{km}(f_k)$.
If $n\ge n'$, restricting an $n$-owner rule to zero reports for the last
$n-n'$ owners and discarding their nonnegative payments yields an eligible
$n'$-owner rule.  All original coalition constraints and unilateral comparisons
remain, so $\mathfrak b_n(f)\ge\mathfrak b_{n'}(f)$ without population incentives.

For an integer $h\ge2$ set
\begin{align*}
 R&=\lceil\log(h+1)\rceil, &s&=\lceil\sqrt R\rceil, &m&=hs,\\
 \alpha&=1/s, &A&=(h+1)m-h^2, &\lambda&=As,\\
 Z&=\sum_{r=0}^h\frac{\lambda^r}{r!}, &K_h&=\lceil2RZ\rceil, &\nu_h&=K_hm.
\end{align*}
Here $s\ge2$, so $m>h$, and $K_hp_0=K_h/Z\ge2R\ge2\log(h+1)$.
Equation~\eqref{eq:weighted-bound}, with hierarchy height $h$, yields
\begin{equation}\label{ext:eq:owner-lower-finite}
 \mathfrak b_{\nu_h}(f_{K_h})
 \ge\frac{h+1-h/s}{(1+1/s)(1+h/(h+1)^2)}=(1-o(1))h.
\end{equation}
Also $\log\lambda=2\log h+O(\log\log h)$ and $\lambda\ge h$.
The terms of $Z$ increase up to index $h$, so
$\lambda^h/h!\le Z\le(h+1)\lambda^h/h!$.
Stirling's formula gives $\log Z=(1+o(1))h\log h$ and hence
$\log\nu_h=(1+o(1))h\log h$.
For fixed $0<\epsilon<1$ choose
$h=\lfloor(1-\epsilon)\log n/\log\log n\rfloor$.
Then $\log\nu_h\le(1-\epsilon+o(1))\log n<\log n$ for all large $n$.
Padding to $n$ owners and \eqref{ext:eq:owner-lower-finite} give
$B_n^{\rm own}\ge(1-\epsilon-o(1))\log n/\log\log n$.
Taking the liminf and then $\epsilon\downarrow0$ proves the all-integer lower.
\end{proof}

\subsection{Marginal floors for the original core}
At a profile define the original efficient core and its coordinate benchmarks
by $\core(b)=\{y:y(N)=T,\ y(S)\ge v_b(S)\ \forall S\}$ and
$H_i(b)=\max_{y\in\core(b)}y_i$.  Put
\begin{equation}\label{ext:eq:marginal}
 U_i=T-f(c-b_i),\qquad
 \rho_i=\max_{e:c_e>0}b_{ie}/c_e\quad(\rho_i=0\text{ if }c=0).
\end{equation}

\begin{lemma}[Marginal contributions]\label{ext:lem:marginal}
For concave $f$, the core is nonempty and compact,
$0\le H_i\le U_i\le T$, and $U$ is eligible.  Furthermore
\[
 \sum_iU_i\le\min\{d,n\}T,\qquad
 \sum_iU_i\le\min\{K,d,n\}T\quad\text{for a finite family of $K$ forms}.
\]
\end{lemma}
\begin{proof}
The grand relative support gives the finite allocation $y_i=\ell_c(b_i)$.
It is nonnegative, sums to $\ell_c(c)=T$, and covers each coalition by
majorization.  Singleton constraints make the core a closed subset of
$[0,T]^n$, hence compact.  Its complementary-coalition constraint gives
$y_i\le T-v_b(N\setminus\{i\})=U_i$.  Thus $U$ dominates a covering core
allocation.  Also $0\le U_i\le T$ and its subtracted term omits the whole
own bundle, proving own-report monotonicity.
Since $c-b_i\ge(1-\rho_i)c$, monotonicity and homogeneity imply
$U_i\le\rho_iT$.  Moreover
\[
 \sum_i\rho_i\le\sum_{e:c_e>0}\sum_i b_{ie}/c_e
   =|\operatorname{supp}(c)|\le d,\qquad \rho_i\le1.
\]
For finite $V$, group owners by a chosen off minimizer $w$ of $w\cdot(c-b_i)$.
For a nonempty group $G$, writing $A=w\cdot c\ge T$ gives
\[
 \sum_{i\in G}U_i=|G|(T-A)+w\cdot\sum_{i\in G}b_i
 \le|G|T-(|G|-1)A\le T.
\]
There are at most $\min\{K,n\}$ groups.  This last cancellation is used only
for ordinary finite forms, never for extended forms taking infinity at $c$.
\end{proof}

\begin{proposition}[Original-core floors]\label{ext:prop:floor}
For concave $f$, let a fixed $D\ge1$ satisfy $\sum_iU_i\le DT$ at every
profile.  One may use $D=d$, $D=n$, or $D=K$ for finite min-linear $f$.
For any eligible $y$, the rule $y_i'=y_i+U_i/D$ is eligible, protects
$y_i'\ge H_i/D$, and costs at most $y(N)+T$ in total.
With $D=K$ this preserves the sharp leading coefficient of
Section~\ref{prod:section}.  The benchmarks use the original core.
\end{proposition}
\begin{proof}
Adding nonnegative own-report monotone payments preserves eligibility.
The inequalities $U_i\ge H_i$ and $\sum_iU_i\le DT$ give the assertions.
One additional factor of $T$ is lower order than $T\log K/\log\log K$.
\end{proof}

\subsection{Resource dimension and technology-aware fallbacks}
\begin{theorem}[Dimension fallback without concavity]\label{ext:thm:dimension}
For every finite monotone homogeneous $f$, even without concavity,
$P_i=T\rho_i$ is eligible and
\[
 \sum_iP_i\le\min\{n,d_+(c)\}T\le dT,
 \qquad d_+(c)=|\operatorname{supp}(c)|.
\]
Thus $\mathfrak b_n(f)\le\min\{n,d\}$.
\end{theorem}
\begin{proof}
With other reports fixed, every ratio is $a/(a+A)$, where $A\ge0$ is fixed
and $a$ increases with the corresponding own coordinate; use $0/0=0$.
If $A=0$, the ratio jumps from zero to one and remains monotone.
Thus both $\rho_i$ and $T$ increase weakly with $b_i$, and their nonnegative
product does too.  For a coalition put $r_S=\sum_{i\in S}\rho_i$.
The coordinate inequalities $b_i\le\rho_ic$ imply
$v_b(S)\le f(r_Sc)=r_ST=\sum_{i\in S}P_i$.
If $r_S=0$ this uses $f(0)=0$; at $T=0$ all worths vanish.
The budget is $\sum_i\rho_i\le\min\{n,d_+(c)\}$, proved above without
concavity.  Zero aggregate coordinates contribute nothing.
\end{proof}
The maximum normalized coordinate share is familiar from multi-resource
allocation~\cite{Ghodsi}.  That literature's demand-report allocation model
and strategy-proofness claim differ from these supplier-side payments.

\begin{proposition}[Exact dimension factor without concavity]\label{ext:prop:sharp-d}
For every dimension $d$, the worst optimal total-payment factor over all
finite monotone homogeneous technologies and all owner counts is exactly $d$.
This does not establish a sharp dimension factor for concave technologies.
\end{proposition}
\begin{proof}
Take $f(x)=\max_ex_e$ and $d$ unit-basis owner bundles.
Then $T=1$ and every singleton has value one, forcing payment at least $d$.
The upper is Theorem~\ref{ext:thm:dimension}.
The example persists with strictly positive bundles: for $0<\epsilon\le1$,
let $b_i=e_i+\epsilon\sum_{e\ne i}e_e$.  Every singleton still has value one,
while $T=1+(d-1)\epsilon$, forcing factor at least
$d/[1+(d-1)\epsilon]$.  With $\epsilon=d^{-2}$ this is $(1-o(1))d$.
For $d\ge2$ this technology is convex and not concave, so the argument does
not apply to $\mathfrak M_d$ or $\mathfrak F_d$.
\end{proof}

\begin{theorem}[Concave dimension gap]\label{ext:thm:dim-gap}
As $d\to\infty$ along all integers,
\[
 (1-o(1))\frac{\log d}{\log\log d}\le B_d^{\rm dim}
 \le\min\{d,1+2\sqrt d\}.
\]
The lower already holds for $f(x)=\min_ex_e$ with axis-supported owners.
The new $O(\sqrt d)$ upper holds for all $\mathfrak F_d$; its unrestricted
optimal order is not determined here.
\end{theorem}
\begin{proof}
The updated upper is Theorem~\ref{general:sqrt}.
Use the parameters $R,s,m,\alpha,A,\lambda,Z,K_h$ from
Theorem~\ref{ext:thm:owner-range}.  For every integer $D\ge K_h$, the weighted
hierarchy with $k=D$ pools still has $Dp_0=D/Z\ge2R$.
The same bound \eqref{ext:eq:owner-lower-finite} therefore yields
$\mathfrak b_{Dm}(f_D)\ge(1-o(1))h$; the owner count $Dm$ is allowed to grow
in $B_D^{\rm dim}$.  Since $\log K_h=(1+o(1))h\log h$, choosing
$h=\lfloor(1-\epsilon)\log D/\log\log D\rfloor$ gives $D\ge K_h$ for all
large $D$.  The resulting liminf is at least $1-\epsilon$;
let $\epsilon\downarrow0$.
\end{proof}

\begin{proposition}[Technology-aware refinements]\label{ext:prop:dual-share}
For finite min-linear $f$ define, with $0/0=0$,
\[
 \theta_i=\max_{w\in V}\frac{w\cdot b_i}{w\cdot c},\qquad Q_i=T\theta_i.
\]
The rule $Q$ is eligible, $\theta_i\le\rho_i$, and
$\sum_iQ_i\le\min\{n,d_+(c)\}T$.  For $|V|=1$ it is efficient.
If every $w\in V$ lies in the nonnegative cone of fixed vectors
$a_1,\ldots,a_r\in\R_+^d$, where $r\ge1$ is an integer, then
\[
 \theta_i\le\sigma_i:=\max_{1\le l\le r}\frac{a_l\cdot b_i}{a_l\cdot c},
 \qquad\sum_i\theta_i\le\sum_i\sigma_i\le r,
\]
so the budget improves to $\min\{n,d_+(c),r\}T$.
The marginal rule also satisfies $U_i\le P_i=T\rho_i$, with total
$\min\{K,d,n\}T$.
\end{proposition}
\begin{proof}
For fixed $w$, its ratio has the same $a/(a+A)$ monotonicity as above.
Taking the maximum and multiplying by monotone nonnegative $T$ preserves it.
For $T>0$, choose a form $w^*$ with $w^*\cdot c=T$.  Then
\[
 v_b(S)\le w^*\cdot\sum_{i\in S}b_i
     =T\sum_{i\in S}\frac{w^*\cdot b_i}{w^*\cdot c}
     \le\sum_{i\in S}Q_i.
\]
At $T=0$ all payments and worths vanish.
Each nonzero-denominator dual ratio is a weighted average of the coordinate
shares $b_{ie}/c_e$, using weights $w_ec_e$, so is at most $\rho_i$.
For one form its shares sum to one at $T>0$, proving efficiency.
If $w=\sum_l\mu_la_l$ with $\mu_l\ge0$, its ratio is similarly a weighted
average of generator ratios, with weights $\mu_l(a_l\cdot c)$.
Zero denominators also have zero numerators.  Finally
$\sum_i\sigma_i\le\sum_l\sum_i(a_l\cdot b_i)/(a_l\cdot c)\le r$.
The marginal claims are Lemma~\ref{ext:lem:marginal}.
No bound in terms of ordinary signed linear rank is asserted.
\end{proof}

\begin{remark}[Rule limitations and ex ante choice]\label{ext:rem:limits}
For $f(x)=\sum_ex_e$ and unit-basis owners, $P$ costs $dT$ whereas $Q$ is
efficient.  For $f(x)=\min_ex_e$ at that profile, $U$ costs $dT$ although
an efficient core allocation exists.  These are limitations of particular
rules, not lower bounds on all eligible rules.
Choosing among the valid rules in advance from fixed $n,d,K$ gives factor
$\min\{n,d,C_j(2^n)\}$ for general concave $f$, and
$\min\{n,d,C_j(\min\{K,2^n\})\}$ for finite min-linear $f$.
Coordinatewise minima or report-dependent switching are not asserted valid.
The convex-analysis, marginal/core, normalized-share, and Chernoff ingredients
are established tools; no historical priority, sharp owner/dimension rate,
or efficient evaluation for arbitrary $f$ is claimed.
\end{remark}

% Insert as a new supplementary section before the bibliography.
% Uses only the existing amsmath/amssymb/amsthm and theorem environments.
% All labels introduced here have the unused prefix arch:.
\section{Limits of fixed-mask common-support mechanisms}\label{arch:section}

This section proves lower bounds for a specified payment architecture.
The owners and technology remain fixed throughout each construction, and only
their reports change. The conclusions do not apply to every eligible mechanism.
In particular, they do not replace the unrestricted logarithmic lower bounds
by square-root lower bounds.

\subsection{The architecture, marginal caps, and profile averaging}

Fix $n$ owners and a continuous finite nonnegative min-linear technology
$F:\R_+^d\to\R_+$: thus $F(x)=\min_{w\in\mathcal W}w\cdot x$ for a fixed
nonempty finite family $\mathcal W\subseteq\R_+^d$.
For a report profile $b=(b_i)_{i=1}^n$, put
\[
 b(S)=\sum_{i\in S}b_i,\qquad x_B=b(B),\qquad
 T=F(b([n])),\qquad v_b(S)=F(b(S)).
\]
A \emph{global linear support at $x$} is a nonnegative finite linear form
$\ell$ such that $\ell(z)\ge F(z)$ for all $z\ge0$ and $\ell(x)=F(x)$.
The support need not belong to the displayed defining family $\mathcal W$.

\begin{definition}[Fixed-mask common-support architecture]\label{arch:def}
Fix any probability law on subsets $B\subseteq[n]$, independently of all
reports. Conditional on $B$, select a global linear support $\ell_B$ at $x_B$
using only the fixed technology, $B$, and the open reports $(b_i)_{i\in B}$.
The selection may randomize. A closed report may not enter that selection.
Choose one of the following two caps:
\begin{align}
 c_i(b)&=T,\label{arch:tcap}\\
 c_i(b)&=V_i(b):=
 \max_{S\subseteq[n]\setminus\{i\}}
       \{F(b_i+b(S))-F(b(S))\}.\label{arch:vcap}
\end{align}
For the chosen cap, the payment map is the expectation of
\begin{equation}\label{arch:maskpay}
 Y_i(B,b)=
 \begin{cases}
 \min\{\ell_B(b_i),c_i(b)\},&i\notin B,\\
 F(x_B)-F(x_B-b_i),&i\in B.
 \end{cases}
\end{equation}
The cap convention and all mask and support-selection laws are fixed before
the profile is evaluated.
\end{definition}

This definition permits arbitrary correlations and asymmetries in the mask
law, any fixed mixture of mask laws, and randomized or redundant supporting
forms. Conditional on $B=\varnothing$, it permits exactly one support law for
the fixed technology and owners: there are no open reports on which that law
could depend. It does not permit a different empty-mask support law for each
profile.

\begin{lemma}[Eligibility and simultaneous marginal truncation]
\label{arch:caplemma}
Both architectures in Definition~\ref{arch:def} are eligible: payments are
nonnegative, cover every coalition, and are nondecreasing in each owner's
own whole bundle with all other reports fixed. Moreover, if a nonnegative
vector $z$ covers every coalition at $b$, then
$(\min\{z_i,V_i(b)\})_{i=1}^n$ also covers every coalition.
\end{lemma}
\begin{proof}
For every $i$, monotonicity of $F$ gives $0\le V_i\le T$. Each term in
the maximum defining $V_i$ is nondecreasing in $b_i$, since its off aggregate
does not involve $i$. Thus $V_i$ is own-bundle monotone.

For simultaneous truncation, fix a coalition $S$ and write
$C=\{i\in S:z_i>V_i\}$ and $R=S\setminus C$.
Coverage gives $\sum_{i\in R}z_i\ge v_b(R)$. Add the members of $C$ to $R$
one at a time. Each increment of coalition value is at most the corresponding
$V_i$, so
\[
 v_b(S)-v_b(R)\le\sum_{i\in C}V_i.
\]
Adding proves $\sum_{i\in S}\min\{z_i,V_i\}\ge v_b(S)$.

For any global support $\ell$, the vector
$A_i=\min\{\ell(b_i),T\}$ covers every coalition. If a member of a coalition
has raw price at least $T$, its capped payment alone covers that coalition;
otherwise the uncapped sum equals $\ell(b(S))\ge F(b(S))$.
The preceding truncation argument shows the same fact for
$\min\{\ell(b_i),V_i\}$. For an open owner $i$, touching and majorization
give
\[
 F(x_B)-F(x_B-b_i)\ge\ell_B(b_i),
\]
and this endpoint marginal is at most both $T$ and $V_i$.
Replacing an open owner's capped support payment by its endpoint marginal
therefore preserves coverage. All payments are nonnegative.

For a closed owner, the selected support is independent of its report, while
its raw price and either cap are nondecreasing in that report. For an open
owner, the off aggregate $x_B-b_i$ is independent of its report and the on
value is nondecreasing. Hence every fixed-mask payment is own-bundle
monotone. Averaging over fixed mask and support randomness preserves the
three requirements.
\end{proof}

\begin{lemma}[Combining regular and empty-mask witnesses]
\label{arch:averaging}
Fix one technology and $n$ owners. Suppose a fixed finite distribution
$\mu$ of positive-value profiles has mean normalized total payment at least
$C>0$ for every nonempty mask and every admissible choice of supports at
those profiles. Suppose a fixed finite distribution $\nu$ of positive-value
profiles has mean normalized empty-mask total payment at least $A>0$ for
every one fixed empty-mask support. Then every rule in the architecture has
some profile with total-payment factor at least
\begin{equation}\label{arch:harmonic}
 \frac{AC}{A+C}.
\end{equation}
The same conclusion holds if the empty support is randomized.
\end{lemma}
\begin{proof}
Give total weight $A/(A+C)$ to $\mu$ and $C/(A+C)$ to $\nu$.
For every nonempty mask the regular contribution alone is at least
$AC/(A+C)$; for the empty mask the anchor contribution alone is at least
that number. Other payments are nonnegative. At the empty mask, the finite
anchor average commutes with integration against its one fixed support law.
The regular bound is valid for every support choice separately, so it also
survives profile-dependent support selection using open reports. Finally
average over the report-independent mask law. The resulting finite weighted
average of normalized expected payments is at least \eqref{arch:harmonic};
at least one profile attains that bound.
\end{proof}

\subsection{A finite technology forcing all nonempty-mask prices}

The following lemma supplies the common finite construction. It is useful
because every supporting form at the regular profile is forced on its active
resource block, regardless of the chosen support-selection rule.

\begin{lemma}[Forced-support realization]\label{arch:realization}
Let $n\ge5$, and let $a_1\ge\cdots\ge a_{n-1}\ge0$ satisfy
\begin{equation}\label{arch:gapassumption}
 v_l:=(n-l)a_l\le\frac{n-l}{n}-\frac1{n^3}
 \qquad(1\le l<n).
\end{equation}
Set $v_0=1$, $v_n=0$, $\delta_l=v_l-v_{l+1}$ for $0\le l<n$, and
\[
 \varepsilon=\frac1{100n^6},\qquad
 t=1+(n-1)\varepsilon,\qquad \eta=\frac1{n^3}.
\]
For every nonempty proper subset $S\subset[n]$ define the form
\[
 w^S_j=
 \begin{cases}n^2/|S|^2,&j\in S,\\a_{|S|},&j\notin S,
 \end{cases}
 \qquad
 f(x)=\min\left\{\frac1n\sum_jx_j,
              w^S\cdot x:\varnothing\ne S\subsetneq[n]\right\}.
\]
At the regular profile
\begin{equation}\label{arch:regularprofile}
 b_i=e_i+\varepsilon\sum_{j\ne i}e_j,
\end{equation}
grand production is $T=t$. If $B$ is nonempty and proper, write
$C=[n]\setminus B$, $l=|C|$ and $r=n-l$. The unique minimizing form at
$x_B$ is $w^C$. Consequently every global linear support at $x_B$ is
this same form, and every closed owner's raw price exceeds $T$.
For every nonempty mask, including $B=[n]$, each open endpoint marginal
divided by $T$ is at least $\delta_l-\eta$.

These conclusions remain valid on the active block if $f$ is one summand
of a sum of technologies on disjoint resource blocks and the other reported
blocks are all zero.
\end{lemma}
\begin{proof}
The function $f$ is continuous, finite min-linear, nonnegative, monotone,
concave, and homogeneous of degree one. The full aggregate is $t\mathbf1$.
Its baseline value is $t$, whereas a subset form contributes at least
$n^2t/|S|>t$, so the baseline is uniquely minimal there.

At a nonempty proper mask, the aggregate has coordinate $r\varepsilon$ on
$C$ and $1+(r-1)\varepsilon$ elsewhere. The value of $w^C$ is
\begin{equation}\label{arch:maskedvalue}
 v_l+\varepsilon E_l,
 \qquad E_l=\frac{n^2r}{l}+(r-1)v_l,
 \qquad 0\le E_l\le2n^3.
\end{equation}
By \eqref{arch:gapassumption} and $2n^3\varepsilon<1/n^3$, this value is
strictly below $r/n$ and hence below the baseline $rt/n$; it is also below 1.
Any subset $S$ meeting an open coordinate has value at least
$n^2/(n-1)^2>1$. If instead $\varnothing\ne S\subsetneq C$, its contribution
from the coordinates in $S$ is $n^2r\varepsilon/|S|$, strictly larger than
$n^2r\varepsilon/l$. Its off-$S$ coefficient is at least $a_l$; it multiplies
all open coordinates and additional nonnegative closed coordinates. Thus it
also costs strictly more than $w^C$. These cases exhaust the competitors.

Every coordinate of $x_B$ is positive. A finite minimum with a unique
minimizing linear form is differentiable there. If any other global linear
form majorizes $f$ and touches there, taking directional derivatives in both
directions forces its coefficients to equal that gradient. This proves the
assertion for all supports, including randomized selections. Each closed
owner has raw price at least
\[
 \frac{n^2}{l^2}\bigl(1+(l-1)\varepsilon\bigr)
 \ge\frac{n^2}{(n-1)^2}>t.
\]

Put $E_0=n-1$ and $E_n=0$. The exact masked values for all mask sizes are
$v_l+\varepsilon E_l$. Since $a_l$ is nonincreasing and nonnegative,
$(v_l)_{l=0}^n$ is nonincreasing and lies in $[0,1]$, so
$0\le\delta_l\le1$. Removing an open owner changes $l$ to $l+1$; its
unnormalized endpoint marginal is at least $\delta_l-2n^3\varepsilon$.
Normalization by $t$ loses at most an additional $(n-1)\varepsilon$:
\[
 \frac{F(x_B)-F(x_B-b_i)}{T}
 \ge\delta_l-(2n^3+n)\varepsilon
 \ge\delta_l-\eta.
\]
The same computation covers $l=0$ and $l=n-1$. Endpoint marginals are
nonnegative even when this displayed lower bound is negative.

For a sum of technologies on separate blocks, restriction of a global
support to the active block majorizes that summand and touches its value
at the active aggregate. The same interior-gradient argument therefore
forces precisely the support just computed. Inactive blocks contribute
zero to every reported bundle and every endpoint value.
\end{proof}

\begin{lemma}[Realizing decreasing endpoint increments]
\label{arch:increments}
Let $n\ge5$ and let $\delta_0\ge\cdots\ge\delta_{n-1}\ge0$ sum to 1, with
$\delta_0\ge1/n+1/n^2$. Put
$v_l=\sum_{j=l}^{n-1}\delta_j$ and $a_l=v_l/(n-l)$ for $1\le l<n$.
Then the assumptions of Lemma~\ref{arch:realization} hold. At its regular
profile, every owner also satisfies $V_i/T\ge\delta_0-\eta$.
\end{lemma}
\begin{proof}
Since $v_{l+1}\le(n-l-1)\delta_l$, we have
$v_l/(n-l)\ge v_{l+1}/(n-l-1)$ whenever $l<n-1$.
Thus $a_l$ is nonincreasing. In particular
$a_l\le a_1=(1-\delta_0)/(n-1)$, whence
\[
 \frac{n-l}{n}-v_l
 \ge\frac{(n-l)(\delta_0-1/n)}{n-1}
 \ge\frac1{n^3}.
\]
Finally, $V_i$ includes the grand-coalition endpoint marginal. The $l=0$
case of Lemma~\ref{arch:realization} bounds it below by
$(\delta_0-\eta)T$.
\end{proof}

\subsection{A sharp owner obstruction already in dimension $d=n$}

\begin{theorem}[Global-cap owner and dimension barriers]
\label{arch:tbarrier}
For each $n\ge9$, let $m$ be the largest integer with $1\le m<n$ such that
\[
 S_m:=\sum_{l=0}^{m-1}\frac{m-l}{n-l}\le1,
 \qquad C_n=m-\frac1{n^2},\qquad A_n=n-2m.
\]
There is one fixed $n$-resource finite nonnegative min-linear technology on
which every global-$T$-cap rule of Definition~\ref{arch:def}, with the same
$n$ owners on all report profiles, has worst-profile factor at least
\begin{equation}\label{arch:tfinite}
 L_n=\frac{C_nA_n}{C_n+A_n}
      =(\sqrt2-o(1))\sqrt n.
\end{equation}
Since $d=n$, this is also a $(\sqrt2-o(1))\sqrt d$ architecture lower bound.
\end{theorem}
\begin{proof}
For $n\ge9$, $S_2<1$, so $m\ge2$. Also
$S_m\ge m(m+1)/(2n)$, implying $m<\sqrt{2n}<n/2$.
Define
\[
 \delta_l=\frac{m-l}{(n-l)S_m}\quad(0\le l<m),
 \qquad\delta_l=0\quad(m\le l<n).
\]
These increments sum to 1 and are nonincreasing, because
$(m-l)/(n-l)$ decreases in $l$ when $m<n$. Moreover
$\delta_0=m/(nS_m)\ge2/n\ge1/n+1/n^2$.
Use Lemma~\ref{arch:increments} to define one technology $f$ and its
regular profile $P$. All closed raw prices at nonempty masks exceed $T$.
For a mask with $0\le l<m$ closed owners, its normalized total payment is
at least
\[
 l+(n-l)(\delta_l-\eta)
 =l+\frac{m-l}{S_m}-(n-l)\eta
 \ge m-\frac1{n^2}=C_n.
\]
The case $l=0$ is included. When $m\le l<n$, closed payments alone give
at least $l\ge m$. Thus every nonempty mask costs at least $C_n$ at $P$.

It remains to control the empty mask in this same technology. Index
coordinates cyclically modulo $n$. For each $j$, let
$S_j=\{j,j+1,\ldots,j+m-1\}$ and let $z_j$ have coordinate
$\tau=m/(4n^3)$ on $S_j$ and 1 elsewhere. At the anchor profile $Q_j$, all
$n$ owners report this same bundle $z_j$. Since $a_m=0$, the form indexed
by $S_j$ gives
\[
 0<T_j=nf(z_j)\le n\frac{n^2\tau}{m}=\frac14.
\]
Strict positivity follows because $z_j$ is strictly positive and every
defining form has a positive coefficient.

For every empty-mask global support $w$, write $s=\sum_e w_e$.
We have $s\ge f(\mathbf1)=1$. Each coordinate lies in exactly $m$ cyclic
windows, so the uniform average of $w(S_j)$ is $ms/n$. Markov's inequality
shows that a fraction at least $1-2m/n$ of windows satisfies
$w(S_j)\le s/2$. At each such anchor,
\[
 w\cdot z_j\ge w([n]\setminus S_j)\ge s/2\ge1/2>T_j.
\]
All $n$ owners are therefore paid $T_j$ at the empty mask. The uniform
anchor average of normalized total payment is at least $n-2m=A_n>0$.
This is a pointwise statement for each one empty support, so it also holds
after randomizing that support according to its fixed law.
Lemma~\ref{arch:averaging} now proves \eqref{arch:tfinite} using only
$P,Q_1,\ldots,Q_n$.

Finally, uniformly for $m=O(\sqrt n)$,
\[
 S_m=\frac{m(m+1)}{2n}+O\left(\frac{m^3}{n^2}\right).
\]
Maximality of $m$ and this expansion give $m\sim\sqrt{2n}$.
Consequently $C_n\sim\sqrt{2n}$ and $A_n\sim n$, proving the stated
asymptotic factor.
\end{proof}

\subsection{Maximum-marginal caps retain the sharp owner obstruction}

The cap $V_i$ is a valid improvement to the payment rule, by
Lemma~\ref{arch:caplemma}. It can improve individual technologies. The next
result shows that it cannot improve the leading owner-count constant of the
fixed-mask architecture over general finite min-linear technologies.

\begin{theorem}[Sharp owner barrier with maximum-marginal caps]
\label{arch:vowner}
For $n\ge5$, let $m$ be the largest positive integer satisfying
$m(m+7)/2\le n-1$, and put
\[
 a=1-\frac{2m+1}{n},\qquad A_0=n-2,
 \qquad L'_n=\frac{amA_0}{am+A_0}.
\]
There is one fixed finite nonnegative min-linear technology on
$nm+n^2$ resources such that every $V_i$-cap rule of
Definition~\ref{arch:def}, with $n$ fixed owners, has worst-profile
total-payment factor at least
\begin{equation}\label{arch:vownerfactor}
 L'_n=(\sqrt2-o(1))\sqrt n.
\end{equation}
\end{theorem}
\begin{proof}
For $k=1,\ldots,m$, use the construction of Lemma~\ref{arch:realization}
with coefficients
\[
 a_{l,k}=\begin{cases}(n-l)/n^2,&l<k,\\0,&l\ge k.
 \end{cases}
\]
They are nonincreasing. If $l<k$, the gap in
\eqref{arch:gapassumption} is $(n-l)l/n^2\ge1/n^3$; if $l\ge k$, the
gap is $(n-l)/n\ge1/n^3$. Denote the resulting technology by $f_k$.
On separate coordinate blocks define
\begin{equation}\label{arch:ownertechnology}
 F(x^{(1)},\ldots,x^{(m)},z)
   =\sum_{k=1}^m f_k(x^{(k)})+\min_{1\le e\le n^2}z_e.
\end{equation}
A sum of finite minima on disjoint blocks is itself a finite minimum:
expand the finite product of form choices. Thus this is one fixed
finite nonnegative min-linear technology.

At profile $P_k$, activate the regular profile
\eqref{arch:regularprofile} only on block $k$, setting every other reported
block to zero. Grand production is $t$. At every nonempty proper mask,
Lemma~\ref{arch:realization} forces the active-block support and gives
closed raw prices exceeding $T=t$.
At the threshold $l=k-1$, with $r=n-l$, the unperturbed on value is
$r^2/n^2$ and the off value is zero. This also holds for $k=1$, because
the full-mask unperturbed value is 1. Thus each threshold endpoint divided
by $T$ is at least $r^2/n^2-\eta$.
The total normalized open payment at that threshold is at least
\[
 r\left(\frac{r^2}{n^2}-\eta\right)
 \ge n-3l-\frac1{n^2}
 \ge n-3m-1=:D.
\]
Here $(n-l)^3/n^2\ge n-3l$.

Every owner has a coalition witnessing this same threshold marginal:
choose $n-k$ of the other owners and add the owner in question. Symmetry
of $P_k$ makes the value independent of those labels. Consequently
\[
 \frac{V_i}{T}\ge\frac{(n-k+1)^2}{n^2}-\eta
 \ge1-\frac{2(k-1)}n-\eta
 \ge1-\frac{2m+1}n=a>0.
\]
The last positivity follows from the defining inequality for $m$.
Since each closed raw price exceeds $T\ge V_i$, every closed owner
receives $V_i$, hence at least $aT$.

Let $S_0=m(m+1)/2$ and put
$\pi_k=(m-k+1)/S_0$ on the regular profiles. These weights sum to 1,
and the choice of $m$ gives $S_0\le D$. For a mask with $0\le l<m$,
the mean normalized total payment across $P_k$ is at least
\[
 al+D\pi_{l+1}
 =al+D\frac{m-l}{S_0}
 \ge al+(m-l)\ge am.
\]
For $m\le l<n$, closed payments alone give at least $al\ge am$.
Thus every nonempty mask, including the full mask, has regular-profile
average at least $am$.

For empty-mask anchors, index the $n^2$ Leontief coordinates cyclically.
For each shift $j$, declare the $n$ coordinates
$S_j=\{j,\ldots,j+n-1\}$ scarce. Owner $i$ holds one unit of coordinate
$j+i-1$, zero of every other scarce coordinate, and two units of every
coordinate outside $S_j$. All other technology blocks are zero.
The grand production is 1. Removing any owner eliminates its unique scarce
resource, reducing production to zero. Hence $V_i=T=1$ for every owner.

For an empty support, let $w$ be its restriction to the Leontief block.
Its nonnegative coefficient sum $s$ is at least 1. Each coordinate lies in
exactly $n$ of the $n^2$ scarce windows, so the uniform mean scarce weight
is $s/n$. Markov's inequality gives
\[
 \Pr_j\{w(S_j)>s/2\}\le2/n.
\]
On every other anchor, each owner's abundant resources have raw price at
least $2(s/2)\ge1$. Its cap is 1, so all $n$ owners receive 1.
The mean normalized empty-mask payment is therefore at least $n-2=A_0$.
This holds for every realization of the one fixed empty-support law.
Lemma~\ref{arch:averaging}, with $C=am$, proves the finite lower bound.

Finally $m=\sqrt{2n}+O(1)$, $am=m-O(1)$, and $A_0=n-2$.
These give \eqref{arch:vownerfactor}.
\end{proof}

\subsection{A resource-dimension obstruction with maximum-marginal caps}

\begin{theorem}[Maximum-marginal dimension barrier]
\label{arch:vdimension}
For each $n\ge9$, put
\[
 m=\lfloor\sqrt{2n}\rfloor,\quad
 H=\sum_{l=1}^{m-1}\frac{m-l}{n-l},\quad a=\frac1{1+H},
 \quad R=\lceil n^{3/4}\rceil,\quad D=n+R,
\]
and define
\[
 C'_n=am-\frac2{n^2},\qquad A'_n=\frac{nR}{n+R}.
\]
There is one fixed finite nonnegative min-linear technology of dimension
$d_n=2n+R$ such that every $V_i$-cap rule of
Definition~\ref{arch:def}, with the same $n$ owners at every profile, has
worst-profile factor at least
\begin{equation}\label{arch:vdimensionfactor}
 \frac{C'_nA'_n}{C'_n+A'_n}
       =(1/2-o(1))\sqrt{d_n}.
\end{equation}
\end{theorem}
\begin{proof}
Set
\[
 \delta_0=a,\qquad
 \delta_l=a\frac{m-l}{n-l}\quad(1\le l<m),\qquad
 \delta_l=0\quad(m\le l<n).
\]
The increments sum to 1 and are nonincreasing. Since $H\le m-1$,
$a\ge1/m$. Since $n\ge9$ and $m\le\sqrt{2n}\le n/2$, we have
\[
 a-\frac1n\ge\frac{n-m}{mn}\ge\frac1{n^2}.
\]
Lemma~\ref{arch:increments} therefore supplies a regular technology $f$.
Add a separate Leontief block:
\[
 F(x,z)=f(x)+\min_{1\le e\le D}z_e.
\]
This is a single finite min-linear technology on $n+D=d_n$ resources.

At the regular profile $P$, activate only the $f$ block. Each closed
raw price at a nonempty mask exceeds $T$, while
$V_i/T\ge a-\eta$. If $1\le l<m$, normalized total payment is at least
\[
 l(a-\eta)+(n-l)(\delta_l-\eta)
 =am-n\eta=am-\frac1{n^2}.
\]
For $m\le l<n$, closed payments alone give
$l(a-\eta)\ge am-1/n^2$. For $l=0$, endpoint payments give
$n(a-\eta)\ge am-1/n^2$. Thus every nonempty mask costs at least the
slightly weaker positive quantity $C'_n=am-2/n^2$.
Positivity follows from $am\ge1$.

For the empty-mask anchors, index the Leontief resources cyclically modulo
$D$. At shift $j$, the $n$ coordinates $S_j=\{j,\ldots,j+n-1\}$ are
scarce. Owner $i$ holds one unit of coordinate $j+i-1$, zero of every
other scarce coordinate, and $D$ units of every abundant coordinate.
All reports on the $f$ block are zero. Again $T=V_i=1$ for every owner.

The restriction $w$ of an empty support to the Leontief block has
nonnegative coefficients summing to at least 1. Some coordinate $e$ has
$w_e\ge1/D$. It is abundant in exactly $R$ of the $D$ cyclic anchors.
On each of these anchors, every owner's raw price is at least
$D w_e\ge1$, so all $n$ capped payments equal 1.
The uniform anchor average is at least $nR/D=A'_n$.
Although the coordinate $e$ can depend on the support realization, the
anchor list and its weights do not. The inequality thus holds pointwise
before averaging the fixed empty-support law.
Lemma~\ref{arch:averaging} proves the finite bound in
\eqref{arch:vdimensionfactor}.

For the asymptotic statement, $m\sim\sqrt{2n}$ and
\[
 H=\frac{m(m-1)}{2n}+O\left(\frac{m^3}{n^2}\right)\longrightarrow1.
\]
Hence $a\to 1/2$ and $C'_n\sim\sqrt{n/2}$. Also
$A'_n\sim n^{3/4}$ and $d_n\sim2n$, yielding
\eqref{arch:vdimensionfactor}.
\end{proof}

\subsection{Scope of the matching statements}

The global owner-count upper factor $n\sqrt{2/(n+1)}$ is realized within
the fixed-mask support/endpoint architecture. Replacing its global cap by
$V_i$ only decreases closed payments and remains eligible by
Lemma~\ref{arch:caplemma}. Thus Theorems~\ref{arch:tbarrier} and
\ref{arch:vowner} show that the leading owner constant $\sqrt2$ is sharp
within the respective architectures. Likewise, the general
$O(\sqrt d)$ dimension upper and Theorems~\ref{arch:tbarrier} and
\ref{arch:vdimension} establish the square-root \emph{order} within those
architectures; they do not identify the optimal dimension constant.

The same lower constructions apply to the explicitly specified canonical
relative support: on the coordinate face of $x\ne0$, take the
minimum-Euclidean-norm nonnegative supporting form, extend its value to
$+\infty$ outside that face, and at zero use zero on the zero bundle and
$+\infty$ on every nonzero bundle. At each nonempty regular mask, the active
resource block is interior and all reported bundles lie in that face;
the unique-gradient argument in Lemma~\ref{arch:realization} therefore
forces the same prices. At the empty mask every nonzero anchor bundle has
infinite raw price. Each global-cap anchor then pays $T$ to all owners,
and each pivotal marginal-cap anchor pays $V_i=T$ to all owners, strengthening
the required anchor estimate. This observation concerns that particular
relative-support convention, not every unspecified extended-valued selector.

These lower bounds exclude improvements obtained solely by changing a
report-independent mask distribution, correlating its bits, choosing different
common supports from the open reports, or replacing the global cap by
maximum coalition marginals. They do not exclude report-dependent masks
whose incentives are separately proved, owner-specific supports coupled by
a new coverage argument, or a different payment architecture. No efficient
evaluation theorem, human specialist endorsement, or historical priority
claim is implied by these constructions.

\section{Finite sampling and exact certification}\label{alg:section}
The expectation defining the hybrid mechanism is finite, but its direct
expansion can have exponentially many terms.  We give a polynomial sample
bound, followed by an exact certification procedure whose preprocessing cost
can be exponential.  These are different computational guarantees.
Throughout this section, fix the owner set $I=[n]$, the resource dimension $d$,
an ordered family $V=(w^1,\ldots,w^K)\subseteq\R_+^d$, and an integer $j\ge2$,
all independently of reports.  Set
\[
 f(c)=\min_{w\in V}w\cdot c,\qquad
 v_b(U)=f\Bigl(\sum_{i\in U}b_i\Bigr),\qquad T=T(b)=v_b(I),
\]
where $b\in(\R_+^d)^n$. Write $M_i=M_i^j$ for the exact hybrid payments
in \eqref{prod:p}--\eqref{prod:payment}, and recall their coalition coverage and budget bound
\begin{equation}\label{alg:exact-budget}
 \sum_{i\in U}M_i(b)\ge v_b(U),\qquad
 \sum_iM_i(b)\le C_j(K)T(b),\qquad
 C_j(K)=j+\frac{\ln K+2\ln j+1}{\ln(j^3/(j^2+j-1))}.
\end{equation}
The family and its priority remain fixed even at zero reports and ties.

\subsection{An exact finite sampling law}
For a mask $S\subseteq I$, put $S_i^-=S\setminus\{i\}$ and
$S_i^+=S\cup\{i\}$.  Let $W^-_{i,S}(b)$ be the first minimizer of the
\emph{raw} score $w\cdot\sum_{h\in S_i^-}b_h$, and define
\begin{equation}\label{alg:sample-components}
 A_i(S,b)=\min\{W^-_{i,S}(b)\cdot b_i,T(b)\},\qquad
 \Delta_i(S,b)=v_b(S_i^+)-v_b(S_i^-).
\end{equation}
Thus capping occurs after selection.  For nonnegative integers $a,k$, set
\begin{equation}\label{alg:integrals}
 J_{a,k}=\int_0^{1/j}q^a(1-q)^k\,dq
 =\sum_{h=0}^k(-1)^h\binom kh\frac{j^{-(a+h+1)}}{a+h+1},
 \quad \pi_k=jJ_{n-k,k},\quad
 \theta_k=\frac{J_{n-k+1,k}}{J_{n-k,k}}.
\end{equation}
Draw a cardinality $k$ with probability $\binom nk\pi_k$, then a uniformly
random $k$-subset $S$.  Define the vector associated with that mask by
\begin{equation}\label{alg:finite-sample}
 X_i(S,b)=\theta_{|S|}A_i(S,b)+(1-\theta_{|S|})\Delta_i(S,b).
\end{equation}
For $N$ independent masks, drawn once before any reports, let
\begin{equation}\label{alg:sample-payment}
 \widehat M_i(b)=\frac1N\sum_{\ell=1}^N X_i(S_\ell,b)
                       +\frac{\eta}{n}T(b),\qquad \eta>0.
\end{equation}

\begin{lemma}[Exact mean, bounded range, and every-seed incentives]
\label{alg:finite-law}
The probabilities in \eqref{alg:integrals} are rational and sum to one over
all masks.  Moreover, $0<\theta_k<1/j$, $0\le X_i(S,b)\le T(b)$, and
$\E_S X_i(S,b)=M_i(b)$.  For every fixed seed $(S_1,\ldots,S_N)$,
\eqref{alg:sample-payment} is nonnegative and componentwise nondecreasing
in each owner's entire own report, simultaneously at all profiles.
\end{lemma}
\begin{proof}
Consider first the experiment that draws $q$ uniformly from $(0,1/j)$ and,
conditionally on $q$, includes each owner independently with probability
$1-q$.  The probability of one particular mask of size $k$ is $\pi_k$,
and
\[
 \sum_{k=0}^n\binom nk\pi_k
 =j\int_0^{1/j}[q+(1-q)]^n\,dq=1.
\]
The posterior mean of $q$ conditional on that full mask is exactly
$\theta_k$.  In the continuous experiment the vector
$qA_i+(1-q)\Delta_i$ has conditional expectation
$q p_i^q+(1-q)m_i^q=r_i^q$: the owner's own sampled bit affects neither
term in \eqref{alg:sample-components}.  Taking expectations over $q$ gives
$M_i$.  Conditional expectation given the full mask gives
\eqref{alg:finite-sample}, hence its mean is also $M_i$.
This reasoning does not assume that $\theta_{|S|}$ is independent of $S$.

Both $A_i$ and $\Delta_i$ lie in $[0,T]$.  The former is capped; for the
latter, monotonicity gives $0\le v_b(S_i^- )\le v_b(S_i^+)\le T$.
The coefficients in \eqref{alg:finite-sample} are nonnegative and sum to one.
With the other reports and $S$ fixed, $W^-_{i,S}$ is independent of $b_i$.
Both its nonnegative linear score on $b_i$ and $T$ are nondecreasing, so
$A_i$ is own-report monotone.  The subtracted value in $\Delta_i$ omits $i$,
whereas its first value is nondecreasing in $b_i$.  The average and the
cushion in \eqref{alg:sample-payment} preserve this property.
At $T=0$ all values, samples, and payments vanish.
\end{proof}

\begin{proposition}[Rational sampling and evaluation]\label{alg:cost}
The mask law has an exact fair-bit sampler with expected cost polynomial in
$n$ and $\log j$.  A seed of length $N$ can be evaluated using at most
$1+N(n+1)$ canonical value-and-minimizer queries.  With a rational production
LP in $d$ resource-price coordinates and lexicographic priority on its dual
vertices, at most $1+N(n+1)(d+1)$ ordinary exact LP calls suffice.
These bit-model assertions concern rationally encoded technology and reports.
\end{proposition}
\begin{proof}
Let $H=(n+2)!j^{n+2}$.  Every denominator of a summand of
$jJ_{n-k,k}$ or $jJ_{n-k+1,k}$ divides $H$.  Therefore
\[
 u_k=H\pi_k>0,\qquad t_k=HjJ_{n-k+1,k},\qquad
 U_k=\binom nk u_k
\]
are integers, $\theta_k=t_k/u_k$, and $\sum_kU_k=H$.
Their bit lengths are $O(n(\log j+\log(n+2)))$; computing the alternating
sums exactly takes $O(n^2)$ operations on polynomial-length integers.
Draw a uniform integer below $H$ to choose a cardinality by the intervals
of lengths $U_k$, then a uniform integer below $\binom nk$ and unrank its
subset.  Rejection from a power-of-two range takes fewer than two attempts
in expectation.  It terminates almost surely, with no bounded worst-case
number of attempts.  Floating-point subtraction is unnecessary.

For one mask, the coalitions $S$ and $S\mathbin{\triangle}\{i\}$,
$i\in I$, supply all off/on values and all off minimizers; there are $n+1$
such coalitions.  Query $T$ once.  An explicitly listed $V$ permits a canonical
query by scanning $K$ scores in $O(Kd)$ arithmetic operations.
For the rational dual polyhedron $Q=\{w\ge0:A^\top w\ge\rho\}$,
first minimize the queried resource objective.  On its optimal face minimize
$w_1$, then $w_2$, and so on, retaining previous optimum values as equalities.
Each optimization is attained and bounded below by zero.  The final point
is unique and is a vertex: any nontrivial feasible convex decomposition
would preserve each successive minimum and finally every coordinate.
It is the lexicographically first optimal vertex, a fixed priority independent
of reports.  This takes at most $d+1$ LP calls per canonical query; the grand
value needs one value-only solve.  Exact rational LP optimization preserves
the required ties.  An arbitrary priority on an implicit vertex family need
not admit this compact oracle.
\end{proof}

\begin{proposition}[One fixed report profile]\label{alg:pointwise}
For $0<\delta<1$, if
\[
 N\ge \frac{n^2}{2\eta^2}\ln\frac{2n}{\delta},
\]
then at any fixed profile, with probability at least $1-\delta$, the rule
\eqref{alg:sample-payment}, with the supplied coefficients $\theta$ and $\eta$,
covers every coalition and has total at most
$(C_j(K)+2\eta)T$.  Its expected total is at most $(C_j(K)+\eta)T$.
\end{proposition}
\begin{proof}
For $T>0$, apply Hoeffding's inequality to $X_i/T\in[0,1]$ and union bound
over the $n$ owners.  With the stated probability each coordinate error is
at most $\eta T/n$.  Adding the cushion makes $\widehat M_i\ge M_i$ and
$\widehat M_i\le M_i+2\eta T/n$ for every $i$; now use
\eqref{alg:exact-budget}.  The expected-total assertion follows from
$\E\widehat M_i=M_i+\eta T/n$.  The case $T=0$ is exact.
\end{proof}
The last proposition holds for a profile fixed independently of the seed.
The next theorem permits the seed to be public and profiles to be selected
adaptively, because its event quantifies over every profile at once.

\subsection{One seed for every real report profile}
\begin{theorem}[Uniform finite-sample implementation]\label{alg:uniform}
Let $0<\eta\le1$, $0<\delta<1$, and put
\[
 p=nd,\quad D=\lceil16p\log_2(8pK)\rceil,\quad
 \varepsilon=\eta/n,\quad A_0=32e/\varepsilon.
\]
If the integer sample size satisfies
\begin{equation}\label{alg:sample-size}
 N\ge\frac{128}{\varepsilon^2}
       \left[D\ln(A_0\ln A_0)+\ln\frac{8n}{\delta}\right],
\end{equation}
then with probability at least $1-\delta$ over the single initial seed,
\emph{simultaneously for all real $b\ge0$ and all $U\subseteq I$},
\[
 \widehat M_i(b)\ge M_i(b)\quad(i\in I),\qquad
 \sum_{i\in U}\widehat M_i(b)\ge v_b(U),\qquad
 \sum_i\widehat M_i(b)\le(C_j(K)+2\eta)T(b).
\]
Incentives hold for every seed, whether or not this event occurs.
The sample size is polynomial in $n,d,\log K,1/\eta,\log(1/\delta)$.
\end{theorem}
\begin{proof}
For owner $i$, consider the $[0,1]$-valued class
\[
 h_{b,i}(S)=\begin{cases}X_i(S,b)/T(b),&T(b)>0,\\0,&T(b)=0.
 \end{cases}
\]
We bound its pseudodimension, meaning the largest number of observations
that can be shattered using a separate real threshold at each observation.
Fix $m$ masks $S_\ell$ and thresholds $a_\ell$.
As functions of the $p$ report coordinates, use the linear forms
\[
 C_v=w^v\cdot\sum_h b_h,\quad
 L^-_{\ell,v}=w^v\cdot\sum_{h\in S_\ell\setminus\{i\}}b_h,
 \quad L^+_{\ell,v}=L^-_{\ell,v}+w^v\cdot b_i.
\]
Partition report space by all pairwise comparisons of the $C_v$, all pairwise
comparisons within each off/on family, and all cap comparisons
$w^v\cdot b_i-C_t$.  The number of hyperplanes is at most
\[
 H_0=(2m+1)\binom K2+K^2\le3mK^2.
\]
Count all ternary sign faces, including lower-dimensional tie faces.
For $H$ hyperplanes in $\R^p$ their number is at most
\begin{equation}\label{alg:faces}
 F(H,p)=\sum_{r=0}^p2^r\binom Hr\le(1+2H)^p.
\end{equation}
Indeed, insertion of the last hyperplane gives the recurrence
$F(H,p)\le F(H-1,p)+2F(H-1,p-1)$, with the usual zero-dimensional
and zero-hyperplane initial conditions; degeneracy only reduces the count.
The displayed sum solves this recurrence.  The second inequality follows
by comparison with the binomial expansion of $(1+2H)^p$.

On each first-arrangement face, all priority winners and cap branches are
fixed, including ties.  Thus $X_i(S_\ell,b)$ is a fixed linear form
$U_\ell(b)$ and $T$ is a fixed form $C_t(b)$ there.
On $T>0$, the subgraph tests become
$U_\ell(b)-a_\ell C_t(b)>0$.  These $m$ additional hyperplanes give at
most $F(m,p)$ label patterns on the face.  All $T=0$ profiles supply one
additional label pattern.  Hence, for $m\ge1$, the subgraph growth function
satisfies
\[
 \Pi_i(m)\le1+F(H_0,p)F(m,p)\le2(21m^2K^2)^p.
\]
If $m$ observations were shattered, taking base-two logarithms gives
$m\le1+p\log_2 21+2p\log_2K+2p\log_2m$.
With $r=m/p$ and $L=\log_2(8pK)\ge3$, this implies
$r<2L+2\log_2r$.  For $r\ge16L$ we have $r\ge48$,
$2L\le r/8$, and $2\log_2r\le r/2$, a contradiction.
Consequently $D$ is an upper bound on the pseudodimension.

The bounded-class uniform-deviation inequality and pseudodimension covering
bound of Haussler~\cite{Haussler1992} give
\begin{equation}\label{alg:uniform-tail}
 \Prob\left\{\sup_b\left|\frac1N\sum_\ell h_{b,i}(S_\ell)
                  -\E h_{b,i}\right|>\varepsilon\right\}
 \le8[A_0\ln A_0]^D e^{-N\varepsilon^2/128}.
\end{equation}
Specifically, Corollary~1 of the January 1992 manuscript gives the factor
$4$, covering radius $\varepsilon/16$, and exponent $128$; its Theorems~6
and~12 bound that covering number by $2[A_0\ln A_0]^D$.
There is no measurability difficulty here: masks, and hence length-$N$
experiments, have finite sample spaces.  The finite-branch parametrization
is also Borel, including the explicitly defined $T=0$ branch.
A union bound over owners and \eqref{alg:sample-size} make the failure
probability at most $\delta$.  On this single event every $i,b$ satisfies
\[
 \left|N^{-1}\sum_\ell X_i(S_\ell,b)-M_i(b)\right|
       \le\eta T(b)/n.
\]
Use the cushion and \eqref{alg:exact-budget} as in
Proposition~\ref{alg:pointwise}.  No union bound over coalitions is required.
\end{proof}
The rule can pay $\eta T/n$ to an owner reporting zero.  If zero payment to
that owner is desired, replace this cushion by
$(\eta T/n)\ind\{b_i\ne0\}$.  Since $X_i=M_i=0$ at $b_i=0$, the same
concentration argument works, and the modified cushion is own-report monotone.
We use the unmodified rule in the certification statement below.

\subsection{Exact certification, including every tie face}
Uniform good-seed existence does not by itself certify a sampled seed.
For this purpose assume that the complete ordered list $V$ is rational.
Fix any masks $S_1,\ldots,S_N$, rational $\theta_0,\ldots,\theta_n\in[0,1]$,
$\eta\ge0$, and rational budget threshold $\overline C$.
The checker below does not require that the seed was randomly generated.

\begin{theorem}[Sound and complete all-profile seed certification]
\label{alg:certification}
There is a terminating exact algorithm deciding whether
\eqref{alg:sample-payment}, with the supplied coefficients $\theta$ and $\eta$,
covers every coalition and pays at most
$\overline C T$, at every real nonnegative report profile.
On rejection it returns a rational violating profile and constraint.
With $p=nd$ and
\[
 H\le 2^n\binom K2+nK^2,
\]
a conservative running-time bound is
\begin{equation}\label{alg:cert-cost}
 (2^n+1)(1+H)^{O(nd)}\operatorname{poly}(N,n,d,K,L),
\end{equation}
where $L$ bounds rational input bit lengths.  In particular this is not a
polynomial preprocessing bound.  Incentives already hold for every seed.
\end{theorem}
\begin{proof}
All relevant functions are positively homogeneous in the entire report
vector.  The zero vector is harmless.  It therefore suffices to check
\[
 \mathcal D=\{b\in\R_+^p:\textstyle\sum_{i,e}b_{ie}=1\}.
\]
Normalize by total report mass, not by $T$: nonzero profiles with $T=0$
remain in this compact simplex.  For every $U\subseteq I$ define
$L_{U,v}(b)=w^v\cdot\sum_{i\in U}b_i$.  Form the arrangement of
\[
 L_{U,v}-L_{U,t}=0\quad(U\subseteq I,\ v<t),\qquad
 w^v\cdot b_i-L_{I,t}=0\quad(i\in I,\ v,t\in[K]).
\]
Omit identically zero forms.  These at most $H$ hyperplanes fix every
coalition's minimizing family and every possible cap branch.

Write the retained forms as $\ell_1,\ldots,\ell_H$.
For a ternary sign pattern $\sigma$, let $C_\sigma\subseteq\mathcal D$
be the region with $\ell_h=0$ when $\sigma_h=0$ and
$\sigma_h\ell_h>0$ otherwise.  Let $P_\sigma$ replace all these strict
inequalities by weak ones.  Determine nonemptiness of $C_\sigma$ by the
rational LP
\[
 \max\tau:\quad b\in\mathcal D,\quad0\le\tau\le1,\quad
 \ell_h(b)=0\ (\sigma_h=0),\quad
 \sigma_h\ell_h(b)\ge\tau\ (\sigma_h\ne0).
\]
Retain exactly those patterns with a positive optimum, storing a rational
strict witness $z_\sigma$.  If there are no strict constraints, the optimum
is one whenever the equality-constrained simplex is nonempty.
The criterion is exact because finitely many positive margins have a positive
minimum.  The compact rational LP supplies rational witnesses.

For a retained pattern,
\begin{equation}\label{alg:closure}
 \overline{C_\sigma}=P_\sigma.
\end{equation}
Indeed, for any $x\in P_\sigma$ and $0<t\le1$, the point
$(1-t)x+t z_\sigma$ satisfies every strict sign and tends to $x$.
This also covers proper coordinate faces and degeneracies.
At $z_\sigma$, record the priority winners and cap branches.  They give
fixed rational linear formulas $V_{U,\sigma}$ for coalition values and
$P_{i,\sigma}$ for payments throughout $C_\sigma$.
For every $U$, and also for the budget, minimize over $P_\sigma$ the slacks
\[
 G_{U,\sigma}=\sum_{i\in U}P_{i,\sigma}-V_{U,\sigma},\qquad
 G_{0,\sigma}=\overline C V_{I,\sigma}-\sum_iP_{i,\sigma}.
\]
Accept exactly when all these LP minima are nonnegative.
Crucially, the objective on the closure is the \emph{frozen formula from
$C_\sigma$}.  Do not recompute a boundary point's tie-resolved payment:
that could belong to a different branch.  Every tie region is checked
separately under its own actual formula.

A linear form is nonnegative on $C_\sigma$ if and only if it is nonnegative
on its closure; \eqref{alg:closure} proves exactness of these tests.
If a rational closure optimizer $x$ has $G(x)<0$, choose
\[
 0<t=\min\left\{\frac12,
       \frac{-G(x)}{2(1+|G(x)|+|G(z_\sigma)|)}\right\}.
\]
Then $(1-t)x+t z_\sigma$ is a rational point of the original strict region
and still has negative slack.  It is an actual violating profile, not merely
a violation of an inactive boundary formula.  Regions partition the simplex,
so the algorithm is sound and complete.  At $T=0$ the actual rule vanishes;
no division by $T$ is used by the checker.

To obtain \eqref{alg:cert-cost}, enumerate signs incrementally.  Insert each
hyperplane, try the three extensions of every retained prefix, and use the
margin LP to discard unrealizable extensions.  Formula~\eqref{alg:faces}
bounds the number of prefixes at each stage by $(1+2H)^p$.
At most $3H(1+2H)^p$ margin LPs and
$(2^n+1)(1+2H)^p$ slack LPs suffice.  Each has at most $p+1$ variables,
$O(p+H)$ constraints, and polynomial input bit length.  Exact rational LP
algorithms yield the bound.  One may retain the exhaustive branching tree,
margin witnesses or rejection certificates, frozen formulas, and rational
LP dual certificates.  Their verification is polynomial in this explicit
certificate's length, which can itself be exponential.
\end{proof}
For the optional zero-reporter cushion, add the $n$ arrangement forms
$\sum_e b_{ie}=0$; these determine its indicator on the nonnegative domain.
If a compact production LP is given instead of $V$, complete vertex
enumeration is additional preprocessing: solve all independent active
$d$-tuples of its dual constraints, retain feasible solutions, and deduplicate.
No polynomial vertex-enumeration claim is needed.

\begin{corollary}[Always-correct finite implementation after preprocessing]
\label{alg:las-vegas}
Fix rational $0<\eta\le1$, $0<\delta<1$, $j$, and a sample size satisfying
\eqref{alg:sample-size}.  Supply a certified rational threshold
$\overline C\ge C_j(K)+2\eta$.
Repeated independent sampling followed by Theorem~\ref{alg:certification}
returns a certified seed almost surely, after at most $1/(1-\delta)$ trials
in expectation.  Every output rule is own-report monotone and satisfies its
coverage and budget guarantees at all real profiles.  Its subsequent
rational-input evaluation obeys Proposition~\ref{alg:cost}.
\end{corollary}
\begin{proof}
Every seed in Theorem~\ref{alg:uniform}'s event passes the exact checker.
Thus each independent trial succeeds with probability at least $1-\delta$.
The geometric tail proves almost-sure termination and the expected-trial
bound.  Every accepted seed has an exact certificate, regardless of the
randomness used to find it.  Sampling and certification precede all reports,
so freezing the accepted seed preserves its every-seed incentives.
\end{proof}
A rational threshold can be certified using rigorous logarithm enclosures,
with arbitrarily small extra budget slack.  A certified upper bound on $K$
may replace $K$ throughout, but this tunes the result to that bound rather
than to an unknown smaller actual count.  For a production LP with $m$
activities, the elementary bound $K\le\binom{m+d}{d}$ is available.
Choosing $j$ and bounding $K$ are technology-side tasks, never report-dependent
adjustments.  Parameters $n,d$ count the fixed domain; zero reporters are not
removed after reports arrive.
The construction has polynomial sample count and polynomial evaluation in
the displayed sample and rational-input parameters, using the stated oracle.
Its generally exponential certification preprocessing is separate.  There is
no practical sample-size, polynomial preprocessing, or efficient deterministic
seed-construction claim.  Nor does this finite-form implementation provide an
oracle for arbitrary concave production functions.

\section{Remaining mathematical and review questions}\label{sec:remaining}
Theorems~\ref{prod:upper} and~\ref{dag:sharp} determine the sharp leading
constant in supporting-form or structural-cut complexity. They do not determine
exact finite-$K$ optima or lower-order terms. The owner-count range remains
between $(1-o(1))\log n/\log\log n$ and $(\sqrt2+o(1))\sqrt n$ for general
concave homogeneous production; the dimension range remains between
$(1-o(1))\log d/\log\log d$ and $\min\{d,1+2\sqrt d\}$.
Section~\ref{arch:section} blocks improvements to square-root order within
the specified common-support mask architectures, even with maximum-marginal
caps. It does not give unrestricted square-root lower bounds. A new rule
escaping those restrictions or a stronger unrestricted lower is a substantive
next target. The exact $d$ lower for the larger
nonconcave class does not settle the concave problem.

The finite implementation preserves all-profile incentives and coverage,
but exact seed certification can be exponential. Polynomial preprocessing,
a direct deterministic good-seed construction, or a precise computational
obstruction would strengthen the result. The actual monotone envelope has a
polynomial lower obstruction and a $KT$ upper; its precise growth remains
open, although that gap does not affect the sharp theorem for other mechanisms.

Historical priority and significance require separate specialist scrutiny.
The current source comparison is bounded and partly access-limited. Internal
model reviews are recorded with their scopes and adjudications; they do not
supply human specialist approval. The manuscript and accompanying short brief
are intended to make those questions concrete and reviewable.


\appendix

\section{Two simultaneously altered binary pools}\label{app:binary}
Use \eqref{eq:abc} and put $F_h=h b_h+(r-h)a_h+rz_h$. For an altered pool
$j$ with $h$ unit owners, let $R^j$ pay $a_h,b_h$ to its own low and unit
owners and $e_h$ to \emph{every} outside owner, including outside low owners.
This vector covers the actual two-altered-pool game: set every outside report
to one, apply Theorem~\ref{thm:escape}, and observe that actual coalition
values only decrease.

For $0\le s\le r-2$, let $q_s=a_s/b_{s+1}$. Equation~\eqref{eq:ratio}
gives $0\le q_s\le1/(r-1)$ and $b_{s+1}\ge(r-1)a_s$.
For $0\le g<h<r$ define
\[
 w(h,g)=\frac12\prod_{s=g}^{h-1}q_s.
\]
A zero factor gives weight zero; no reciprocal of $q_s$ is used.
For two altered pools with equal counts, use their equal mixture. With counts
$h>g$, give the $h$-pool component weight $w(h,g)$ and the $g$-pool component
the remaining weight. With zero or one altered pool use the previous rule.

\begin{theorem}\label{thm:binary-two}
This assignment covers all coalitions, is own-monotone on every retained
binary edge, preserves Theorem~\ref{thm:escape} on its entire boundary,
and has base factor at most $2$, or factor $3$ with the floor.
\end{theorem}
\begin{proof}
Coverage follows from the component argument and convexity. Check the changing
owner's payment on each possible increase.
\begin{itemize}
\item Equal counts $s,s$ become $s+1,s$: the old payment is
$(a_s+e_s)/2$, and the new one is
$(q_s/2)b_{s+1}+(1-q_s/2)e_s=a_s/2+(1-q_s/2)e_s$, which is larger.
\item An already larger count $h>g$ increases internally: writing
$w=w(h,g)$ and $w'=wq_h\le w$, old and new payments are
$wa_h+(1-w)e_g$ and $w'b_{h+1}+(1-w')e_g$
$=wa_h+(1-w')e_g$, respectively.
\item A smaller count $g<h$ increases: put $w'=w(h,g+1)$, or $1/2$ at a
tie, so $w=w'q_g\le w'\le1/2$. The old payment is
$(1-w)a_g+we_h$ and the new one is $(1-w')b_{g+1}+w'e_h$.
For $r\ge3$, $b_{g+1}\ge2a_g$ makes the new own-component term at least
$a_g$, while the outside term increases. For $r=2$, the only case is
$g=0,h=1$, where $q_0=1$ and $b_1=a_0=1/4$, giving equality.
\item The last low owner reaches one: $a_{r-1}=0$, so its old payment is
at most the other component's outside payment, which it receives exactly
after its pool becomes unaltered. With no other altered pool its payment
instead increases from zero to $1/k$.
\end{itemize}
One-altered internal edges were proved in Theorem~\ref{thm:escape}.

For the budget we use, for integers $0\le g<h<r$,
\begin{equation}\label{eq:slack}
 2T_g-F_g\ge1/r,\quad F_h-F_g\le2(h-g),\quad
 a_g>0\Longrightarrow F_g=T_g\bigl(2-(g+1)/r\bigr).
\end{equation}
Here is a direct verification. Extend $h$ to real $x\in[0,r)$ in the formulas.
Where $\delta_x\ge z_x$, $F(x)=T_x(2-(x+1)/r)$ and $F'(x)\le2$;
its slack is $T_x(x+1)/r\ge1/r$. Where $\delta_x\le z_x$,
$F(x)=x+(r-x)T_x/r$ and
$F'(x)=2(r-1)(1-x/r)/r\le2$. The formulas agree at the transition.
At integer $g$ in the latter region,
$2T_g-F_g=(r-g)/r+gT_g/r\ge1/r$. This proves \eqref{eq:slack}.

Equal counts cost $F_g\le2T_g$. Otherwise actual $T=T_g$, and total cost is
$F_g+w(h,g)(F_h-F_g)$. A nonpositive difference is harmless. Set $d=h-g$.
The excess is at most $d/(r-1)^d$ by \eqref{eq:slack} and the product bound.
If $r\ge4,d\ge2$, this is at most $2/(r-1)^2\le1/r$, within the slack.
If $d=1,g\ge1$, either $a_g=0$ and the weight is zero, or the slack is
$T_g(g+1)/r\ge2/r$ and the excess is at most $1/(r-1)\le2/r$.
If $d=1,g=0,r\ge3$, then
\[
 w=\frac1{2(r-1)},\quad F_0=2-1/r,\quad
 F_1-F_0=\frac{(2r-1)(r-2)}{r^2},
\]
and the excess is at most $1/r$ since
$(2r-1)(r-2)\le2r(r-1)$.
The two remaining cases are $r=2,h=1,g=0$, with cost $13/8\le2$, and
$r=3,h=2,g=0$, with cost
$5/3+(1/16)(10/9)=125/72\le2$. These exhaust the parameters.
Apply Lemma~\ref{lem:floor}.
\end{proof}
The one-altered $h=0$ boundary costs $(2-1/r)T$, tending to $2T$.
Thus factor two is sharp \emph{for rules required to preserve this boundary}.
No minimax lower bound over freely chosen nonlinear boundaries follows.

\section{A sharper pairwise probability bound}\label{app:smallball}
Retain the capped Banzhaf rule of Section~\ref{sec:banzhaf}. Normalize its
capacities by $T$ and its pool totals by $\rho_j=D_j/T$, so a shortest pool
$s$ has $\rho_s=1$. For $\rho\ge2$, write $n=\lfloor\rho\rfloor$ and
$\theta=\rho-n$, and define
\[
 q(n)=\frac1{(n+1)!},\qquad
 q(n+\theta)=\frac{1-(1-\theta)^{n+2}}{\theta(n+2)!}
 \quad(0<\theta<1).
\]
\begin{theorem}\label{thm:smallball}
For every pool $j$ with $\rho_j\ge2$,
$\Prob(X_j<X_s)\le q(\rho_j)$, and this pairwise bound is attainable.
Consequently the base rule has total factor at most
\[
 Q(k)=2+\int_2^\infty\min\{1,(k-1)q(t)\}\,dt<U(k),\qquad k\ge2.
\]
No optimality or fixed additive improvement is asserted for $Q(k)$.
\end{theorem}
\begin{proof}
Here normalize $X_j$ by $T$. A sum of independent uniform summands has
a symmetric unimodal density. To check the needed property, center each
summand: if $f$ is even and nonincreasing on the positive half-line, its
convolution with the uniform density on $[-a/2,a/2]$ is even, with derivative
$[f(x+a/2)-f(x-a/2)]/a\le0$ for almost every $x\ge0$.
Induction starts with a single uniform density; zero coefficients can be
dropped. Thus the distribution function $F_j$ is convex on $[0,1]$ when
$\rho_j\ge2$, since this interval lies to the left of its center
$\rho_j/2$. The shortest pool has coefficients $b_i\ge0$ summing to one.
Pointwise Jensen and independence give
\[
 \Prob(X_j<X_s)=\E F_j\Bigl(\sum_i b_iU_i\Bigr)
 \le\sum_i b_i\E F_j(U_i)
 =\int_0^1F_j(u)\,du=\E(1-X_j)_+.
\]
The last expression is a convex symmetric function of the coefficients of
$X_j$. On the compact polytope $\{a\in[0,1]^r:\sum_i a_i=\rho_j\}$,
its maximum occurs at a vertex. Every vertex has $n$ unit coordinates,
one coordinate $\theta$ if $\theta>0$, and zeros: two fractional coordinates
would admit a nontrivial opposite perturbation. This distributional bound
does not split strategic owners.

For $S_n=\sum_{i=1}^nU_i$ and $0\le t\le1$, simplex integration gives
$\E(t-S_n)_+=t^{n+1}/(n+1)!$. At $\theta=0$ this gives $q(n)$;
conditioning on the additional $\theta U$ gives the stated fractional formula.
A singleton shortest pool and these extremal coefficients attain the
pairwise bound. Coupling shows $q$ is nonincreasing, and its formulas agree
at the integer endpoints.

At most $k-1$ pools have $\rho_j>t$ for $t\ge2$, so a union bound gives
$\Prob(\rho_W>t)\le\min\{1,(k-1)q(t)\}$. Integrating, and bounding
the lower tails by one, proves $Q(k)$. For $t\ge2$ and $n=\lfloor t\rfloor$,
\[
 q(t)\le\frac1{(n+1)!}
 <\left(\frac e{n+1}\right)^{n+1}
 \le\left(\frac et\right)^t.
\]
The strict factorial inequality follows from the exponential series; the
last inequality uses the decrease of $(e/x)^x$ for $x>1$.
On $[2,e]$ bound the integrand by one, then use the tail argument of
Theorem~\ref{thm:upper} to obtain
\[
 Q(k)\le e+\int_e^\infty\min\{1,k(e/t)^t\}\,dt
 \le R_k+1/\log R_k.
\]
The integrand comparison is strict for $t>R_k$, so the final inequality is
strict. The optional floor again costs at most one by Lemma~\ref{lem:floor}.
\end{proof}

\section{Reproducible evidence}\label{app:verification}
Table~\ref{tab:historical-checks} records the October~7 verification package,
retained in the historical archive. Programs use standard Python, exact
rationals or integers, and explicit checks under \texttt{-O}. These are
historical finite counts; the current Online Resource 2 has a narrower scope.
\begin{table}[htbp]
\centering
\caption{Historical finite checks from the October 7 package.}
\label{tab:historical-checks}
\small
\begin{tabular}{p{0.29\linewidth}p{0.62\linewidth}}
\toprule
Program & Principal covered claims\\
\midrule
\texttt{poisson\_certificate.py} & 24 exact weighted-profile duals over all
histograms in each instance; 168 specified corruptions rejected.\\[3pt]
\texttt{hierarchy\_certificate.py} and \texttt{verify.py} & Fourteen exact
hierarchy duals, including padded instances; 98 specified corruptions rejected.\\[3pt]
\texttt{verify\_small\_ball.py} & 24 exact rational pairwise probability
comparisons, including fractional totals and attained equality.\\[3pt]
\texttt{verify\_thresholds.py} & Capped Banzhaf probabilities by exact polynomial
integration; 4,227 full profiles, 221,253 coalition checks, 16,040 own-report
comparisons; four additional tail stress cases.\\[3pt]
\texttt{verify\_lower.py} & 53 exact unrestricted lower certificates and 16
specified rejected corruptions.\\[3pt]
\texttt{verify\_beyond.py} & Six proportional certificates, eight corruptions,
one-pool nonlinear coalition and edge checks; auxiliary archive checks.\\[3pt]
\texttt{verify\_two\_pools.py} & 1,330 complete labeled binary profiles,
3,129,152 coalition/profile pairs and 5,008 directed report edges.\\[3pt]
\texttt{verify\_continuous.py} & 24,135 labeled profiles on fifteen grids,
3,463,360 coalition/profile pairs and 235,523 directed report edges, including
zero and below-$1/r$ reports.\\[3pt]
\texttt{verify\_scope.py} & Exact substitutions for the continuous rule's
many-altered-pool limitation and an alternative mixture's failure.\\
\bottomrule
\end{tabular}
\end{table}
The command \texttt{python3 -B -O verification/run\_all.py} executes the six
earlier scripts in isolated output directories and verifies report/source hashes.
The command \nolinkurl{python3 -B -O verification/run_hierarchy.py} reproduces
the hierarchy and pairwise-probability checks in fresh output directories.
The command \texttt{python3 -B -O verification/run\_sharp.py} checks the
weighted-profile duals in fresh output directories.
The verification archive records its file manifest and clean replay result.
These checks supplement rather than replace each proof. They do not certify
novelty, rounded implementations, or the universal theorem without its proof.


\section{Verification of the consolidated extensions}\label{app:consolidated-checks}
The October~7 portable packet preserved repository-relative paths and
provided commands for the general production, concave, network, sampling,
certificate and historical routines. Its runner executed 26 standard-library
checkers with \texttt{-B -O}, recording dependency hashes and full outputs;
the earlier hierarchy and finite-pool checks were separate. Those dated
records bind the historical inputs. The current Online Resource 2 instead
replays five scoped checks for the central theorem and projection comparison,
including the repaired two-owner example and small-network symmetries. It
does not claim a fresh replay of every historical supplementary check.
These finite checks corroborate identities and explicit certificates.
Universal and asymptotic statements rely on their written proofs; all-profile
seed certification is a mathematical algorithm theorem, not a claim that
large generic instances were solved. Historical source comparison and
independent human validation remain separate obligations.

% Consolidated earlier incentive results; insert before reproducible evidence.
\section{Finite-pool mechanisms and limits of restricted rules}\label{hist:appendix}
This appendix collects nonredundant results from the earlier incentive
campaigns. All guarantees use the fixed-owner, binding-underreport,
payment-only model of Section~\ref{sec:model}. The subsidy factor is
$\beta=B-1$, not $B$. The mechanisms and their restricted-class comparisons
remain useful even where later theorems improve the asymptotic frontier.
The weighted-glove model, ordinary Shapley value, and continuous Banzhaf
functional are established ingredients; the comparisons below make no
historical-priority claim.

\subsection{Aggregate prices and an exact restricted-class tradeoff}
\begin{proposition}\label{hist:aggregate}
For $T>0$, write $C_j=c(N_j)$ and define
\[
 p_j=\frac1{C_j}\int_0^T\prod_{\ell\ne j}(1-t/C_\ell)\,dt,
 \qquad y_i=c_ip_j\quad(i\in N_j).
\]
The rule covers coalitions, is own-report monotone, protects $H_i/k$, and
has total payment at most $(k+1)T/2$. This budget bound is a supremum for
this rule. For two pools and $\alpha\in[0,1]$, the optimal worst-case
subsidy factor among rules $y_i=c_ip(C_1,C_2)$ in pool one and
$y_i=c_iq(C_1,C_2)$ in pool two is
$\max\{1/2,\alpha\}$, provided the same nonnegative prices are used for
\emph{every} finite ownership partition and the rule protects $\alpha H$.
\end{proposition}
\begin{proof}
Differentiating $P(t)=\prod_j(1-t/C_j)$ and integrating gives
$\sum_jp_j=P(0)-P(T)=1$. Thus payments to a coalition are a convex
combination of its pool totals. For an owner reporting $x$ with fixed
same-pool remainder $a$, set
$g(t)=\prod_{\ell\ne j}(1-t/C_\ell)_+$. Its payment is
\[
 \frac{x}{x+a}\int_0^{x+a}g(t)\,dt.
\]
Both factors are nonnegative and nondecreasing in $x$. Assign zero at
$x=a=0$; if another pool has zero total all payments stay zero. This proves
monotonicity at every boundary as well.
The complementary-coalition constraint gives
$H_i\le T-v_c(N\setminus\{i\})\le c_iT/C_j$.
For the last inequality, it is equality as an upper bound when $C_j=T$;
otherwise the marginal is $(c_i-C_j+T)_+\le c_iT/C_j$.
Also $p_j\ge T/(kC_j)$, since every factor in the integral is at least
$1-t/T$. Choose a pool of total $T$. Its payment is at most $T$;
every other pool receives at most $\int_0^T(1-t/T)dt=T/2$.
Totals $(1,M,\ldots,M)$ approach this bound as $M\to\infty$.

For the two-pool upper bound, the displayed rule gives $\beta=1/2$ and
$H/2$. The auxiliary rule $z_i=c_iT/C_j$ gives $H$, covers coalitions,
and spends $2T$. Its own payment is $x\min\{1,D/(x+a)\}$, where $D$ is
the other pool total, so it is nondecreasing. For $\alpha\ge1/2$, the
convex combination $(2-2\alpha)y+(2\alpha-1)z$ gives $\alpha H$ and
$\beta\le\alpha$; for smaller $\alpha$, use $y$ itself.

For the lower bound, one owner per pool at totals $(1,1)$ gives
$p,q\ge\alpha$ and $p+q\ge1$. One price is at least
$r=\max\{1/2,\alpha\}$; suppose it is $p$.
Raising that owner's report to $M$ gives $Mp(M,1)\ge r$.
Now partition the left total into owners of capacities $1,M-1$.
Partition-independent prices are unchanged. The unit left owner together
with the right owner imply $p(M,1)+q(M,1)\ge1$. Hence total payment is
at least $1+(M-1)p(M,1)\ge1+r(1-1/M)$. Let $M\to\infty$.
This argument uses the stated comparison \emph{across} ownership
partitions; it is not an unrestricted fixed-owner lower bound.
\end{proof}

For clarity, these pool totals are themselves the established absolute
continuous Banzhaf functional \cite[Definition~6.11]{Kurz}. For
$F(z)=\min_j C_jz_j/T$, the zero endpoint is zero and the other endpoint
has tail integral
\[
 T\operatorname{BZI}_j(F)
 =\E\min\{C_j,\min_{\ell\ne j}C_\ell U_\ell\}
 =\int_0^T\prod_{\ell\ne j}(1-t/C_\ell)dt=C_jp_j.
\]
Here the coordinates are \emph{pools}. The owner-coordinate identification
in Section~\ref{sec:banzhaf} concerns a different continuous game.

\subsection{What ordinary Shapley scaling can and cannot achieve}
Let $\phi$ be the ordinary individual Shapley value, with fixed weights
$|U|!(n-|U|-1)!/n!$ on predecessor sets $U\subseteq N\setminus\{i\}$.
Nonnegativity and efficiency follow by monotonicity and telescoping along
each arrival order. Each own marginal is nondecreasing in its owner's
capacity, so $\phi_i$, and every nonnegative constant multiple of it, is
own-report monotone; this truthfulness argument is established in
\cite[Proposition~1]{Huang}.

\begin{proposition}\label{hist:shapley-half}
In every finite two-pool game,
$\phi_i\ge M_i/2\ge H_i/2$ and $\phi(S)\ge v_c(S)/2$.
Both constants $1/2$ are sharp over finite ownerships. Consequently a
constant multiple $\lambda\phi$ covers every coalition at every profile
and ownership if and only if $\lambda\ge2$. Even a profile-dependent
scalar multiplier cannot have a uniform budget factor less than two.
\end{proposition}
\begin{proof}
Fix an owner of capacity $a$, same-pool remainder $A$, and opposite total
$D+A$. Put $f_a(z)=\min\{a,z_+\}$. A predecessor subset with opposite-minus-
same capacity $z$ gives marginal $f_a(z)$, while its complement gives
$f_a(D-z)$. For all real $z,D$,
\[
 f_a(z)+f_a(D-z)\ge f_a(D).
\]
If $D\le0$ this is immediate; if either argument is at least $a$ it is
also immediate. Otherwise the positive parts are untruncated and their
sum is at least $D$. Complementary predecessor sets have equal Shapley
weights, so $2\phi_i\ge f_a(D)=M_i$. The profile $(1\mid1)$ attains
$\phi_i=H_i/2$. If $L$ is a globally minimum-total pool, every $i\in L$
has $M_i=c_i$; hence
$\phi(S)\ge\phi(S\cap L)\ge c(S\cap L)/2\ge v_c(S)/2$.

For sharpness of the coalition constant, choose an integer $q\ge2$,
$m=q^4$, $r=q^4-q^3$, and a game with $m$ left unit owners, $r$ right
unit owners, and one right owner $b$ of capacity $m$. Let $S$ contain all
left owners and $b$, so $v_c(S)=T=m$. Independent uniform arrival times
represent a uniform random order. For a left unit arriving at time $t$,
its marginal is one if $b$ has arrived; otherwise it is one precisely
when $Y>X$, with independent $X\sim\operatorname{Bin}(m-1,t)$ and
$Y\sim\operatorname{Bin}(r,t)$. Thus
\[
 \phi_{\rm left}=\tfrac12+\eta,
 \qquad \eta=\int_0^1(1-t)\Prob(Y>X)dt.
\]
Writing $D_0=m-1-r=q^3-1$, the difference has mean $-D_0t$ and variance
at most $m/2$. Chebyshev gives
$\Prob(Y>X)\le\min\{1,m/(2D_0^2t^2)\}$, whose integral is at most
$\sqrt{2m}/D_0$. At $b$'s arrival, its marginal is $(X-Y)_+$ for counts
with parameters $m,r$; Cauchy--Schwarz bounds its expectation by
$(m-r)t+\sqrt{m/2}$. Therefore
\[
 \frac{\phi(S)}m\le\frac12+
 \frac{\sqrt2q^2}{q^3-1}+\frac1{2q}+\frac1{\sqrt2q^2}\longrightarrow\frac12.
\]
This proves an infimum over finite games, without splitting any actual
owner. Coverage forces any scalar multiplier to be at least
$v_c(S)/\phi(S)$, and efficiency identifies that multiplier with its
budget factor, proving the last two claims.
\end{proof}

\begin{proposition}\label{hist:four-shapley}
With exactly two owners in each of two pools, $(6/5)\phi$ covers all
coalitions, is own-report monotone, protects $3H/5$, and spends $6T/5$.
The multiplier $6/5$ is sharp for constant Shapley scaling on this domain.
It fails for some five-owner games.
\end{proposition}
\begin{proof}
Write capacities $(a,b\mid c,d)$, set $m_{ij}=\min(c_i,c_j)$, and
$w_i=v_c(N\setminus\{i\})$. Directly grouping the predecessor subsets gives
\[
 \phi_1=\frac{m_{13}+m_{14}+w_4-m_{23}+w_3-m_{24}+w_2}{12}
          +\frac{T-w_1}{4},
\]
with the other coordinates obtained by pool-preserving relabelings.
Only cross pairs and triples need a nontrivial coverage check.
For $S=\{2,3\}$,
\[
 6\phi(S)=m_{23}+w_4+w_1-m_{14}+3T-w_2-w_3.
\]
Exchange pools if needed so $a+b\le c+d$. Then $T=a+b,w_1=b,w_2=a$,
and
\[
 3T+w_1+w_4-w_2-w_3-m_{14}
 \ge3b+\min(a+b,c)\ge4\min(b,c).
\]
Here $w_3\le a+b$ and $m_{14}\le a$. Thus $\phi(S)\ge5v_c(S)/6$.
For $S=\{1,2,3\}$ the desired inequality, after using efficiency, is
\[
 K=9T-7w_4-w_1-w_2-w_3-m_{14}-m_{24}+m_{13}+m_{23}\ge0.
\]
Put $L=a+b,R=c+d$. If $L\le R$ and $c\ge L$, this expression is
$[L-\min(L,d)]+[a-\min(a,d)]+[b-\min(b,d)]\ge0$.
If $L\le R$ and $c<L$, bounding $w_3,m_{14},m_{24}$ by $L,a,b$ gives
$K\ge6(L-c)+\min(a,c)+\min(b,c)-c\ge0$.
If $L>R$, use $w_4=c,w_3=d$,
$\min(x,c+d)\le\min(x,c)+d$ for $x=a,b$, and $m_{14},m_{24}\le d$;
then $K\ge2c+4d\ge0$. These cases include zeros and ties.
Every other cross pair or triple follows by relabeling.
Proposition~\ref{hist:shapley-half} supplies the claimed fairness floor.
At $(2,1\mid1,1)$,
$\phi=(7/12,1/4\mid7/12,7/12)$, so a pair containing the unit left
owner and a right owner has worth one and payment $5/6$.
At $(1,1\mid1,1,2)$ the Shapley vector is
$(41/60,41/60\mid11/60,11/60,4/15)$, obtained by the same predecessor
formula. The two left owners and the capacity-two right owner have worth
two but receive only $49/25$ after scaling by $6/5$.
\end{proof}

\subsection{The ownership-dependent exponential two-pool rule}
\begin{theorem}\label{hist:exponential}
For arbitrary finite owners in two pools, let $d_i=\min(c_i,T)$.
Call a minimum-total pool short, put
$\rho=d(N_{\rm long})/T$, and $p=\tfrac12e^{1-\rho}$.
Pay $(1-p)c_i$ to short owners and $pd_i$ to long owners, with zero
payments at $T=0$. This rule covers coalitions, is own-report monotone,
protects $M_i/2\ge H_i/2$, and has subsidy at most $1/(2e)$.
The constant is attained for this rule. Every admissible rule, even
without a floor, has worst-case subsidy at least $1/8$ on two pools
with two fixed owners each.
\end{theorem}
\begin{proof}
Capping preserves worth by Lemma~\ref{lem:cap}, and $\rho\ge1$, so the
payments are a convex combination of the two capped pool totals. At a
tie, each owner is paid $c_i/2$, independently of the selected short pool.
Writing $u=\rho-1\ge0$, total payment is
$T[1+ue^{-u}/2]\le T[1+1/(2e)]$, with equality at $(1,1\mid1)$.
Short owners have $M_i=c_i$ and $1-p\ge1/2$. For a long owner, let $a$
be the total of its same-pool companions. If $a\ge T$ its marginal is
zero. Otherwise its companions are uncapped, and with $m=d_i$,
$M_i=T-a\le m\le T$ and $u=(m-M_i)/T$. Since $t e^{-t/T}$ increases
on $[0,T]$, its payment is at least $M_i/2$.

For own-monotonicity, fix a left owner reporting $x$, same-pool remainder
$a$, and other pool total $D>0$. In the long region $x+a\ge D$, put
$m=\min(x,D)$ and $\gamma=\sum_{h\ne i,\,h\in L}\min(c_h,D)$.
Its payment is $\tfrac12e^{1-\gamma/D}m e^{-m/D}$, nondecreasing in $x$.
In the short region $0<x+a\le D$, put $A=x+a$ and
$\rho(x)=\sum_{h\in R}\min(c_h/A,1)$.
Away from finitely many clipping points,
$0\ge\rho'\ge-\rho/A$, and the derivative of $x(1-p)$ is
\[
 1-p+xp\rho'\ge1-p(1+\rho)\ge0,
\]
where the last inequality follows because
$(1+\rho)e^{1-\rho}/2$ decreases from one on $[1,\infty)$.
The payment is continuous at all clipping points and at $A=D$, where
both formulas give $x/2$. If $D=0$ payments remain zero; if $A=0$ the
assigned zero payment is no greater than any later payment.

For the lower bound, use fixed owners $0,1\mid2,3$ and profiles
$c=(1,1\mid1,1)$, $a=(1,1\mid1,2)$, $b=(2,1\mid1,1)$.
At $a$, coverage of $\{0,2\},\{1,2\},\{0,1,3\}$ implies
$2y(a)(N)-y_3(a)\ge4$. The budget $y(a)(N)\le2+2\beta$ gives
$y_3(a)\le4\beta$. At $b$, coverage of
$\{1,2\},\{1,3\},\{0,2,3\}$ similarly gives $y_0(b)\le4\beta$.
Own-monotonicity and the base coalition $\{0,3\}$ give
$1\le y_0(c)+y_3(c)\le y_0(b)+y_3(a)\le8\beta$.
\end{proof}
The resulting interval $[1/8,1/(2e)]$ holds for the two-pool minimax
subsidy with any fixed floor fraction $0\le\alpha\le1/2$.
Neither endpoint is asserted to be the unrestricted optimum. The price
here depends on individual clipping, so Proposition~\ref{hist:aggregate}'s
partition-independent-price lower bound does not apply. The smaller
budget does not dominate Proposition~\ref{hist:four-shapley}'s stronger
$3H/5$ fairness guarantee.

\subsection{Small-stage guarantees and failed exponential extensions}
\begin{proposition}\label{hist:softmax}
For $k=2$ or $3$, put $\rho_j=\sum_{i\in N_j}d_i/T$, $u_j=\rho_j-1$,
$w_j=e^{-u_j}$, $p_j=w_j/\sum_\ell w_\ell$, and $y_i=d_ip_j$.
This rule is own-report monotone, covers coalitions, protects $M_i/k$,
and has exact worst-case subsidy $b_k$, where
$b_ke^{b_k}=(k-1)/e$. The same unit-rate formula fails own-report
monotonicity when $k=4$.
\end{proposition}
\begin{proof}
Coverage is immediate from capping and $\sum p_j=1$.
For an owner with $M_i>0$, same-pool remainder $a<T$, and $m=d_i$, one
has $u_j=(m-M_i)/T$ and $\sum w_j\le k$. Hence
$y_i\ge m e^{-(m-M_i)/T}/k\ge M_i/k$ by the preceding monotonicity of
$t e^{-t/T}$. The assertion is immediate if $M_i=0$.
In the long region of an own report, $T$ is fixed, and differentiation
with respect to $m$ gives $p_j[1-(m/T)(1-p_j)]\ge0$.
In the short region, $T=x+a$, $u_j=0$, and
$0\ge\rho_\ell'\ge-\rho_\ell/T$ off clipping points. Therefore
\[
 \frac{d}{dx}(xp_j)
 \ge p_j\left[1-\sum_{\ell\ne j}p_\ell\rho_\ell\right]\ge0,
\]
since the last bracket is nonnegative precisely when
$\sum_{\ell\ne j}u_\ell e^{-u_\ell}\le1$, and each term is at most
$1/e$. For $k\le3$, $(k-1)/e<1$. Continuity handles clipping points
and the bottleneck seam; the zero boundary is handled by $0\le y_i\le x$.

The subsidy is $F=\sum_j u_je^{-u_j}/\sum_j e^{-u_j}$, with some
$u_s=0$. For $b=b_k$, the elementary maximum
$(u-b)e^{-u}\le e^{-1-b}$ implies
\[
 \sum_j(u_j-b)e^{-u_j}\le-b+(k-1)e^{-1-b}=0.
\]
Thus $F\le b$. Equality holds when $u_s=0$ and every other $u_j=1+b$:
use one unit owner in pool $s$ and three owners of capacity $(2+b)/3<1$
in each other pool. Here $b<1$ follows from $k\le3$.

For failure at $k=4$, use one owner reporting $x$ in the first pool and
three owners of capacity $2/3$ in each other pool. Near $x=1$ the first
owner's payment is
$y(x)=x/[1+3e^{1-2/x}]$, whose derivative at one is
$(1-3/e)/(1+3/e)^2<0$. Thus some strict underreport increases payment.
\end{proof}
A different extension also fails: assigning each nonshort pool price
$e^{-u_j}/k$ and the residual to the short pool is not truthful even for
three pools. At $(x\mid1/2,1/2\mid1/2,1/2)$, the first payment is
$1/3$ when $x=1$, but at $x=99/100$ it is
\[
 x\left[1-\tfrac23 e^{-(1-x)/x}\right]
 \ge x(1-2x/3)=1683/5000>1/3,
\]
using $e^{-q}\le1/(1+q)$. These are obstructions to the stated formulas,
not to all mechanisms with more pools.

\begin{proposition}\label{hist:small-stage-lower}
For $k\ge3$ and two owners per pool, every admissible rule has
\[
 \beta\ge\max\left\{\frac{k-1}{2(k+3)},\frac{k-2}{2k+3}\right\}.
\]
In particular the lower bounds are $1/6$ at $k=3$ and $3/14$ at $k=4$;
the latter exceeds the two-pool upper bound $1/(2e)$.
\end{proposition}
\begin{proof}
Average over pool permutations and within-pool swaps, preserving every
guarantee on the same fixed ownership. The first displayed bound is
Theorem~\ref{thm:unrestricted} with $r=2,t=2$.
For the second, take profile $X$ with one pool $(1,1)$ and $k-1$ pools
$(1,2)$. Let $a$ be the former owners' common payment, and $b,d$ the
unit and capacity-two payments in each mixed pool. Coalitions selecting
all unit owners or both short owners and every capacity-two owner give
\[
 a+(k-1)b\ge1,\quad 2a+(k-1)d\ge2,\quad
 2a+(k-1)(b+d)\le2+2\beta.
\]
Subtracting the first inequality from the budget and comparing with the
second gives $a\ge1-2\beta$. At the profile $Z$ of all mixed pools,
write the payments $b_Z,d_Z$. Then $kb_Z\ge1$ and
$k(b_Z+d_Z)\le3+3\beta$, so $kd_Z\le2+3\beta$.
Raising one unit owner in $X$'s short pool to two produces $Z$;
own-monotonicity gives $d_Z\ge a$. Thus
$k(1-2\beta)\le2+3\beta$. Finally $e>5/2>7/3$ proves
$3/14>1/(2e)$ without numerical approximation.
\end{proof}

\subsection{An obstruction to capped-total-only prices}
\begin{proposition}\label{hist:capped-totals}
Suppose payments have the form $y_i=d_ip_j(\rho_1,\ldots,\rho_k)$,
where $\rho_j=d(N_j)/T$ and the nonnegative prices do not depend on
within-pool capacity distributions. Fixed owner counts and pool labels
may enter the prices. Coalition coverage and own-monotonicity imply,
for $k=q^2$ and $n$ owners per pool with $q\mid n$, the budget bound
\[
 B\ge\frac{k}{2\sqrt{k}-1}>\frac{\sqrt{k}}2.
\]
Prices need not sum to one. Zero reports are allowed.
\end{proposition}
\begin{proof}
At the profile with every owner reporting $1/n$, all ratios equal one
and the budget gives $\sum_jp_j\le B$. Choose a pool $j$ with
$p_j\le B/k$. Within this same fixed ownership, concentrate its capacity
in one owner, leaving its others zero and every other pool unchanged.
At the concentrated report one, the ratio vector is still all ones.
Lowering this owner's report to $x=1/q$ therefore gives
$xp_j(\rho^{\rm low})\le B/k$. A coalition containing this owner and
$n/q$ owners in each other pool has worth $x$ and payment
$x\sum_\ell p_\ell(\rho^{\rm low})$, so that sum is at least one.
The low-profile total payment is consequently at least
\[
 xp_j+\sum_{\ell\ne j}p_\ell
 \ge1-(1-x)p_j\ge1-\frac{(1-x)B}{kx}.
\]
Its budget is at most $Bx$, proving
$B\ge[x+(1-x)/(kx)]^{-1}=k/(2\sqrt{k}-1)$.
No independently strategic clones or comparison across ownerships is
used here. The restriction to ratio-only prices is essential; the
owner-coordinate Banzhaf rule uses the full within-pool distribution.
\end{proof}

\section{Effective capacities: a successful class and its boundary}
\label{net:effective}
These results analyze an earlier continuous-random-cut rule.  The general
hybrid in Theorem~\ref{prod:upper} does not require its structural incentive
argument.  All networks below are finite directed acyclic $s$--$t$ graphs,
every edge lies on an $s$--$t$ path, and each edge has one distinct fixed owner.
The topology and owners are fixed during every report comparison.  Parallel
edges are permitted unless a graph is explicitly called simple.
Write $F(c)$ for maximum flow, $T=F(c)$, and $v_c(S)=F(c\ind_S)$.
Let $\mathcal C$ be the inclusion-minimal structural edge cutsets and
$K=|\mathcal C|$; this family is independent of capacities, including zero
capacities.  Then $F(c)=\min_{C\in\mathcal C}\sum_{e\in C}c_e$.

Define $D_e(c)$ as the greatest flow through edge $e$ among all feasible
$s$--$t$ flows, without requiring their total value to be maximal.
For independent uniform $U_e\in[0,1]$, choose the first cut $W$ minimizing
$\sum_{e\in C}D_e(c)U_e$, using a fixed priority on $\mathcal C$, and put
\begin{equation}\label{net:base-rule}
 z_e(c)=D_e(c)\Prob(e\in W).
\end{equation}
These are deterministic expectations; the uniforms are an integration device.
A useful graph formula, also proving continuity, is
$D_e=\min\{c_e,F_{s,\mathrm{tail}(e)},F_{\mathrm{head}(e),t}\}$,
omitting a prefix or suffix with coincident terminals.  An upstream path and
a downstream path cannot share a vertex, since that would create a cycle
through $e$.  Scaled maximum prefix and suffix flows therefore join through
$e$ and attain the formula; the reverse bounds are immediate.

\begin{lemma}[Effective-capacity preservation and budget]
\label{net:preservation}
On every such DAG, $D\le c$, each $D_e$ is nondecreasing in every report,
$0\le D_e\le T$, and replacing $c$ by $D(c)$ preserves the complete feasible
flow set of every coalition.  Consequently \eqref{net:base-rule} is
nonnegative and coalition-covering.  For $K\ge2$ let $R_K>e$ solve
$R_K\ln(R_K/e)=\ln K$ and set
\[
 U(K)=R_K+1/\ln R_K.
\]
Then $\sum_ez_e(c)\le U(K)T$; for $K=1$ the rule is efficient.
Here $U(K)\sim\ln K/\ln\ln K$.  No incentive assertion is made in this lemma.
\end{lemma}
\begin{proof}
Inclusion of feasible flow sets proves monotonicity of $D$.  A DAG has no
circulations, so path decomposition bounds edge flow by total flow and hence
by $T$.  Every coalition flow is grand-network feasible, so it respects all
$D_e$; the reverse inclusion uses $D\le c$.
For every selected cut and every coalition,
$\sum_{e\in S\cap W}D_e\ge F(D\ind_S)=v_c(S)$.
The equality follows from preservation and max-flow/min-cut on the fixed
graph after giving nonmembers capacity zero.  Taking expectations proves
coverage.  If $T=0$, all $D_e$ and payments are zero.

For $T>0$ put $a_e=D_e/T\in[0,1]$ and $\rho_C=\sum_{e\in C}a_e$.
Some cut $C_0$ has $\rho_{C_0}=1$.  Its random score is at most one, so the
winning cut has $X_W=\sum_{e\in W}a_eU_e\le1$.
For $t>0$, concavity of $x\mapsto x^{a_e}$ gives
\[
 \E e^{-tX_C}\le\prod_{e\in C}(\E e^{-tU_e})^{a_e}
 \le t^{-\rho_C},\qquad
 \Prob(X_C\le1)\le e^t t^{-\rho_C}.
\]
Set $t=\rho_C$.  Since $(e/s)^s$ decreases for $s>1$, a union bound gives
$\Prob(\rho_W>s)\le\min\{1,K(e/s)^s\}$ for $s\ge e$.
Beyond $R_K$ the logarithm of the second bound has derivative
$-\ln s\le-\ln R_K$.  Therefore
\[
 \frac{\sum_ez_e}{T}=\E\rho_W
 \le R_K+\int_{R_K}^\infty e^{-(s-R_K)\ln R_K}\,ds=U(K).
\]
Cut scores may overlap; no independence between cuts was used.
With one cut, its $D$-sum is $T$.  Finally the defining equation implies
$R_K\sim\ln K/\ln\ln K$, and $1/\ln R_K=o(R_K)$.
\end{proof}

\subsection{Series-parallel networks and co-cut independence}
A two-terminal series-parallel (TTSP) network is built from directed edge
leaves by series and parallel composition of components with disjoint edges
and internal vertices.  Fix its decomposition tree.  Component flow satisfies
$F_{A\parallel B}=F_A+F_B$ and $F_{A;B}=\min(F_A,F_B)$.
Minimal cuts form disjoint alternatives at a series node and unions of one
cut from each child at a parallel node.  Thus
$K_{A;B}=K_A+K_B$ and $K_{A\parallel B}=K_AK_B$.

\begin{theorem}[TTSP incentives]\label{net:ttsp}
On every fixed TTSP network, the rule \eqref{net:base-rule} is own-report
monotone at every nonnegative real profile, and hence satisfies coverage and
the $U(K)T$ budget bound with binding-underreport incentives.
\end{theorem}
\begin{proof}
For an edge $e$, let $L_e$ be the minimum of $F_B(c)$ over the sibling
components $B$ at all series ancestors of $e$, with empty minimum $+\infty$.
Then
\begin{equation}\label{net:ttsp-cap}
 D_e=\min(c_e,L_e).
\end{equation}
For attainability, route exactly this amount through the leaf.  At a parallel
ancestor use zero flow in its other child; at a series ancestor route the same
amount through the sibling by scaling its maximum flow.  Each ancestor's
bound is respected.  Conversely, every path carrying flow through $e$ crosses
each such series sibling, giving all the displayed upper bounds.

If distinct edges $e,h$ share a minimal cut, their lowest common ancestor
must be parallel: at a series node a minimal cut lies wholly in one child.
Formula~\eqref{net:ttsp-cap} shows that $D_h$ depends on $c_e$ only if $e$
lies in a sibling at a series ancestor of $h$.  This is impossible for a pair
sharing a cut.  Thus $D_h$ is independent of $c_e$ for co-cut pairs.

Fix $c_{-e},U_{-e}$ and write $x=c_e$, $d(x)=D_e(c)$.  Put
\[
 A=\min_{C\ni e}\sum_{h\in C\setminus\{e\}}D_hU_h,\qquad
 B(x)=\min_{C\not\ni e}\sum_{h\in C}D_hU_h,
\]
with $B=+\infty$ if no excluding cut exists.
Co-cut independence makes $A$ constant in $x$; $B$ and $d$ are nondecreasing.
Integration over the single uniform $U_e$ gives
\begin{equation}\label{net:conditional}
 D_e\Prob(e\in W\mid U_{-e})=\min\{d(x),(B(x)-A)_+\}.
\end{equation}
When $d>0$ a tie between the containing and excluding scores occurs at at
most one value of $U_e$, hence has measure zero.  The formula also holds for
$d=0$ and $B=+\infty$ under their evident conventions.
Its right side is nondecreasing in $x$.  Average over $U_{-e}$.
\end{proof}
The same proof applies to any fixed cut representation with a normalization
that preserves every coalition, is coordinatewise nondecreasing and bounded
by $T$, and has co-cut independence.  The condition is sufficient, not necessary.
Serial layers of parallel edges recover multi-glove games.  Their lower
bound and Theorem~\ref{net:ttsp} give sharp worst-case TTSP order
$(1+o(1))\ln K/\ln\ln K$.  This lower construction is inherited; the general
upper and simple-DAG lower are now Theorems~\ref{prod:upper} and~\ref{dag:sharp}.

\begin{proposition}[A global cap does not replace effective capacities]
\label{net:global-cap}
Globally replacing each edge by $\min(c_e,T)$ before applying the continuous
random-cut rule can violate own-report monotonicity even on TTSP graphs
with three structural cuts and arbitrarily many owners.
\end{proposition}
\begin{proof}
Use the flow function
$F(x,a,b,d_1,\ldots,d_n)=\min\{x+\min(a,b),\sum_jd_j\}$.
For $n\ge4$, set $\epsilon=1/10$, $a=\epsilon$, $b=n+\epsilon$,
$d_j=1$, and compare $x_0=n-2\epsilon$ with $x_1=n-\epsilon$.
Then $T_0=n-\epsilon$, $T_1=n$.  Under the global cap only $b$ is changed,
to $T$.  If $S_n$ is a sum of $n$ independent unit uniforms and
$A_T=\min(\epsilon U,TV)$ is independent, the focal payment is
\[
 z(x,T)=\E\min\{x,(S_n-A_T)_+\}.
\]
Write $\delta=n-T\in\{\epsilon,0\}$, $t=\epsilon+\delta$, so $x=n-t$.
For $0\le u\le1$, simplex integration gives
$\E(u-S_n)_+=u^{n+1}/(n+1)!$; also $n-S_n$ has the same law as $S_n$.
Since $0\le A_T\le\epsilon\le t<1$, subtraction of the two positive parts
therefore yields
\[
 z(x,T)=\frac n2-\E A_T+
 \frac{\E A_T^{n+1}-\E(t-A_T)^{n+1}}{(n+1)!}.
\]
The survival law $(1-u/\epsilon)(1-u/T)$ on $[0,\epsilon]$ gives
\[
 \E A_T^r=\frac{\epsilon^r}{r+1}
 -\frac{r\epsilon^{r+1}}{(r+1)(r+2)T},\quad r\ge1.
\]
Thus the mean increases by $\epsilon^3/[6n(n-\epsilon)]$ between profiles.
Bounding the two positive remainder terms separately gives
\[
 z(x_1,T_1)-z(x_0,T_0)
 \le-\frac{\epsilon^3}{6n(n-\epsilon)}
       +\frac{2(2\epsilon)^{n+1}}{(n+1)!}<0.
\]
For the last inequality, the ratio of the positive bound to
$\epsilon^3/(6n^2)$ is
$r_n=6n^2 2^{n+2}\epsilon^{n-2}/(n+1)!$.
It equals $64/125<1$ at $n=4$ and decreases, since
$r_{n+1}/r_n=2\epsilon(n+1)^2/[n^2(n+2)]<1$.
Effective capping instead fixes the detour capacity at $a$, as required by
Theorem~\ref{net:ttsp}.  This is a failure of a specific formula, not a
lower bound for all mechanisms.
\end{proof}

\subsection{A bundled Wheatstone bridge works, but composition fails}
Let $H_k$ have vertices $s,u,v,q$ and edges
$x:s\to u$, $y:s\to v$, $z_j:u\to q$ ($1\le j\le k$),
$w:v\to q$, and $h:u\to v$.  These are distinct fixed owners.
Its four minimal cuts are
$\{x,y\}$, $\{z_1,\ldots,z_k,w\}$,
$\{y,z_1,\ldots,z_k,h\}$, and $\{x,w\}$.

\begin{proposition}[Bundled bridge]\label{net:bundle}
For every $k\ge1$, effective-capacity payments are own-report monotone on
$H_k$, at every nonnegative real profile, and pay at most $U(4)T$.
\end{proposition}
\begin{proof}
Writing $C=\sum_jz_j$, the effective capacities are
\[
 \begin{aligned}
 d_x&=\min(x,C+\min(h,w)),&d_y&=\min(y,w),\\
 d_{z_j}&=\min(z_j,x),&d_w&=\min(w,y+\min(x,h)),\\
 d_h&=\min(h,x,w).
 \end{aligned}
\]
Each formula is the maximum amount that can be routed through its specified
edge: the displayed upstream/downstream bounds are necessary and attained
by routing that amount along the indicated direct and bridge paths.
Put $X=d_xU_x$, $Y=d_yU_y$, $Z=\sum_jd_{z_j}U_{z_j}$,
$W=d_wU_w$, $H=d_hU_h$.  Formula~\eqref{net:conditional}, applied to the
four cuts and integrating only the focal uniform, gives the integrands.
Here the containing intercept may depend on the report: only the conditional
probability identity is used, and monotonicity is checked directly below.
The resulting expressions are
\[
 \begin{aligned}
 q_x&=\min(d_x,Z+\min((W-Y)_+,H)),\\
 q_w&=\min(d_w,Y+\min((X-Z)_+,H)),\\
 q_y&=\min(d_y,(W-\min((X-Z)_+,H))_+),\\
 q_{z_j}&=\min(d_{z_j},(X-Z_{-j}-\min((W-Y)_+,H))_+),\\
 q_h&=\min(d_h,(X-Z)_+,(W-Y)_+).
 \end{aligned}
\]
Here $Z_{-j}$ excludes the focal summand.  For example, for $z_j$ the
least containing score without $z_j$ is $Z_{-j}+\min(W,Y+H)$ and the least
excluding score is $X+\min(Y,W)$, which gives the stated cancellation.
The other expressions follow by the same subtraction of the four cut scores.
When $x$ grows, only $d_x,Z,W,H$ can increase and $Y$ stays fixed; all occur
positively in $q_x$.  When $w$ grows, $d_w,Y,X,H$ can increase and $Z$ is
fixed.  When $y$ grows, only $d_y,W$ can increase in its expression.
When $z_j$ grows, only $d_{z_j},X$ can increase; $Z_{-j},W,Y,H$ are fixed.
When $h$ grows, $d_h,X,W$ can increase and $Y,Z$ remain fixed.
Every integrand is therefore own-report monotone.  Average and apply
Lemma~\ref{net:preservation}.
\end{proof}
This is an unbounded-owner, four-cut class outside TTSP, not a theorem about
arbitrarily many interacting bridge structures.

\begin{theorem}[Actual expected-payment failures]
\label{net:wheatstone-failure}
There are fixed edge-owned DAGs with five structural cuts and arbitrarily
many owners on which \eqref{net:base-rule} decreases an owner's expected
payment when only its own report increases.  There are also simple DAGs
with growing vertex and owner counts with the same property.  In both
constructions the two grand values are one and all capacities are positive.
\end{theorem}
\begin{proof}
Add an edge $p:s\to q$ in parallel with $H_k$, then two parallel edges
$r_1,r_2:q\to t$.  The structural cuts are the four bridge cuts with $p$
adjoined, and $\{r_1,r_2\}$.  Set
\[
 \epsilon=1/20,\quad x_0=9/10,\quad x_1=19/20,\quad y=\epsilon,
 \quad z_j=w=h=p=2,\quad r_1=r_2=1/2.
\]
Only $x$ changes.  Both values are one.  For $x\in[x_0,x_1]$ the effective
capacities are
\[
 D_x=x,\ D_y=\epsilon,\ D_{z_j}=D_h=x,\ D_w=x+\epsilon,
 \ D_p=1,\ D_{r_i}=1/2.
\]
The incoming $x$ and $x+y$ cuts and downstream capacity one give the upper
bounds; routes through an individual $z_j$, through $h,w$, and through $p$
attain them.  Let independent uniforms define
\[
 \begin{gathered}
 Y=\epsilon U_y,\quad W=(x+\epsilon)U_w,\quad H=xU_h,\quad Z=x\sum_{j=1}^kU_{z_j},\\
 P=U_p,\quad R=(U_{r_1}+U_{r_2})/2,\quad A=\min(Y,W).
 \end{gathered}
\]
The containing-cut score without $x$ is $P+A$; the least excluding score
is $\min\{P+Z+\min(W,Y+H),R\}$.  Thus its actual payment is $\E q_k(x)$,
where
\begin{equation}\label{net:wheatstone-q}
 q_k(x)=\min\{x,[\min\{Z+\min(H,(W-Y)_+),R-P-A\}]_+\}.
\end{equation}
Compare $I(x)=\E\min\{x,(R-P-A)_+\}$.  The integrands agree when $Z\ge1$
and are otherwise between zero and one, so
\begin{equation}\label{net:ideal-error}
 I(x)-\Prob(Z<1)\le z_x(x)\le I(x),\qquad
 \Prob(Z<1)\le\frac1{k!x^k}.
\end{equation}
The second inequality bounds the relevant portion of the unit cube by the
nonnegative simplex of coordinate sum below $1/x$; it is not an equality.

For $b=x+\epsilon$ write
\[
 \begin{gathered}
 M_0(b)=1,\qquad M_j(b)=\E A^j=\frac{\epsilon^j}{j+1}
       -\frac{j\epsilon^{j+1}}{(j+1)(j+2)b},\quad j\ge1,\\
 P_4(t,b)=\sum_{j=0}^4\binom4j(-1)^jt^{4-j}M_j(b).
 \end{gathered}
\]
For $0\le a\le\epsilon$, convolution of independent uniforms of widths
$1/2,1/2,1$ (or simplex inclusion-exclusion) gives
\[
 \E(R-P-a)_+=\frac{(1-a)^4-2(1/2-a)^4}{6}.
\]
At $x_0,x_1$, its truncated part is
$\E(R-P-a-x)_+=(1-x-a)^4/6$.
Consequently
\[
 I(x)=\frac{P_4(1,b)-2P_4(1/2,b)-P_4(1-x,b)}6,
 \qquad I(x_0)-I(x_1)=\frac{553}{136800000}.
\]
By \eqref{net:ideal-error}, the actual payment drop is at least
$553/136800000-(10/9)^k/k!$.
At $k=10$ this is
$97726081746439/30050502681578400000>0$; the error decreases thereafter.
This proves the first assertion for every $k\ge10$.

For the simple version, replace each $z_j$ by a private two-edge path with
both capacities two and distinct owners.  Replace each $r_i$ by a private
two-edge path with both capacities $1/2$.  The other edges are unchanged.
There are $k+7$ vertices and $2k+9$ owners; there are no parallel edges.
The effective capacities of each pair equal the original effective capacity,
and \eqref{net:wheatstone-q} holds with
\[
 Z=x\sum_{j=1}^k\min(U_j,V_j),\qquad
 R=\bigl(\min(U'_1,V'_1)+\min(U'_2,V'_2)\bigr)/2.
\]
Each minimum has density $2(1-u)\le2$, so the error in
\eqref{net:ideal-error} is now at most $2^k/(k!x^k)$.
Each downstream path has density $4-8r$ on $[0,1/2]$.
Convolution gives
\[
 F_R(r)=8r^2-\tfrac{32}{3}r^3+\tfrac83r^4\quad(0\le r\le1/2),
 \qquad 1-F_R(r)=\tfrac83(1-r)^4\quad(1/2\le r\le1).
\]
Since $P$ is uniform, $g(a)=\E(R-P-a)_+=\E(R-a)_+^2/2$.
Integrating $g''=1-F_R$, with $g(0)=5/72$, $g'(0)=-1/3$, yields
\[
 g_0(a)=\tfrac5{72}-\tfrac a3+\tfrac{a^2}2-\tfrac23a^4
               +\tfrac8{15}a^5-\tfrac4{45}a^6\quad(0\le a\le1/2),
 \qquad g_1(a)=\tfrac4{45}(1-a)^6\quad(1/2\le a\le1).
\]
Using $A\le1/20$ and $A+x\ge9/10$ therefore gives
\[
 I(x)=\E g_0(A)-\tfrac4{45}\E(1-x-A)^6,
\]
where every moment is the rational $M_j(b)$ above.  Expansion gives
\[
 I(x_0)=\frac{11809546127}{191520000000},\quad
 I(x_1)=\frac{14797313}{240000000},\quad
 I(x_0)-I(x_1)=\frac{1290353}{191520000000}.
\]
The actual drop is at least
$1290353/191520000000-(20/9)^k/k!$.
At $k=14$ this equals
\[
 \frac{700533736635095043154291}{118415105665157830969440000000}>0,
\]
and it increases thereafter.  The structural cut count is
$2^{k+1}+6$: each of the two cut types involving the $z$ bundle chooses one
edge of each path, the other two upstream types are unchanged, and the
final pair contributes four cuts.  This proves the simple-DAG assertion.
\end{proof}
The obstruction concerns an actual integrated payment, not only a conditional
integrand.  Coverage and the cut-count budget still hold by
Lemma~\ref{net:preservation}; incentives fail.  In particular, the successful
bundled bridge rule is not preserved by arbitrary series/parallel composition
with other modules.  This does not preclude other mechanisms on these graphs.

\section{The own-report envelope: a linear upper bound and growing obstruction}
\label{net:envelope}
For the rule \eqref{net:base-rule}, its actual own-report monotone envelope is
\[
 y_e(c)=\sup_{0\le t\le c_e}z_e(t,c_{-e}).
\]
It is automatically own-report monotone and dominates the covering vector $z$.
The supremum is attained and defines a continuous function: the graph formula
above makes every $D_e$ continuous.  The conditional formula
$z_e=\E\min\{D_e,(B-A)_+\}$ makes $z_e$ continuous by dominated
convergence, with $A,B$ the containing/excluding scores at the given profile.
Writing $t=\lambda c_e$ with $0\le\lambda\le1$ gives a maximum of a
continuous function over a fixed compact interval, which is continuous.
These regularity observations are not needed for the budget proof.
The next bounds concern this specific repair; they are not bounds on the
optimal subsidy among all truthful mechanisms.

\begin{lemma}[A shared-uniform comparison]\label{net:uniform-comparison}
For finite nonnegative coefficient vectors $a,b$ and independent unit
uniforms, if $A=\sum_i a_i$, $B=\sum_i b_i$, then
\[
 A\Prob\Bigl(\sum_i a_iU_i\le\sum_i b_iU_i\Bigr)\le B.
\]
\end{lemma}
\begin{proof}
If $A\le B$, the assertion is immediate.  Otherwise put
$N=\sum_i(b_i-a_i)_+$ and $W=\sum_i|a_i-b_i|=A-B+2N>0$.
Reflect uniforms whose coefficient $a_i-b_i$ is negative.  The event becomes
$\sum_i|a_i-b_i|V_i\le N$, where $N<W/2$.
The resulting positive weighted uniform sum has a symmetric unimodal density
about $W/2$.  To check this directly, a centered uniform is even and
nonincreasing on the positive half-line.  If $f$ has these properties,
convolution with the centered uniform on $[-r,r]$ has derivative
$[f(x+r)-f(x-r)]/(2r)\le0$ for $x\ge0$ almost everywhere and is even;
absolute continuity proves the assertion by induction.
Its CDF is convex on $[0,W/2]$, vanishes at zero, and equals $1/2$ at $W/2$.
Thus its value at $N$ is at most $N/W$.
Since $N\le B$ and $A>B$, $N/(A-B+2N)\le B/(A+B)$.
Multiplication by $A$ proves the claim, including $B=0$.
\end{proof}

\begin{theorem}[Linear budget for the actual envelope]\label{net:envelope-upper}
For every fixed edge-owned DAG and every nonnegative real profile,
$\sum_ey_e(c)\le KT(c)$.  More precisely, if $T>0$ and
$d_C=\sum_{h\in C}D_h(c)$, then
\begin{equation}\label{net:envelope-coordinate}
 y_e(c)\le T D_e(c)\sum_{C\ni e}\frac1{d_C}.
\end{equation}
\end{theorem}
\begin{proof}
Fix $c$, with $T>0$, and a minimum cut $A$ for $D(c)$, so $d_A=T$.
Write $Z_C=\sum_{h\in C}D_h(c)U_h$ using one common uniform vector.
If $D_e=0$, all own underreports pay zero.
Reports $t\ge D_e$ leave every feasible flow unchanged, so it suffices to
consider $0<t\le D_e$.  Set $x=t/D_e$ and $c^t=(t,c_{-e})$.
Feasible-set inclusion and scaling a flow attaining each effective cap give
\[
 xD_h(c)\le D_h(c^t)\le D_h(c),\qquad D_e(c^t)=t.
\]
If the selected cut at $c^t$ is $C$, its score is no larger than that of $A$,
and hence $xZ_C\le Z_A$.  Lemma~\ref{net:uniform-comparison}, with totals
$xd_C$ and $T$, yields
\[
 t\Prob(W(c^t)=C)\le TD_e/d_C.
\]
Sum over cuts containing $e$ before taking its own supremum, obtaining
\eqref{net:envelope-coordinate}.  Summing the resulting bounds gives
$\sum_ey_e\le T\sum_C\sum_{e\in C}D_e/d_C=KT$.
For $T=0$, all effective capacities vanish.  Coverage and incentives follow
from domination and the expanding supremum domain.
\end{proof}
For comparison with earlier accounting, let $L$ be the longest path length,
$h(C)=\max_P|P\cap C|$, and $H=\sum_Ch(C)$.
Since $H\ge K$, Theorem~\ref{net:envelope-upper} immediately implies the
archived weaker bound $2HT$.  Also $H\le KL$ and $L\le K$: for each edge
of a fixed path, deleting that edge and every edge outside the path contains
a minimal cut whose trace on the path is exactly that edge, giving $L$
distinct cuts.  Hence the former $2KLT$ and $2K^2T$ conclusions are preserved.
The earlier structural estimate $H=O(K^{3/2})$ and its sharp Fibonacci-family
example concern the statistic $H$; their archived proofs are independent
of the improved envelope estimate.  They do not supply an envelope lower
bound and are not repeated here.  The sharp general mechanism theorem uses
a different payment rule.

\begin{theorem}[Growing lower bounds for the envelope]
\label{net:envelope-lower}
There are explicit fixed-owner DAG families at positive profiles with $T=1$
and growing cut count $K$ for which
\[
 \sum_ey_e=\Omega\!\left(\frac{K^{1/401}}{\ln K}\right).
\]
There are also explicit simple-DAG families with
$\sum_ey_e=\Omega(K^{1/19201}/\ln K)$.
Both statements apply to every own-report monotone pointwise majorant of $z$.
They do not apply to arbitrary truthful coalition-covering mechanisms.
\end{theorem}
\begin{proof}
First take an integer $r\ge4$ and set
\[
 N=2^r,\qquad m=400r,\qquad k=10m,\qquad
 a=6/5,\qquad b=3/10,\qquad t=b/m.
\]
Place $N$ modules in series.  Module $j$, with entry $s_j$, exit $q_j$,
vertex $A_j$ and vertices $B_{j,i}$, has edges and capacities
\[
 \begin{array}{c|c|c}
 \text{edge}&\text{endpoints}&\text{capacity}\\\hline
 e_j&s_j\to A_j&1\\
 y_{j,i}&s_j\to B_{j,i}&a/m\\
 g_{j,i}&A_j\to B_{j,i}&2\\
 h_{j,i}&B_{j,i}\to q_j&2\\
 Z_{j,\ell}&A_j\to q_j&2\quad(1\le\ell\le k).
 \end{array}
\]
After the modules add a final parallel bundle of $100000$ edges, each of
capacity $1/100000$.  All listed edges have distinct original owners.
For each $j$ separately compare the common original profile with the profile
lowering only $e_j$ to $t$.

A module has $2^{m+1}$ minimal cuts.  A cut containing $e$ chooses one of
$y_i,h_i$ in each branch.  A cut excluding $e$ contains all $Z$ edges and
chooses either $h_i$ or both $y_i,g_i$ in each branch.
These are necessary to block the direct and branch paths, sufficient to block
all paths, and minimal because each chosen edge has an unblocked-path witness
when removed.  A series-composition minimal cut lies in one component, so
\begin{equation}\label{net:aggregate-K}
 K=N2^{m+1}+1=2^{401r+1}+1.
\end{equation}
If a module's focal report is $q\in\{1,t\}$, its raw value is $a+q$:
the source cut is an upper bound, attained by routing $q$ along $e,Z$ and
each $y_i$ along $h_i$.  Thus all modules remain above one, and the final
bundle fixes global value one at every comparison.
The effective capacities inside this module are
\[
 D_e=D_{g_i}=D_{Z_\ell}=q,\quad D_{y_i}=a/m,\quad
 D_{h_i}=\min(1,q+a/m).
\]
Incoming cuts and the final bound give necessity.  Prefix/suffix routes
attain each bound; all other modules retain raw value greater than one, so
changing the focal report changes no effective capacity in those modules.

Let
\[
 A_q=\sum_{i=1}^m\min((a/m)U_{y_i},D_{h_i}U_{h_i}),
 \qquad Z_q=q\sum_{\ell=1}^kU_{Z_\ell}.
\]
The least containing score is $qU_e+A_q$; every excluding score is at least
$Z_q$.  For $0<u\le v$, integrating the product survival function gives
$\E\min(uU,vV)=u/2-u^2/(6v)$.  Therefore
\[
 \E A_{\rm high}=\frac35-\frac6{25m}\ge\frac{59}{100},
 \qquad \E A_{\rm low}=\frac{11}{25}.
\]
The summands lie in $[0,a/m]$.  The bounded-sum inequality
\[
 \Prob\!\left(\sum_i(X_i-\E X_i)\ge s\right)
 \le \exp\!\left(-\frac{2s^2}{\sum_i\operatorname{range}(X_i)^2}\right),
\]
and its lower-tail version give
\[
 \Prob(A_{\rm high}\le.53)\le e^{-m/200},\quad
 \Prob(A_{\rm low}\ge.47)\le e^{-m/800},\quad
 \Prob(Z_{\rm low}\le.53)\le e^{-9409m/4500}\le e^{-2m}.
\]
Here $\E Z_{\rm low}=3/2$ and its sum of squared ranges is $9/(10m)$;
high $Z$ stochastically dominates low $Z$.
For completeness, the bounded-sum inequality follows by convexity reducing
each variable to its endpoints: the centered Bernoulli log-mgf has second
derivative at most one quarter, so it is at most $\lambda^2/8$ after unit
scaling.  Multiply the mgf bounds and optimize the exponential Markov bound.
The final-bundle random score $Q$ has mean $1/2$ and variance $1/1200000$.
Chebyshev gives $\Prob(Q\notin[.49,.51])\le1/120$.

At the profile lowering $e_j$, union bound the event that its $A_{\rm low}$
is below $.47$, every other $A_{\rm high}$ exceeds $.53$, all $Z$ exceed
$.53$, and $Q\in[.49,.51]$.  Its failure probability is at most
\[
 e^{-m/800}+Ne^{-m/200}+Ne^{-2m}+1/120
 \le1/4+1/16+1/16+1/120=23/60.
\]
The last inequalities follow from $m=400r$, $r\ge4$, and $e>2$.
On the good event the focal containing score is at most $.47+t<.48$,
whereas its excluding and all other module scores exceed $.53$ and the
final score is at least $.49$.  The global winner therefore contains $e_j$
for every realization of its own uniform.  Thus
\[
 z_{e_j}(c^{t,j})\ge37t/60>t/2,\qquad
 \sum_ey_e(c)\ge Nt/2=\frac{3\,2^r}{8000r}.
\]
Together with \eqref{net:aggregate-K} this proves the first bound.
The separate underreport comparisons justify different own suprema at the
one original profile; they do not sum actual payments from different profiles.

For the simple construction set $m=1600r$, retaining $N,k,a,b,t$ as above.
Replace every $y_i,h_i,Z_\ell$ by a private two-edge path, each copy with
the original capacity and its own fixed owner.  Keep $e,g_i$ unsplit.
Use $2000$ private final two-edge paths in parallel, each copy of capacity
$1/2000$.  All private vertices are distinct.  There are
$N(25m+1)+4000$ owners, and the graph is simple and acyclic.
Raw values and effective-capacity bounds are unchanged for each copy.
The random capacity of a replaced path is its cap times
$V=\min(U_1,U_2)$, with mean $1/3$, variance $1/18$, and survival $(1-v)^2$.
Different paths give independent variables.  Direct integration yields
\[
 \E\min(uV_1,vV_2)=u\left[\frac13-\frac{u}{6v}
                                      +\frac{u^2}{30v^2}\right],\quad u\le v.
\]
Thus
\[
 \E A_{\rm high}=\frac25-\frac6{25m}+\frac{(6/5)^3}{30m^2}
       \ge.39,\qquad \E A_{\rm low}=166/625.
\]
The same range bounds and bounded-sum inequality give
\[
 \begin{gathered}
 \Prob(A_{\rm high}\le.36)\le e^{-m/800},\quad
 \Prob(A_{\rm low}\ge.30)\le e^{-1849m/1125000}\le e^{-m/800},\\
 \Prob(Z_{\rm low}\le.36)\le e^{-1024m/1125}\le e^{-m/2}.
 \end{gathered}
\]
Now $\E Z_{\rm low}=1$, with squared ranges still $9/(10m)$.
The final score has mean $1/3$ and variance $1/(18\cdot2000)$, so
$\Prob(Q\notin[47/150,53/150])\le5/72$.
The corresponding good event fails with probability at most
\[
 1/16+1/256+1/16+5/72=457/2304<1/2.
\]
Indeed $Ne^{-m/800}\le2^{-r}$, $e^{-m/800}\le2^{-2r}$, and
$Ne^{-m/2}\le2^{-r}$.  The focal score is at most
$.30+t<.31<47/150$; its excluding and all other module scores exceed
$.36>53/150$.  Hence it wins with probability greater than one half and
\[
 \sum_ey_e(c)\ge\frac{3\,2^r}{32000r}.
\]
Containing cuts now have four choices per branch.  Excluding cuts have
four per branch and two choices on each of the $k$ direct paths.  The final
paths contribute $2^{2000}$ cuts.  Therefore
\[
 K=N4^m(1+2^k)+2^{2000}
   =2^{19201r}+2^{3201r}+2^{2000},\qquad
 2^{19201r}\le K\le3\,2^{19201r},
\]
which proves the second asymptotic bound.
Finally, if a rule $p$ is own-report monotone and $p\ge z$ pointwise, then
$p_e(c)\ge p_e(t,c_{-e})\ge z_e(t,c_{-e})$ for every $t\le c_e$;
taking suprema gives $p\ge y$.
\end{proof}

\begin{remark}[Why this is a mechanism-specific separation]
\label{net:fixed-cut}
Fix a cut $C$ from topology alone and pay
$p_e=\ind\{e\in C\}\min(c_e,T(c))$.
Both arguments of the minimum are own-report nondecreasing.  For a coalition,
\[
 \sum_{e\in S}p_e\ge
 \min\left\{\sum_{e\in C\cap S}c_e,T\right\}\ge v_c(S),
\]
because either a capped summand already reaches $T$, or none is capped.
Total payment is at most $|C|T$.  The two lower-bound families above have
fixed final cuts of sizes $100000$ and $2000$, respectively, and therefore
admit constant-factor truthful covering rules on their entire report domains.
The envelope nevertheless grows without bound on them.  Its exact worst-case
growth between Theorems~\ref{net:envelope-upper} and~\ref{net:envelope-lower}
remains open.  The latter blocks the sharp-order target for every monotone
pointwise majorant of the original rule, while Theorem~\ref{prod:upper}
attains that target by departing from this class.
\end{remark}

\section{A diagnostic for projection onto the full efficient core}
\label{net:projection}
The following examples distinguish the full efficient core from a smaller
set of dual-price-generated allocations.  They concern the specific rule
that projects grand-minus-complement marginal payments onto that full core.
They are source-comparison diagnostics, not new general subsidy lower bounds.

For four fixed owners $A=\{0,1\}$, $B=\{2,3\}$, initially put
\[
 v_c(S)=\min\left\{\sum_{i\in S\cap A}c_i,
                         \sum_{j\in S\cap B}c_j\right\},\qquad
 \mathcal C(c)=\{x:x(N)=T,\ x(S)\ge v_c(S)\text{ for every }S\}.
\]
Singleton constraints imply nonnegativity.  Define
$U_i(c)=T-v_c(N\setminus\{i\})$ and let $P(c)$ be the unique minimizer
of $\frac12\|x-U(c)\|_2^2$ over $\mathcal C(c)$.
Compare the profiles
\[
 L=(1/2,1,1,1),\qquad H=(1/2,1,2,1),
\]
which differ only in owner 2's own capacity.

\begin{proposition}[Linear full-core projection failure]
\label{net:linear-projection}
The stated projection pays $P_2(L)=1/4$ and $P_2(H)=0$.
These games arise from four unique resource types and four fixed linear
production activities.  Their dual-price allocation set, however, is already
a singleton at the low profile, so this example does not refute projection
onto that smaller set.
\end{proposition}
\begin{proof}
At $L$, $T=3/2$ and $U=(1/2,1,1/2,1/2)$.
Orthogonal projection onto the efficiency hyperplane subtracts $1/4$ from
every coordinate, giving
\[
 x^L=(1/4,3/4,1/4,1/4).
\]
It lies in the full core.  Empty and single-side coalitions have value zero.
The mixed-pair values are $1/2,1/2,1,1$, equal to their payments.
The triples $\{0,1,j\}$ have value one and payment $5/4$;
$\{0,2,3\}$ has value $1/2$ and payment $3/4$;
$\{1,2,3\}$ has value one and payment $5/4$.
The grand value and payment are $3/2$.
Thus the hyperplane minimizer is feasible for the core and is its unique
projection by strict convexity.

At $H$, the marginal vector is $(1/2,1,1/2,0)$, not the core allocation.
The coalition $\{0,1,2\}$ has value $3/2$, so efficiency and singleton
coverage force $x_3=0$ in every core allocation.  The coalitions
$\{0,3\}$ and $\{1,3\}$ then give $x_0\ge1/2$ and $x_1\ge1$;
efficiency and nonnegativity force $x=(1/2,1,0,0)$.
This vector covers every coalition because the sum of its present $A$-side
capacities bounds production from above.  It is therefore the singleton high
core, proving the payment drop.

For an explicit production representation, activity $y_{ij}$ consumes one
unit of each resource $i\in A,j\in B$ and earns one unit of profit.
The four constraints are
\[
 y_{02}+y_{03}\le c_0,\quad y_{12}+y_{13}\le c_1,\quad
 y_{02}+y_{12}\le c_2,\quad y_{03}+y_{13}\le c_3.
\]
The maximum total activity equals the minimum of the two pool totals:
each total is an upper bound, attained by a bipartite transportation flow.
Giving nonmembers capacity zero realizes exactly the characteristic game.

At both profiles the unique optimal nonnegative resource-price vector is
$(1,1,0,0)$.  Indeed, let $a=\min(p_0,p_1)$ and $b=\min(p_2,p_3)$.
Dual feasibility $p_i+p_j\ge1$ implies $a+b\ge1$.
Writing the pool totals as $A_0=3/2$ and $B_0>A_0$, the objective is at least
$A_0a+B_0b\ge A_0+(B_0-A_0)b$.
Equality at the optimum forces $b=0,a=1$ and, because all capacities are
positive, forces both prices within each pool to equal their minima.
Thus dual-price-generated owner payments are $(1/2,1,0,0)$ even at $L$.
This is strictly smaller than the low full core used by the diagnostic.
\end{proof}

\begin{proposition}[Strictly concave utility does not remove the failure]
\label{net:concave-projection}
Keep the same four owners, activities, resource constraints, and reports,
but maximize $\sum_{i\in A,j\in B}u(y_{ij})$ with
$u(t)=10t/(10+t)$.  Projection of the corresponding marginal vector onto
the full efficient core still violates own-report monotonicity: owner 2's
payment falls by at least $1270/9471$ between $L$ and $H$.
This nonlinear production value is not asserted to be homogeneous.
\end{proposition}
\begin{proof}
Here $u(0)=0$, $u'(t)=100/(10+t)^2>0$, and
$u''(t)=-200/(10+t)^3<0$ on the nonnegative half-line.
Set
\[
 a=2u(1/4)=20/41,\qquad b=2u(1/2)=20/21,\qquad
 T=a+b=1240/861.
\]
At both profiles split each $A$ row equally between the two $B$ columns;
each column then uses $3/4$.  The split is feasible and attains the upper
bound from row-wise concavity, so it is optimal.
At $L$, deleting owners $0,1,2,3$ gives values $(b,a,b,b)$.
For example, with only one $B$ column of capacity one, two rows split its
capacity equally, which is feasible and optimal by concavity.
Thus $U(L)=(a,b,a,a)$ and its efficiency-hyperplane projection is
\[
 x^L=(a/2,b-a/2,a/2,a/2).
\]
It is nonnegative.  Mixed pairs involving owner 0 have value $u(1/2)\le a$
and payment $a$; those involving owner 1 have value $u(1)\le b$ and payment
$b$.  Triples $\{0,1,j\}$ have value $b$ and payment $b+a/2$;
$\{0,2,3\}$ has value $a$ and payment $3a/2$;
$\{1,2,3\}$ has value $b$ and payment $b+a/2$.
Other proper coalitions have value zero.  Efficiency holds, so $x^L$ is the
unique full-core projection.

At $H$, coalition $\{0,1,2\}$ has value $u(1/2)+u(1)$, because owner 2
can accommodate both full rows.  Put
\[
 \delta=T-u(1/2)-u(1)=\frac{520}{9471}>0.
\]
For every high core allocation, $x_3\le\delta$.
The coalitions $\{0,3\}$ and $\{1,3\}$ give
$x_0+x_3\ge u(1/2)$ and $x_1+x_3\ge u(1)$.
Summing and using efficiency implies
$x_2\le\delta+x_3\le2\delta=1040/9471$.
The high core is nonempty: $(a,b,0,0)$ covers every coalition, since its
present $A$ rows contribute at most their row-concavity bounds $a,b$.
The unique high projection therefore exists and satisfies the same bound,
regardless of its other coordinates.  Owner 2's low payment is
$a/2=10/41$, so its decrease is at least
$10/41-2\delta=1270/9471>0$.
\end{proof}

\paragraph{Original-source comparison (14 October 2026).}
The freely available author manuscript of Amigo, Belzarena and Vaton~\cite{ABV}
uses a fixed nonnegative service-resource matrix, scalar owner stocks and
strictly concave service utilities (Sections 2.1--2.2, pp.\ 4--5).
Section 3.2 (pp.\ 7--8) defines the same own-resource monotonicity, and
Sections 5.1--5.2 (pp.\ 18--20) project onto the full efficient core.
The strictly concave example above satisfies this inspected general model.
Table 5 and the discussion on pp.\ 26--27 report simulation evidence for
projection monotonicity. Our example rules out a universal extension of that
evidence; it does not refute a proved theorem. The publisher-formatted final
version has not been compared in full. The linear example alone would not
meet strict concavity, and neither projection diagnostic establishes a claim
about an unspecified restricted allocation set. The nonlinear diagnostic is
outside the homogeneous main theorem.

% Additional Oct7 structural results and boundaries; not separate paper claims.
\section{Structural facts for effective-capacity cut payments}\label{noise:structure}
Throughout these appendices, $G$ is a finite directed acyclic $s$--$t$
multigraph with one fixed owner per edge, and every edge lies on an $s$--$t$
path. Let $F(c)$ be maxflow, $T=F(c)$, $D_e(c)$ the largest feasible flow
through edge $e$, $\mathcal C$ the structural inclusion-minimal cuts,
$K=|\mathcal C|$, and $L$ the longest path length. With independent uniforms
$U_e$ and a fixed report-independent cut priority, let $W$ minimize
$\sum_{e\in C}D_eU_e$ and put
\[
 z_e=D_e\Prob(e\in W),\qquad y_e(c)=\sup_{0\le t\le c_e}z_e(t,c_{-e}).
\]
Thus $y$ is the least own-report-monotone pointwise majorant of $z$.
For other noise laws the superscript specifies the law. Set $U(1)=1$ and,
for $K\ge2$, $U(K)=a_K+1/\log a_K$, where
$a_K\log(a_K/e)=\log K$ and $a_K>e$.
The usual uniform-noise tail calculation gives $z(E)\le U(K)T$:
normalized weights lie in $[0,1]$, some cut has total one, and
$\Prob(\sum a_eU_e\le1)\le(e/\rho)^\rho$ when $\sum a_e=\rho$.
These appendices distinguish bounds on a particular rule from bounds on
all eligible mechanisms. The latter sharp theorem is Theorem~\ref{prod:upper}.

\begin{lemma}[Sensitivity and a simultaneous flow]\label{noise:sensitivity}
Writing $c^t=(t,c_{-e})$ and $\delta_{eh}=D_h(c)-D_h(c^0)$, one has
\[
 D_h(c^t)=\min\{D_h(c),t+D_h(c^0)\},\quad
 0\le\delta_{eh}\le D_e(c).
\]
For every path antichain $A$, $\sum_{e\in A}\delta_{eh}\le D_h(c)$;
for all edges the analogous bound is $\sum_e\delta_{eh}\le LD_h(c)$.
Furthermore, if
$r_e=2\E[F(DU)-F((DU)_{-e},0)]$, then
$z\le r\le2z$, and a single feasible flow $\bar f$ at capacities $D$
satisfies $r\le\bar f$. Consequently $r(A)\le T$ and
$r(E)\le\min\{L,2U(K)\}T$.
\end{lemma}
\begin{proof}
On a DAG, the maximum possible flow through $h=(u,v)$ is the minimum of
$c_h$, the $s$--$u$ maxflow and the $v$--$t$ maxflow, omitting trivial
terminal terms. The prefix and suffix regions are disjoint by acyclicity,
so their attaining flows concatenate through $h$. Each relevant cut
contains the changing edge at most once. Its minimum therefore has form
$\min\{A+t,B\}$, which gives the displayed hinge on $[0,c_e]$;
for $h=e$ it is $\min\{t,D_e(c)\}$.
Choose one flow attaining $D_h$ and retain only its paths through $h$.
Deleting paths through $e$ proves $\delta_{eh}\le f_e$ simultaneously
for every $e$. Summing over an antichain or all edges gives the bounds.

At capacities $X=DU$, choose one maximum flow $f(X)$. Removing its paths
through $e$ proves $F(X)-F(X_{-e},0)\le f_e(X)$ simultaneously.
Conditionally on $U_{-e}$, the gain from setting the focal multiplier to
$u$ is a nonnegative increasing concave function $g(u)$ with $g(0)=0$.
Its endpoint is $g(1)=\min\{D_e,(B_e-A_e)_+\}$, where $A_e$ is the
cheapest containing-cut intercept and $B_e$ the cheapest excluding-cut
score. Conditional integration gives $\E g(1)=z_e$ and
$g(1)/2\le\int_0^1g(u)du\le g(1)$. Thus $z\le r\le2z$.
The flow $\bar f=2\E f(DU)$ conserves flow and obeys
$\bar f_e\le2D_e\E U_e=D_e$. A measurable selection is obtained from
the finitely many affine LP bases by lexicographic tie-breaking.
Path decomposition now gives the remaining claims.
\end{proof}
The hinge also gives a valid, though noncollective, gain estimate. Put
\[
 \Delta_e(t,U)=\max_{C\ni e}\sum_{h\in C\setminus\{e\}}
                         (\delta_{eh}-t)_+U_h.
\]
The excluding score cannot increase under an underreport, and the containing
intercept decreases by at most $\Delta_e$. The clamp is increasing and
one-Lipschitz, whence
\[
 0\le y_e-z_e\le\sup_{0\le t\le D_e}\E\min\{t,\Delta_e(t,U)\}.
\]
Reports above $D_e$ preserve the feasible-flow set. If no excluding cut
exists, the excess is zero. An expectation of a maximum cannot be replaced
by a maximum of expectations; no sharp aggregate gain bound follows here.

\begin{proposition}[Antichains and a conditional reduction]\label{noise:antichain}
For incomparable edges $e,f$, $y_e+y_f\le R(\{e,f\})$, where $R(A)$ is
the maximum flow on paths meeting the antichain $A$. For a complete parallel
bundle $X$ from $u$ to $v$, with total capacity $q$ and external prefix/suffix
bottleneck $B$, one has $y(X)\le\min\{q,B\}$.
If an inclusion-minimal antichain $A$ violates $y(A)\le R(A)$, then
$|A|\ge3$ and the simultaneous profile setting $c_e=y_e(c)$ for $e\in A$
preserves $R(A)$. At that profile every cut of $G_A$ meeting $A$ costs
strictly more than $R(A)$, and $G_A$ has a common unchanged minimum cut
disjoint from $A$. Here $G_A$ is the union of paths meeting $A$; this cut
need not disconnect the full network.
This reduction does not assert existence or impossibility of such an $A$.
\end{proposition}
\begin{proof}
The union $G_A$ of paths meeting $A$ has no path avoiding $A$: vertices before
and after $A$ cannot overlap, and an edge crossing between those regions
must belong to $A$, by acyclicity and the antichain property.
Fix $e\in A$, put $R=R(A)$ and $R_-=R(A\setminus\{e\})$. Under an
own underreport, if the effective focal capacity is at most $R-R_-$ its
payment is at most that quantity. Otherwise every containing cut in $G_A$
has capacity greater than $R$, since its other edges block an $R_-$ flow.
A minimum cut $C_0$ of $G_A$ excludes $e$ and has capacity at most $R$.
It blocks every path through $e$, so the endpoint flow marginal defining
$z_e$ is at most its random capacity. Taking its mean gives $R/2$. Hence
\[
 y_e\le\min\{D_e,\max(R-R_-,R/2)\}.                 \tag{A}
\]
For two edges, $R_-=D_f$. If either $D_f\le R/2$, then
$y_e\le R-D_f$ and $y_f\le D_f$; if both effective capacities exceed
$R/2$, each payment is at most $R/2$.

For a parallel bundle, $D_i=\min(c_i,B)$. External feasible flows depend
on its reports only through $\min(q,B)$. If $q\le B$, use $y_i\le c_i$.
Otherwise write $W=\sum_iD_i$. Every minimal cut contains the whole bundle
or none, so $z_i=D_ip$ for a common $p$; the simultaneous flow lemma gives
$Wp\le B$. With $s_i=q-c_i$ and $\delta_i=(B-s_i)_+$, underreports
$t\ge\delta_i$ keep every other effective cap fixed, and conditional
clamping implies payment at most $z_i(c)$. Smaller reports pay at most
$\delta_i$. Thus $y_i\le\max(z_i,\delta_i)$. If $\delta_i>0$, then
$W=D_i+s_i$ and $D_i\ge B-s_i$, giving
$\delta_iW\le BD_i$. Summing $y_i\le BD_i/W$ proves the bundle claim.

For a minimal violating $A$, write $Y(B)=\sum_{e\in B}y_e$ and
$\delta_e=R(A)-R(A\setminus\{e\})$. Minimality gives $y_e>\delta_e$.
The own-report maxflow hinge in $G_A$ is
$\min\{R(A),R(A\setminus\{e\})+t\}$, so any witness with
$z_e(c^t)>\delta_e$ preserves $R(A)$. All minimum cuts in the rest of this argument are in
$G_A$. If an original minimum cut contains $e$, deleting $e$ lowers its
capacity from $R(A)$ to $R(A)-c_e$, so
$R(A\setminus\{e\})\le R(A)-c_e$. Removing the paths through $e$
from an original maximum flow gives the reverse inequality. Hence
$\delta_e=c_e$, contradicting $c_e\ge y_e>\delta_e$.
Thus every original minimum cut $C^*$ of $G_A$ is disjoint from all of $A$.
For a cut $C$ of $G_A$ let $B=C\cap A$. Its unchanged non-$A$ capacity is at least
$R(A\setminus B)$, by a flow meeting $A\setminus B$ and avoiding $B$.
At the joint $y$-profile its capacity, if $B\ne\varnothing$, is at least
$Y(B)+R(A\setminus B)\ge Y(A)>R(A)$. Cuts avoiding $A$ retain their
capacities and $C^*$ stays tight. It is saturated by a through-$A$ flow,
so its full-network effective capacities remain its original reports.
Near-maximizing separate witnesses suffice; no attainment is needed.
\end{proof}

\begin{proposition}[Cut height]\label{noise:height}
For $h(C)=\max_P|P\cap C|$ and $H=\sum_Ch(C)$,
\[
 h(C)\le\left\lfloor\frac{\sqrt{8K+17}-3}{2}\right\rfloor,
 \qquad H=O(K^{3/2}).
\]
The exponent for $H$ is attained in order by the classical Fibonacci DAG:
vertices $0,\ldots,2r-1$ and arcs of lengths one and two have
$K=r^2$ and $H=r(r+1)(2r+1)/6$.
\end{proposition}
\begin{proof}
For a minimal cut $C$, let $R$ be the vertices reachable after deleting it.
Minimality gives, for each $e\in C$, a private path meeting $C$ only in $e$;
its prefix lies in $R$ and suffix outside $R$, so $C=\delta^+(R)$.
Fix a path $P$ meeting $C$ in $e_i=(a_i,b_i)$, $i=1,\ldots,r$, and a
topological order $\tau$. For $i\le j$, the vertex set
\[
 S_{ij}=\{v:\tau(v)\le\tau(a_i)\}
       \cup\{v\in R:\tau(v)<\tau(b_j)\}
\]
has boundary trace on $P$ exactly $e_i,\ldots,e_j$. Each corresponding
private path meets that boundary only at its own edge, so every minimal
subcut contains the full interval. Distinct intervals give distinct cuts.
Every edge outside $C$ on $P$ also has a minimal cut with singleton $P$ trace,
obtained by a topological-prefix boundary. If $P$ has length $\ell$,
$K\ge\ell+r(r-1)/2$; successive exits from $R$ imply $\ell\ge2r-1$.
This proves the stated finite bound.

For the Fibonacci DAG, a canonical cut partition in vertex order has word
$S^a(TS)^bT^c$ with $a,c\ge1$: two consecutive $T$'s prevent a later
reachable $S$, while two consecutive $S$'s after the first $T$ prevent its
private suffix to the sink. Conversely such words make every $S$ reachable
within $S$ and every $T$ able to reach the sink within $T$, proving minimality.
The height is $h=b+1$, achieved by the backbone path. There are
$2r-2h+1$ such cuts, giving both sums. The graph and quadratic cut count
are established; no priority is asserted for them or the height comparison.
\end{proof}

\begin{proposition}[Serial blocks and failure of a maximum-module bound]
\label{noise:series}
If an edge is replaced by $k$ serial owners with no internal attachments,
their total envelope payment is at most their common effective capacity.
Nevertheless no constant multiple of the largest standalone module payment
bounds envelope payment after arbitrary series composition.
\end{proposition}
\begin{proof}
Let $q$ be the least report in the serial block. All its effective capacities
are a common $d(q)$, and external capacities depend only on $q$. Minimal cuts
use either no block edge or one arbitrary member. The smallest of its $k$
uniforms selects a uniformly distributed member, independently of that
minimum and all external uniforms. Thus every member has base payment
$Z(q)/k$, where $Z(q)=d(q)\Prob(W\text{ meets the block})$.
Each member's underreports realize exactly bottlenecks $q'\in[0,q]$;
therefore the sum of their envelopes is $\sup_{q'\le q}Z(q')\le d(q)$.
The collapsed noise is $\operatorname{Beta}(1,k)$, not uniform.

For the second claim, fix $m\ge2$, $r=4^m m!$, $M=16r$. Each of $r$
modules has two serial bundles: $M$ owners of capacity $1/M$ and $m$
owners of capacity one. Let their uniform scores be $L_i,H_i$.
Standalone payment is $1+(m-1)\Prob(H_i<L_i)\le1+(m-1)/m!\le2$.
In the composition, some $H_i\le1/4$ has probability at least $1/2$,
whereas $\Prob(\min_iL_i<1/4)\le r\,4/(3M)=1/12$.
A heavy cut therefore wins with probability at least $5/12$, giving payment
at least $5m/12$. The base equals its envelope at these profiles: underreporting
a light owner leaves all other owners of its containing cut unclipped, as
the resulting grand flow stays at least $1/M$; underreporting a heavy owner
leaves its pool total at least one. The containing intercept is fixed and
the other caps are nondecreasing, so conditional clamping proves monotonicity.
For general serial modules, $K=\sum_iK_i$ and global caps are
$D_{i,e}=\min\{d_{i,e},\min_jT_j\}$; this bookkeeping alone supplies no
maximum-module estimate.
\end{proof}

\section{What fixed noise and Bernoulli conditioning can achieve}\label{noise:laws}
\begin{theorem}[Conditioned-zero binary production]\label{noise:bernoulli}
Consider finitely many nonempty resource bundles $P$, each consuming one
unit of its resources and producing one unit of value. Each resource has
one owner. Let $V$ be the vertices of
$\{w\ge0:\sum_{e\in P}w_e\ge1\ \forall P\}$ and $K=|V|$.
Let $D_e$ be maximum feasible resource usage and $w^*$ the fixed-priority
minimizer of $\sum_eD_eB_ew_e$, where $\Prob(B_e=0)=q\in(0,1)$ independently.
Then
\[
 p_e=D_e\E[w_e^*\mid B_e=0]
\]
is own-report monotone, covers every coalition, and for every $\lambda>0$
obeys
\[
 p(E)\le\frac{T(\log K+\lambda+1)}
 {q[-\log(q+(1-q)e^{-\lambda})]}.
\]
For fixed $q=e^{-1}$ this is $(e+o(1))T\log K$. In the network subclass,
optimizing one common $q=q(K)$ cannot improve the leading constant $e$.
This class bound is superseded as an unrestricted existence bound by
Theorem~\ref{prod:upper}, but its formula and class limitation are distinct.
\end{theorem}
\begin{proof}
LP duality gives $F(c)=\min_{w\in V}c\cdot w$ and $0\le w_e\le1$:
a vertex coordinate exceeding one could be perturbed both ways without
violating an active constraint. Nonempty activities bound every primal
variable. Effective usage satisfies $D_e\le\min(c_e,T)$ and preserves
the entire feasible polytope, also after any binary mask. Convex mixtures
of feasible plans at own reports zero and $u$ imply, for $0<t\le u$,
\[
 D_h(t)\ge D_h(0)+(t/u)[D_h(u)-D_h(0)].
\]
Thus $D_h(t)>D_h(0)$ is independent of the positive own report. With $B_e=0$
the minimum masked value is independent of that report. A vertex remains
optimal precisely when it was optimal at zero and every coordinate with
$B_hw_h>0$ has no effective-usage increase. Hence the whole minimizing
family and its fixed-priority selection are invariant on positive own
reports. Multiplication by nondecreasing $D_e$ proves incentives, including
zero where the payment vanishes.
If $w_0,w_1$ are selected with $B_e=0,1$, their two optimality inequalities
give $w_{0e}\ge w_{1e}$ when $D_e>0$. Therefore for
$z_e^q=D_e\E w_e^*$, one has $z_e^q\le p_e\le z_e^q/q$.
For every dual vertex $w$ and coalition $S$, binary-mask preservation gives
$\sum_{e\in S}D_ew_e\ge F(D\mathbf1_S)=F(c\mathbf1_S)$.
Thus every selected dual allocation covers the original coalition values,
so the same is true of $z^q$ and $p$.

For $a_e=D_ew_e/T\in[0,1]$ and $\rho_w=\sum a_e$, convexity gives
$\log(q+(1-q)e^{-\lambda a_e})\le-\alpha a_e$, where
$\alpha=-\log(q+(1-q)e^{-\lambda})$. Some unmasked objective equals $T$,
so the winning masked objective is at most $T$. Chernoff and a union over
vertices yield $\Prob(\rho_{w^*}>s)\le\min(1,Ke^{\lambda-\alpha s})$.
Integrating gives $\E\rho_{w^*}\le(\log K+\lambda+1)/\alpha$.
Take $q=e^{-1}$ and $\lambda=\log(2+\log K)$ for the asymptotic bound.

In graphs, $V$ consists exactly of minimal-cut vectors: maxflow/mincut
identifies the cut dominant with the path-cover polyhedron by separation
with nonnegative objectives; each minimal cut is uniquely exposed by
cost one on itself and a sufficiently large cost outside. With directed
cycles, define usage by path packings, excluding circulations. The same
binary theorem also covers unit-value splittable multicommodity paths,
where dual vertices can be fractional and $K$ is not a cut count.

For the sharp constant, the network identity
\[
 p^q(E)=q^{-1}\E\sum_{e\in W}D_e(1-B_e)                 \tag{B}
\]
is immediate by conditioning. With $K$ serial two-owner pools, each report
$1/2$, $T=1$ and
$p^q(E)=[1-\tfrac12(1-q^2)^K-\tfrac12(1-q)^{2K}]/q
\ge K/(1+2Kq)$, by $(1-q)^{-2K}\ge1+2Kq$.
Put $N=K-1$, $A=\log N\ge16$. If $q\ge A^{-2}$, write
$h=\log(1/q)$, $m=\lfloor(A-\sqrt A)/h\rfloor$,
$R=\lceil\sqrt A/h\rceil$. Use $N$ serial pools of $m$ unit owners
and a final pool of $R$ owners of report $1/R$.
When the final score is positive and some heavy pool is all closed, the
winning cut has closed weight $m$. By (B),
\[
 p^q(E)\ge\frac m q(1-q^R)[1-(1-q^m)^N]
 \ge e(A-\sqrt A-2\log A)(1-e^{-\sqrt A})(1-e^{-e^{\sqrt A}}),
\]
since $Nq^m\ge e^{\sqrt A}$, $q^R\le e^{-\sqrt A}$ and
$q\log(1/q)\le1/e$. If $q<A^{-2}$ the two-owner family gives
$p^q(E)\ge A^2/3\ge eA$. The lower is uniform in $q$, proving class
optimality on multigraphs, without claiming its simple-DAG conversion.
\end{proof}

\begin{theorem}[Fixed iid-law classification]\label{noise:classification}
Fix a law $X$ on $[0,1]$, independent of graph size, and replace uniforms
in $z,y$ by iid $X$. On DAGs the worst envelope factor has these properties:
identically zero noise has no owner-independent bound already at $K=2$;
constant positive noise has exact worst factor $K$ on multigraphs;
a nondegenerate law on $\{0,c\}$ has order $\Theta_q(1+\log K)$;
and a nonconstant conditional positive law has a lower bound
$\Omega_X(K^{\eta_X}/\log K)$ for some $\eta_X>0$, even on simple DAGs.
Every nonzero law has an $O_X(K)$ upper bound. These statements concern
pointwise-majorant repairs of that fixed base, not arbitrary mechanisms.
\end{theorem}
\begin{proof}
Let $\mu=\E X>0$. For a cut $C$, set $w_h=D_h/T\le1$,
$d=\sum_Cw_h$ and $S=\sum_Cw_hX_h$. At a report $t=xD_e$, $0<x\le1$,
homothetic scaling of feasible flows gives $D_h(c^t)\ge xD_h(c)$.
Selection of $C$ implies $xS\le1$. For $u=1/x$, if $u\ge\mu d/2$,
$(d/u)\Prob(S\le u)\le2/\mu$; otherwise Hoeffding gives a bound
$d e^{-\mu^2d/2}\le2/(e\mu^2)$. Thus
\[
 t\Prob(W(c^t)=C)\le A_\mu T\frac{D_e}{D(C)},\qquad
 A_\mu=\max\{2/\mu,2/(e\mu^2)\}.
\]
Summing over containing cuts before taking suprema proves $y(E)\le A_\mu KT$.
Constant positive noise makes the same argument deterministic with constant
one. Equality is approached with $K$ serial two-owner pools of reports
$1/2$: lowering either owner makes its pool uniquely cheapest. On simple
DAGs $N$ serial diamonds give $K=4N$ and factor at least $N$, by lowering
an owner of each module's first-priority local cut.
Zero noise always selects the fixed first cut: two serial pools of $m$
owners, with reports one in that cut and $1/m$ in the other, give factor
$m$ at $K=2$. Subdivision into two-edge private paths gives a simple-DAG
lower of order $\log K$.
For $X\in\{0,c\}$, Theorem~\ref{noise:bernoulli} majorizes its envelope
and proves the logarithmic upper. For the lower, take
$N=\lceil q^{-m}\rceil$ pools of $m$ unit owners and a final unit singleton.
When the singleton is open and some heavy pool is closed, the winning
uncensored cut weight is $m$, so $z(E)\ge(1-q)(1-e^{-1})m$.
Subdivision of heavy edges replaces $q$ by $q'=1-(1-q)^2$ and gives
$K=N2^m+1$, still $m=\Theta_q(\log K)$ on simple DAGs.

It remains to handle a nonconstant conditional positive law. Let
$s=\Prob(X>0)$, $\nu=\E\min(X,X')$ and $d=\mu-\nu/s>0$;
strictness follows by conditioning on both draws being positive. Put
$a=(1/s+\mu/\nu)/2$, $b=d/(4\mu)$, $\delta=d/40$.
Then
\[
 \E\min(aX,(a+b)X')\le a\nu+b\mu=\mu-10\delta,
 \qquad \E\min(aX,mX')\uparrow a\mu s>\mu+20\delta.
\]
Choose $m_0$ making the second expectation at least $\mu+10\delta$;
integers $B,C,M$ satisfying $C\ge(\mu+10\delta)/(b\mu)$, $M\ge\delta^{-2}$; and
\[
 \gamma=\min\{32\delta^2/a^2,32\delta^2/(Cb^2)\},\qquad
 B\ge\max\{1,m_0,a+b,b/\delta,2/\gamma\}.
\]
For integer $r\ge4$ use $N=2^r$ modules in series, $m=Br$,
$k=Cm$, and a final $M$-owner pool of total one. In each module put
$e:s\to A$ of report one, $y_i:s\to B_i$ of report $a/m$,
$g_i:A\to B_i,h_i:B_i\to t$ of report two, and $k$ parallel
$Z:A\to t$ owners of report two. Under a separate focal report $q$,
$T=1$ and
$D_e=D_g=D_Z=q$, $D_y=a/m$, $D_h=\min(1,q+a/m)$.
These values follow from explicit paths; each raw module has flow $a+q>1$.
There are $2^{m+1}$ cuts per module, according as $e$ is included or
excluded, so $K=N2^{m+1}+1$.

Lower one focal report to $t=b/m$. Its cheapest containing score is
$tX_e+A_{\rm low}$, where
$A_q=\sum_i\min((a/m)X_{y_i},D_{h_i}X_{h_i})$.
The high and low means are at least $\mu+10\delta$ and at most
$\mu-10\delta$. Every excluding score is at least $q\sum_ZX_Z$,
whose low mean is at least $\mu+10\delta$.
The squared-range sums are $a^2/m$ and $Cb^2/m$. Hoeffding bounds each
wrong-side deviation to $\mu\pm6\delta$ by $e^{-\gamma m}$.
The final score is an average of $M$ draws and lies within $2\delta$
of $\mu$ except with probability $1/16$. A union over at most $2N$
other failures costs at most $1/8$. Since $t\le\delta$, every winning
cut contains the lowered focal owner on this event. Hence at the single
original profile $y(E)\ge Nb/(2m)$, proving a positive polynomial exponent.

For a simple DAG, choose the constants afresh for $V=\min(X_1,X_2)$,
whose conditional positive law is still nonconstant. Subdivide each
$y,h,Z$ and final edge into a private two-edge path, retaining $e,g$.
The primitive owner noise is still $X$, while those path scores use $V$.
The proof used only $X_e\le1$ and nonnegativity of $X_g$ for the remaining
edges. Now
$K=2^{(1+2B)r}+2^{[1+(C+2)B]r}+2^M$ and the same payment lower holds,
giving exponent $1/[1+(C+2)B]>0$. Strict cut-class separation makes all
arguments valid with atoms and any fixed tie priority.
\end{proof}

The classification fixes the law as $K$ grows. One particular variable-law
escape also fails: if $q=K^{-1/2}$ and
$X\sim(1-q)\delta_1+q\operatorname{Unif}(0,\epsilon)$ with
$0<\epsilon\le1/8$, use $K$ serial pools of reports $(3/4,1/4)$.
Lower each small owner separately to $1/8$. Its pool wins strictly when
the companion has low noise, the focal noise is one, and no other pool has
two low noises. Its score is at most $7/32$, while competitors cost at least
$1/4$. Effective capacities remain the reports. Summing gives
\[
 y(E)\ge(K/8)q(1-q)(1-q^2)^{K-1}\ge\sqrt K/64\quad(K\ge4).
\]
This does not classify arbitrary $K$-dependent noise laws.

\begin{proposition}[Uniform mixture of Bernoulli parameters]\label{noise:uniform-mixture}
The truthful covering rule $\bar p=\int_0^1p^q\,dq$ has worst-case
budget order $\Theta(K)$ on multigraphs.
\end{proposition}
\begin{proof}
Under a uniform random arrival order, before owner $e$ arrives select the
first cut minimizing the revealed $D$-weight and charge $D_e$ if it contains
$e$. Conditioning on its arrival time gives exactly $\bar p$.
At the last positive charge assigned to a particular cut, its revealed score
was at most $T$ before arrival and increases by at most $D_e\le T$.
All charges previously assigned to that cut are included in its revealed
score, even if the cut ceased to win and returned. Thus each cut is charged
at most $2T$, proving $\bar p(E)\le2KT$.
For the lower, take $K$ serial pools of $B^{K-1}$ owners, with $B\ge3$.
Enumerate pools in their fixed priority order. In pool $i$ give
$m_i=B^{K-i}$ owners report one and all others zero. Then $T=1$ and $D=c$.
Let $\tau_i$ be its first positive-owner arrival. If
$\tau_i>\max_{j<i}\tau_j$, cut $i$ is selected just before that arrival
and charges one. Since
$\Prob(\tau_i<\tau_j)=m_i/(m_i+m_j)\le B^{-(i-j)}$ for $j<i$,
\[
 \bar p(E)\ge K-\sum_i\sum_{d=1}^{i-1}B^{-d}
             \ge K(1-1/(B-1)).
\]
Zero owners remain part of the fixed topology. This construction has
$n=KB^{K-1}$ owners and does not refute an $O(\log n)$ bound.
\end{proof}

\section{Truthful repairs that lose their budget or coalition guarantee}
\label{noise:repairs}
\begin{proposition}[Deletion repairs can have linear cost]\label{noise:deletion}
Both of the following rules are truthful pointwise majorants of $z$, and
hence cover coalitions and majorize $y$:
\[
 p_e^{\rm pot}=2[\Psi(c)-\Psi(c^0)],\quad
 \Psi(c)=\E F(D(c)U),
 \qquad
 p_e^{\rm one}=\E\min\{D_e(c),(B_e(c)-A_e(c^0))_+\}.
\]
Here $A_e$ omits the focal term from containing cuts, and $B_e$ minimizes
over excluding cuts. If no excluding cut exists, set $p_e^{\rm one}=D_e$.
Each rule nevertheless has examples with total payment $\Omega(KT)$.
\end{proposition}
\begin{proof}
$\Psi$ is coordinatewise increasing, so its own deletion difference is
own-monotone. At fixed other uniforms, concavity of the focal flow gain
implies its average is at least half its endpoint. Since recomputing $D$
after deletion only decreases other capacities, $p^{\rm pot}\ge z$.
For the one-sided rule, $A_e(c^0)$ is independent of $c_e$, whereas $D_e$
and $B_e(c)$ are nondecreasing. Also $A_e(c^0)\le A_e(c)$, proving domination.

For the potential example, put a branch of $k$ unit serial owners in parallel
with a single owner of capacity $\beta=1/2$, followed by a final unit edge.
Then $T=1,K=k+1$. With $A_k=\min_{i\le k}U_i$,
$\Psi=\E\min(A_k+\beta U,V)$. Deleting any serial owner leaves effective
capacities $\beta$ on the other branch and final edge, so its deletion
potential is $\beta/3$. Meanwhile
$\Psi\ge\E\min(\beta U,V)=\beta/2-\beta^2/6$.
Each serial owner receives at least $\beta(1-\beta)/3=1/12$.
Indeed $A_k\to0$ and deletion of the parallel owner leaves potential
$1/(k+2)$, yielding
$p^{\rm pot}(E)=2[(k+2)\Psi-k/6-1/(k+2)]\sim k/12$.
The original envelope on this TTSP graph equals $z$ and has total at most
$3/2$, the largest effective weight of a selected cut.

For the one-sided example, put $k$ unit serial owners $X$ on $s\to A$,
add $b:s\to B$ of capacity $\beta\in(0,1)$, and unit edges
$a:A\to t,g:A\to B,h:B\to t$. Then $T=1+\beta$, $K=2k+2$; cuts are
$\{x_i,b\},\{x_i,h\},\{a,h\},\{a,b,g\}$.
Deleting $x_i$ leaves only $b+h$, so
$A_i(c^0)=\beta\min(U_b,U_h)$, while
$B_i(c)\ge m=\min(\beta U_b,U_h)$. Therefore
\[
 p_i^{\rm one}\ge\E[m-\beta\min(U_b,U_h)]
 =\beta/2-\beta^2/6-\beta/3=\beta(1-\beta)/6.
\]
For $k\ge2$, $B_i\le m+\min_{j\ne i}U_j$ and the expected error is
$1/k$. The four other payments sum to at most $3+\beta$, so
$p^{\rm one}(E)=k\beta(1-\beta)/6+O(1)$.
At $\beta=1/2$ the normalized asymptotic is $K/72$.
By contrast, when $x_i$ underreports, all serial effective caps become equal;
its selection requires the smallest of their uniforms. Hence $y_i\le1/k$
and $y(E)\le4+\beta$ on this family.
\end{proof}

\begin{proposition}[Common scales]\label{noise:scales}
Common clipping $a_e=\min(c_e,T)$ preserves coalition coverage and the
$U(K)T$ base budget but its uniform-cut payment need not be own-monotone.
The truncated scale envelope
\[
 P_e=\sup_{0\le r\le T}\min(c_e,r)
       \Prob\bigl(e\in W(\min(c,r)U)\bigr)
\]
is truthful and covering, but can cost at least $5KT/8$ with $O(K^7)$ owners.
\end{proposition}
\begin{proof}
Every feasible DAG flow has each coordinate at most its total, at most $T$,
so common clipping preserves all feasible flows. A minimum cut retains
weight $T$, proving the usual tail-budget premise.
For nonmonotonicity, use $x:s\to A$ of capacity one, $z:A\to t$ of
capacity three, $g:A\to B$ of capacity $\delta=1/10$, focal $h:B\to t$,
and $N$ parallel $s\to B$ owners of report $1/N$.
The four cuts are $\{x\}\cup Y$, $\{x,h\}$, $\{z,h\}$,
$\{z,g\}\cup Y$. For $h\in[3/4,1]$, $T=1+h$.
The focal payment is
\[
 w_h(h)=\E\min(h,Y+W_T),\qquad
 W_T=\min(\delta U_g,(U_x-TU_z)_+),\quad
 Y=N^{-1}\sum U_i.
\]
Direct tail integration gives
$\Prob(W_T>r)=(1-r/\delta)(1-r)^2/(2T)$ on $[0,\delta]$ and
$\E W_T=187/(8000T)$. At $h=3/4$, clamping loses at most
$\Prob(Y>13/20)\le1/2700$ when $N=10000$, by variance $1/(12N)$.
Ignoring the high-profile clamp only raises its payment. Thus
\[
 w_h(3/4)-w_h(1)\ge187/112000-1/2700=3929/3024000>0.
\]
A smaller exact certificate uses $N=6$. With
$F_N(u)=N^N\sum_{j=0}^N(-1)^j\binom Nj(u-j/N)_+^N/N!$ and
$L_N(u)=N^N\sum_{j=0}^N(-1)^j\binom Nj(u-j/N)_+^{N+1}/(N+1)!$,
the exact expression is
$w_h=1/2-L_N(1-h)+\int_0^\delta
(1-r/\delta)(1-r)^2F_N(h-r)/(2T)\,dr$.
Piecewise polynomial integration gives the positive difference
$587847649425223/914457600000000000$.

For a fixed external $r$, conditional clamping proves own-monotonicity;
the interval $[0,T]$ expands with the own report. The supremum is therefore
truthful and includes the covering rule at $r=T$.
For its lower bound, use $K$ serial bundles. Put $M=64K^4$, $N=256K^6$,
$b_j=3/4+j^2/(4K^2)$. Bundle $j$ has $N$ owners of report $b_j/N$ and
$n_j=M(K-j)$ owners of report one. Then $T=1$.
At scale $r_i=i/(2K^2M)$ all light owners are unclipped, and the clipped
bundle totals $g_j=b_j+n_jr_i$ satisfy
\[
 g_j-g_i=(j-i)^2/(4K^2),\qquad g_i=1-(K-i)^2/(4K^2)\ge3/4.
\]
Each random score has variance at most $1/(12N)+1/(48KM)$.
Chebyshev and a union bound with $\eta=1/(16K^2)$ give failure probability
at most $1/(12K)+1/12\le1/6$ for all deviations being less than $\eta$.
The mean gaps are at least $2\eta$, so bundle $i$ wins with probability
at least $5/6$. Summing its scale-$r_i$ payments over disjoint bundles
proves $P(E)\ge5K/8$.
The same obstruction can be made balanced: replace heavy reports in
$j<K$ by $(K+j)/(4K^2M)$. Every raw pool then totals one; for $i<K$
the clipped totals become $\min(1,g_j)$, preserve the stated gap, and have
no larger variance. This gives $P(E)\ge5(K-1)/8$.
\end{proof}

\begin{proposition}[Personal ceilings lose coverage]\label{noise:ceilings}
Let $B_e=\min_{C\not\ni e}c(C)$, allowing $\infty$, and use for owner
$e$ the clipped vector $a_h^{(e)}=\min(c_h,B_e)$ and payment
$p_e=a_e^{(e)}\Prob(e\in W(a^{(e)}U))$.
This is truthful and spends at most $(U(K)+1)T$, but no universal constant
multiple covers every coalition, even on TTSP graphs.
\end{proposition}
\begin{proof}
$B_e$ is independent of $c_e$, proving incentives by conditional clamping.
If $B_e=T$, its coordinate equals the common-clipped payment. If $B_e>T$,
every minimum cut contains $e$, so all such exceptional edges belong to one
minimum cut and their payments sum to at most $T$. The rest sum to at most
$U(K)T$. At $T=0$ the corresponding endpoint marginals vanish.

For a counterfamily, fix $q\ge4$, $m=q^2$, $\delta=q^{-2}$,
$N=m!q^m$. The initial component has parallel branches: an edge $e$ of
capacity $\delta$, and a serial path with first capacity $1-\delta$ and
$m$ further owners of report $m$. Follow it by $N$ serial bundles of $m$
unit owners each. Cuts are $\{e,a\}$, $\{e,h_i\}$ and those bundles;
$T=1,K=N+m+1$. Personal ceilings are $B_e=B_a=m$ and all others one.
Take $S$ to contain $e$ and one owner from every bundle, so $v(S)=\delta$.
With $Y=\min_{i\le m}U_{h_i}$,
$Z=\min_{j\le N}\sum_{k\le m}U_{jk}$ and
$X_r=\delta U_e+\min((1-\delta)U_a,rY)$, its payment is exactly
\[
 p(S)=\delta\Prob(X_m<Z)+\Prob(Z<X_1).
\]
The simplex CDF $z^m/m!$ for $z\le1$ gives
$\Prob(Z<1/(2q))\le2^{-m}$ and
$\Prob(Z>2/q)\le(1+2^m)^{-1}$.
On $X_m<Z\le2/q$, either $(1-\delta)U_a<2/q$ or $mY<2/q$.
Also $X_1\le\delta+Y$. Hence
\[
 \frac{p(S)}{v(S)}\le
 \frac{2q}{q^2-1}+\frac2q+\frac1{1+2^{q^2}}
 +q^2\left[2^{-q^2}+
  \left(1-\frac{q-2}{2q^2}\right)^{q^2}\right]
 =\frac{4+o(1)}q\longrightarrow0.
\]
Indeed the last power is at most $e^{1-q/2}$.
Stirling gives $\log K/\log\log K\sim3q^2/2$, so the necessary scalar
coverage repair is at least order $\sqrt{\log K/\log\log K}$.
This is a limitation of owner-dependent normalization, not a lower bound
on optimal mechanisms.
\end{proof}

\subsection{Failure of coordinatewise and cross-owner comparisons}
\begin{proposition}\label{noise:pointwise}
Neither $y_e\le Cz_e$ nor $y_e\le Cr_e$ holds for a universal constant
$C$. This already fails for DAGs of longest path four and cut height at
most two. The frozen marginal $r$ itself need not be own-monotone.
\end{proposition}
\begin{proof}
Use vertices $s,A,B_1,\ldots,B_m,q,t$: focal $x:s\to A$;
$y_i:s\to B_i$ of report $\epsilon=1/10$;
$g_i:A\to B_i,h_i:B_i\to q$ of report two;
$k$ parallel $Z:A\to q$ owners of report two;
$N$ direct $P:s\to q$ owners of report $(1-2b)/N$;
and $N$ final $R:q\to t$ owners of report $1/N$.
Whenever $m\epsilon\ge2b$ and $x\in[0,9/10]$, $T=1$ and
$D_x=x,D_y=\epsilon,D_g=D_Z=x,D_h=x+\epsilon$, with $D_P,D_R$ unchanged.
Cut families are $R$; $P\cup\{x\}$ plus one of $y_i,h_i$ per branch;
and $P\cup Z$ plus one of $h_i,\{y_i,g_i\}$ per branch.
Thus $K=2^{m+1}+1$, and private paths verify every listed cut is minimal.

Let $B=R-P$ now denote random score sums, so $\E B=b$ and
$\operatorname{Var}B\le1/(6N)$. Put
$A_i=\min(\epsilon U_i,(x+\epsilon)V_i)$ and
$J_x(s)=\E\min(x,(s-\sum_iA_i)_+)$.
If the random $Z$ score is at least one, the only competing excluding
threshold is $R$. Conditional integration and Lipschitzness therefore give
\[
 |z_x-J_x(b)|\le1/\sqrt{6N}+10^k/k!.
\]
For $b\le\min(\epsilon,x)$ the density of each $A_i$ on the relevant
simplex is $\alpha_x-\beta_xa$, with
$\alpha_x=10+1/(x+\epsilon)$, $\beta_x=20/(x+\epsilon)$. Thus
\[
 J_x(b)=\sum_{j=0}^m\binom mj\alpha_x^{m-j}(-\beta_x)^j
                      \frac{b^{m+j+1}}{(m+j+1)!}.       \tag{C}
\]
Set $b=1/(100m)$, $I=b^{m+1}/(m+1)!$, $Q=11^mI$,
$N=\lceil16/Q^2\rceil$, $k=2\lceil8/Q\rceil$.
At $x_0=1/10$, density is at least $15-1/m$, so
$J_{x_0}\ge(14/15)15^mI$; at $x_1=9/10$ it gives
$(54/55)Q\le J_{x_1}\le Q$.
The error is at most $Q/4$: the variance term is at most $Q/8$, and
$10^k/k!\le(200/k)^{k/2}\le2^{-k}\le1/k\le Q/16$ since $k\ge800$.
Consequently $z_{x_1}>0$ and
\[
 \frac{y_x(c^{x_1})}{z_x(c^{x_1})}
 \ge\frac{56}{75}(15/11)^m-\frac15\longrightarrow\infty.
\]
The comparison with $r$ fails because $r\le2z$.
An exact smaller factor-two certificate uses $m=3,b=1/20,N=100000,k=50$;
(C) gives $J_{x_0}-2J_{x_1}=7627/134400000$.
The convolution density bound implies $|J_x''|\le\alpha_x$, so centering
at $b$ improves the error to $\alpha_x/(12N)$ and yields
$z_{x_0}-2z_{x_1}\ge1161/44800000-10^{50}/50!>0$.

For nonmonotonicity of $r$, use the same graph with $m=1,b=1/20$,
$N=100000,k=20$, and compare $x_0=1/2,x_1=3/5$.
Conditional frozen payment is
$\ell_x(t)=0$ for $t\le0$, $2t-t^2/x$ for $0<t<x$, and $x$ thereafter.
With $A=\min(\epsilon U,(x+\epsilon)V)$, direct integration gives
\[
 \E\ell_x(b-A)=(\epsilon^{-1}+(x+\epsilon)^{-1})b^2
 -\frac{2b^3}{3\epsilon(x+\epsilon)}
 -\frac{(\epsilon^{-1}+(x+\epsilon)^{-1})b^3}{3x}
 +\frac{b^4}{6\epsilon(x+\epsilon)x}.
\]
Its low-minus-high value is $23/100800$. The averaged function has
second derivative $2f_A(s)-(2/x)\Prob(s-x<A<s)$, bounded in absolute value
by $24$, so centering error per profile is at most $2/N$.
Truncation error is at most $2^{20}/20!$ per profile. Therefore
\[
 r_x(x_0)-r_x(x_1)\ge23/100800-4/100000-2^{21}/20!>0.
\]
These are finite fixed-owner examples, not limit-only counterexamples.
\end{proof}

Even a constant through-antichain value and a common disjoint minimum cut
do not imply cross-substitutes. Here is an exact separating construction.
Take two source owners $e,h$ of reports $1,h$ in parallel, followed by two
owners of report two in parallel. In parallel with that whole branch put a
100-edge path with every report ten. Follow both by $N=10000$ parallel
owners of total $21/2$. The source antichain is $\{e,h,p_1\}$.
On the joint box $e\in[.99,1]$, $h\in[.1,1]$, $p_1\in[9.9,10]$,
every upstream cut costs at least $10.99$, so the final cut fixes both
$T$ and through-antichain rank at $10.5$.
At $e=1,p_1=10$, the isolated first branch gives
\[
 q_h=\Prob(U_e+hU_h<(1+h)(U_{b_1}+U_{b_2}))
     =5/6+h/[12(1+h)^2],\quad q_1-q_{.1}=27/1936.
\]
In the full graph the additional condition is that this winning score plus
$10\min_{j\le100}U_{p_j}$ is below the final score. Failure probability
is at most $(9/10)^{100}+49/(27N)<1/4000$, by Chebyshev below three for
the final score. Thus $z_e(h=1)-z_e(h=.1)>27/1936-1/2000>0$.
The graph is TTSP and its envelope is safe; this only rejects the proposed
comparison and does not contradict Proposition~\ref{noise:antichain}.

\section{Shapley obstructions beyond two pools}\label{noise:shapley}
\begin{proposition}\label{noise:shapley-paths}
Ordinary individual Shapley value can require a coverage multiplier
$\Omega(K^{1/3})$ on unit-capacity simple TTSP graphs. Even on serial
parallel pools all of total one, no positive constant coalition-coverage
fraction is valid over all finite ownerships.
\end{proposition}
\begin{proof}
For the first statement, place a unit veto edge before two internally disjoint
unit paths of lengths $n,\ell\ge2$. There are $K=1+n\ell$ structural cuts.
Let $S$ consist of the veto and the length-$n$ path, so $v(S)=1$.
By independent uniform arrival times, the veto share is
$1/(n+1)+1/(\ell+1)-1/(n+\ell+1)$, and every outsider on the other path
has share $1/(\ell+1)-1/(n+\ell+1)$: it is pivotal exactly when the veto
and its companions have arrived but the other path is incomplete.
Efficiency gives
$\phi(S)=1/(\ell+1)+\ell/(n+\ell+1)$.
For $n=\ell^2$ this is less than $2/\ell$, proving the claimed multiplier.
For example $n=9,\ell=3$ gives $K=28$ and $\phi(S)=25/52<1/2$.

For balanced pools fix integers $M,L\ge2$, put $a=1/(4M^2)$,
$q_n=2M(n+1)-1$, and take $K_n=L^{q_n}$ pools of type $n=1,\ldots,M$.
Each has one distinguished owner of report $a$ and $n$ outsiders of report
$b_n=(1-a)/n$. All pool totals are one. Let $S$ consist of the distinguished
owners, so $v(S)=a$. We prove
\[
 \begin{gathered}
 \frac2{M+1}\le\frac{\phi(S)}a\le\frac2{M+1}+E(M,L),\\
 \begin{aligned}
 E={}&(M-1)L^{-1/[M(M+1)]}+L^{-1/(2M)}+(M-1)L^{-2M}\\
 &+\frac{M(M-1)}2(2^{M+1}+1)L^{-1/2}.
 \end{aligned}
 \end{gathered} \tag{D}
\]
Write $v(X)=\int_0^1\ind\{\text{every pool has $X$-capacity}\ge\theta\},d\theta$.
Finite Shapley summation commutes with this integral. Pool completion times
are independent continuous order statistics; the last completing pool
supplies the global pivotal owner. A distinguished type-$n$ owner matters
only on bands $(jb_n,jb_n+a)$, $j=0,\ldots,n$. On such a band the Boolean
pool game requires $j+1$ of its $n+1$ owners, so labels are symmetric;
in the gaps the distinguished owner is a dummy.

For an interior band $1\le j<n$, put $d=n-j$, $r=d+1$ and
$m=\lceil n/d\rceil-1<n$. Then $md<n\le(m+1)d$ and
$j/n-(m-1)/m\ge1/M^2$. Since $a<(1-a)/M^2$, the whole band lies in
the type-$m$ gap requiring all $m$ outsiders, where its distinguished owner
is a dummy. A type-$n$ completion is late beyond $1-\delta$ with probability
at most $2^{n+1}\delta^r$. Set
$Q=q_mr-q_n\ge M+1$ and $\beta=(q_n/r+q_m)/2$,
$\delta=L^{-\beta}$. A union over type-$n$ pools bounds the probability
any is late by $2^{n+1}L^{-Q/2}$. All type-$m$ pools complete before that
threshold with probability at most $e^{-mL^{Q/(2r)}}$.
Thus the global last pool has type $n$ with probability at most
$(2^{M+1}+1)L^{-1/2}$. There are $M(M-1)/2$ interior bands, each of
length $a$; their total contribution gives the last error in (D).

On the common lower endpoint band $(0,a)$ each pool is an OR game.
Its completion late tail is $\delta^{n+1}$. Set
$\alpha_n=q_n/(n+1)=2M-1/(n+1)$ and
$\beta=(\alpha_M+\alpha_{M-1})/2$, $\delta=L^{-\beta}$.
The probability any lower-type pool is late is at most
$(M-1)L^{-1/[M(M+1)]}$, and that no type-$M$ pool is late is at most
$e^{-L^{1/(2M)}}\le L^{-1/(2M)}$.
Conditional on the latest pool's type $n$, the pivotal label is uniform
among $n+1$ owners, independently of its completion time. Its distinguished
share is at least $1/(M+1)$ and at most that plus these errors.
On the upper endpoint band $(1-a,1)$, the entire game is unanimity; the
$S$ share is $\sum K_n/\sum(n+1)K_n$, between $1/(M+1)$ and that plus
$(M-1)L^{-2M}$. These bands each have length $a$; nonnegative interior
contributions prove (D). For each fixed $M$, let $L\to\infty$, and then
let $M\to\infty$. All witnesses are finite. Explicitly $M=4,L=10^{40}$
gives a ratio below $431/1000$, and taking
$L=[4(M+1)M(2^{M+1}+1)]^{2M(M+1)}$ gives ratio at most $3/(M+1)$.
\end{proof}
The ordinary Shapley rule remains own-report monotone; these statements
concern its coverage, not its incentives. A smaller noncore calibration is
three balanced pools $(1,1,1)$, $(1,1,1)$, $(1,2)$: the six first owners
receive $9/28$ each, and the last two $7/20,101/140$. A coalition with one
unit owner from each pool receives $139/140$ for worth one. For direct
verification, at arrival time $u$ let $p_1=1-(1-u)^3$, $p_2=3u^2-2u^3$,
$p_3=u^3$. The last unit owner's marginal integrates
$(1-u)p_1^2+u^7$, and the weight-two owner's integrates
$(1-u)(p_1^2+p_2^2)+u(p_2^2+p_3^2)$; efficiency determines the others.
Balanced noncore examples are already in the literature; no priority
claim for this calibration or the quantitative families is made.

\section{Limits of single-base lower-bound certificates}\label{noise:stars}
\begin{proposition}\label{noise:star-domain}
Fix a DAG and one base profile $c$ of value $T>0$. Let $X(c)$ consist of
$c$ and all its unilateral underreports. There are coalition-covering
payments, monotone along every unilateral comparison within $X(c)$, with
budget factor $B$ whenever
\[
 B>1,\qquad B\ge1+\frac{LB}{4(B-1)^2}.
\]
In particular $B=2+\sqrt L/2$ suffices. This is a restricted-domain
feasibility result, not a globally defined mechanism.
\end{proposition}
\begin{proof}
Write $M_e=T-F(c_{-e},0)$. Removing paths through $e$ from a single
maximum flow gives $M_e\le f_e$ simultaneously, so $\sum_eM_e\le LT$.
Fix a base minimum cut $C$ and its capacity allocation $a$, costing $T$.
It covers the base game and every underreported game.
For $e\notin C$, the own-report value is
$T_e(t)=\min(T,b_e+t)$. If $b_e\ge T$ put $h_e=0$; otherwise set
\[
 h_e=\sup_{0<t<T-b_e}
 t\left[\frac{T-BT_e(t)}{T-T_e(t)}\right]_+.
\]
Then $h_e\le M_e$. If $h_e>0$, $b_e<T/B$ and
$M_e>(1-1/B)T$, so there are at most $LB/(B-1)$ such owners.
Writing $x=T_e(t)/T$ and using $t\le Tx$ gives
\[
 h_e/T\le\max_{0<x<1/B}\frac{x(1-Bx)}{1-x}
 =(\sqrt B-\sqrt{B-1})^2\le\frac1{4(B-1)}.
\]
Thus $\sum_{e\notin C}h_e\le LTB/[4(B-1)^2]$.
Pay $a_e$ on $C$ and $h_e$ outside it at the base.

For a proper underreport by $e\in C$, the old cut now costs less than $T$
while every avoiding cut costs at least $T$. Every new minimum cut contains
$e$, so use a minimum-cut allocation paying it $t$ and costing $T_e(t)$.
For $e\notin C$ with $T_e(t)=T$, reuse the base allocation. Otherwise
choose a minimum-cut allocation $z(t)$, whose focal coordinate is $t$.
For $t>0$ pay $\alpha z(t)+(1-\alpha)a$, with
$\alpha=\min(h_e,t)/t$. It covers the lowered game and costs
$T-\alpha(T-T_e(t))\le BT_e(t)$ by the definition of $h_e$ when
$BT_e(t)<T$, and trivially otherwise. At zero use any minimum-cut allocation.
The focal payment is $\min(h_e,t)$ until the value plateaus, then $h_e$;
since $h_e\le M_e=T-b_e$, this is monotone. Different proper rays differ
in two coordinates, so no other unilateral comparisons arise.
\end{proof}

For completeness, the direct path-star certificate has an explicit bound.
For each target $e$ choose an underreport $t_e$ and $m_e\ge2$ path coalitions
through $e$, pairwise disjoint outside it, each of worth at least $\delta_e$.
Writing $T_e$ for the lowered grand value, summed coverage and budget force
at the base only
\[
 p_e(c)\ge f_e=\left[\frac{m_e\delta_e-BT_e}{m_e-1}\right]_+.
\]
With a fixed base mincut $C$, its targets contribute at most $T$.
For $e\notin C$, disjointness gives $m_e\delta_e\le T$.
A positive $f_e$ requires $T_e<T/B$, hence $M_e>(1-1/B)T$; there are
at most $LB/(B-1)$ such targets. Moreover
$f_e\le T(m_e-B)/[m_e(m_e-1)]\le T/[4(B-1)]$, by maximizing
$(u-(B-1))/u^2$ for $u=m_e-1$.
Thus $\sum_ef_e\le T[1+LB/(4(B-1)^2)]$; this specified certificate
cannot contradict budget $B$ whenever $4(B-1)^3\ge LB$.
The stronger star-domain proposition permits arbitrary coalitions.

A separate useful stopping case concerns the complete ordered DAG with
$r$ internal vertices. It has $K=2^r$ minimal cuts: each internal-vertex
subset defines a source side, and every crossing edge has a private path
of length at most three. If every internal crossarc has capacity at least
$T$, shortcut each maximum-flow path to its first and last internal
vertices. Boundary flows are preserved and each shortcut carries at most
$T$, so the resulting maximum flow has total edge flow at most $3T$.
The globally truthful marginal rule $M_e=F(c)-F(c_{-e},0)$ is coalition
covering and satisfies $M_e\le f_e$ for every maximum flow: on any original
minimum cut it pays the full cut capacities, which already cover coalitions.
Consequently $\sum_eM_e\le3T$ on all such high-crossarc profiles.
A lower-bound certificate supported only on those profiles cannot force a
factor above three. Multi-base transport with genuine internal congestion
is not excluded by either stopping statement.

\subsection{Fixed monotone statistics beyond the minimum-sum rule}
\begin{lemma}[Independent-threshold closure]\label{hist:threshold-class}
Fix a partition $N=N_1\sqcup\cdots\sqcup N_k$ with $k\ge2$ and the
weighted multi-glove game
$v_c(S)=\min_j\sum_{i\in S\cap N_j}c_i$. Write $T=v_c(N)$ and
$d_i=\min(c_i,T)$. For each pool choose a fixed finite-real-valued,
Borel-measurable, coordinatewise nondecreasing function
$\Phi_j:\R_+^{N_j}\to\R$. The functions may depend on the fixed pool
labels, ownership and $k$, but their only varying inputs are the vectors
$(d_iU_i)_{i\in N_j}$, where all $U_i$ are independent uniforms.
Let $W$ minimize these $k$ statistics, with a fixed total priority for ties.
Then the deterministic payments
\[
 b_i(c)=d_i\Prob(W=j),\qquad i\in N_j,
\]
cover every coalition and are nondecreasing in each owner's own report.
Set them to zero when $T=0$. If $M_i=T-v_c(N\setminus\{i\})$ and
$H_i$ is the original entire-core coordinate maximum, then
$\max\{b_i,M_i/k\}$ also protects $H_i/k$, preserves these properties,
and increases total payment by at most $T$.
\end{lemma}
\begin{proof}
For every set $A$, $\min(c(A),T)=\min(d(A),T)$: either a member has
capacity at least $T$, or no member is changed. A minimum-total pool
has capped total $T$, so $v_c(S)=\min_jd(S\cap N_j)$.
Since the winning probabilities sum to one, $b(S)$ is a convex combination
of these capped pool totals and therefore covers $v_c(S)$.

Fix an owner $i\in N_j$, its report $x$, the fixed remainder
$a=\sum_{h\in N_j\setminus\{i\}}c_h$, and the least other pool total
$B$. If $B=0$, every payment remains zero. Otherwise
$T(x)=\min(x+a,B)$ and $d_i(x)=\min(x,B)$.
The other capped coordinates in pool $j$ are constant: if $a\le B$,
then $T(x)\ge a\ge c_h$ for every such owner; if $a>B$, $T(x)=B$.
Every other pool's capped coordinates are nondecreasing in $x$.

Condition on $U_{-i}$. Let $E_x\subseteq\R_+$ be the set of focal
coordinate values $z$ for which pool $j$ wins, keeping its other coordinates
fixed. This is a Borel down-set because $\Phi_j$ is nondecreasing in $z$.
As $x$ increases, all competing statistics weakly increase, so $E_x$
expands, with the same conclusion at equality by the fixed tie priority.
For $d_i(x)>0$, the conditional payment is exactly the Lebesgue length
of $E_x\cap[0,d_i(x)]$. Both this set and its upper endpoint expand with
$x$, proving conditional and hence expected monotonicity. At $d_i=0$,
its length and payment are zero. No continuity of the statistics or
atom-free distribution of their values is required.

The complementary-coalition constraint gives $H_i\le M_i$, and
$M_i$ is own-report monotone because its deletion value is independent
of that report. In pool $j$, write $s=C_j-T\ge0$, where
$C_j=\sum_{i\in N_j}c_i$. Then $M_i=(c_i-s)_+$.
If $r\ge1$ of these terms are positive, their sum is at most
$C_j-rs\le T$; if none is positive, it is zero.
Thus $\sum_iM_i/k\le T$. Taking the coordinatewise maximum preserves
coverage and monotonicity and adds at most that amount. At $T=0$, the
original core and every floor are zero.
\end{proof}
This closure statement gives no budget bound for an arbitrary statistic.
The minimum-sum statistic yields the previously analyzed continuous
Banzhaf rule. The earlier class-wide growth-order obstruction is now
implied by the stronger unrestricted lower bound; no separate optimality
or novelty claim for this broader sufficient-condition lemma is needed.

\begin{thebibliography}{99}
\bibitem{MorettiNorde}
S. Moretti and H. Norde.
A note on weighted multi-glove games.
\emph{Social Choice and Welfare} 57:721--732, 2021.
\href{https://doi.org/10.1007/s00355-021-01337-8}{doi:10.1007/s00355-021-01337-8}.
\bibitem{Huang}
S. Huang, W. Mei, X. Yang and Z. Cao.
Mechanism Design in Max-Flows.
\href{https://arxiv.org/abs/2502.08248}{arXiv:2502.08248v1}, 2025.
\bibitem{Kurz}
S. Kurz. Measuring voting power in convex policy spaces.
\href{https://arxiv.org/abs/1312.5936}{arXiv:1312.5936v1}, 2013.
\bibitem{Bachrach2009}
Y. Bachrach, E. Elkind, R. Meir, D. Pasechnik, M. Zuckerman, J. Rothe,
J. S. Rosenschein.
The cost of stability in coalitional games.
\href{https://arxiv.org/abs/0907.4385v1}{arXiv:0907.4385v1}, 2009.
\bibitem{Bachrach}
Y. Bachrach, E. Elkind, E. Malizia, R. Meir, D. Pasechnik,
J. S. Rosenschein, J. Rothe and M. Zuckerman.
Bounds on the Cost of Stabilizing a Cooperative Game.
\emph{Journal of Artificial Intelligence Research} 63:987--1023, 2018.
\href{https://doi.org/10.1613/jair.1.11270}{doi:10.1613/jair.1.11270}.
\bibitem{Branzei}
R. Br\^anzei, T. Solymosi and S. Tijs.
Type monotonic allocation schemes for a class of market games.
\emph{TOP} 15:78--88, 2007.
\url{https://doi.org/10.1007/s11750-007-0002-7}.
Final author-uploaded text inspected; Definition 1 and Theorems 5--6 concern
population changes among unit owners, not own-capacity reports.
\bibitem{ABV}
I. Amigo, P. Belzarena and S. Vaton.
Revenue sharing in network utility maximization problems.
\emph{NETNOMICS} 17:255--284, 2016.
\url{https://doi.org/10.1007/s11066-016-9112-4}.
Inspected author manuscript: \url{https://hal.science/hal-01400498v1}.
\bibitem{Immorlica}
N. Immorlica, M. Mahdian and V. S. Mirrokni.
Limitations of Cross-Monotonic Cost-Sharing Schemes.
\emph{ACM Transactions on Algorithms} 4(2), Article 24, 2008.
\href{https://doi.org/10.1145/1361192.1361201}{doi:10.1145/1361192.1361201}.
\bibitem{Owen}
G. Owen. On the core of linear production games.
\emph{Mathematical Programming} 9:358--370, 1975.
\href{https://doi.org/10.1007/BF01681356}{doi:10.1007/BF01681356}.
The current source audit inspected the publisher record, not the full original.
\bibitem{SametZemel}
D. Samet and E. Zemel. On the core and dual set of linear programming games.
\emph{Mathematics of Operations Research} 9(2):309--316, 1984.
\url{https://doi.org/10.1287/moor.9.2.309}.
\href{https://www.tau.ac.il/~samet/papers/lp\%20games.pdf}{Author-hosted original}.
\bibitem{BoydSubgradients}
S. Boyd, J. Duchi, M. Pilanci and L. Vandenberghe.
\emph{Subgradients}. Stanford EE364b lecture notes, Section 2.1.
\url{https://web.stanford.edu/class/ee364b/lectures/subgradients_notes.pdf}.
\bibitem{Ghodsi}
A. Ghodsi, M. Zaharia, B. Hindman, A. Konwinski, S. Shenker and I. Stoica.
Dominant resource fairness: Fair allocation of multiple resource types.
\emph{Proceedings of NSDI}, 2011.
\url{https://www.usenix.org/conference/nsdi11/dominant-resource-fairness-fair-allocation-multiple-resource-types}.
\bibitem{Haussler1992}
D. Haussler. Decision theoretic generalizations of the PAC model for neural
net and other learning applications.
\emph{Information and Computation} 100:78--150, 1992.
Constants here refer to Corollary 1 and Theorems 6 and 12 of the January 1992 manuscript,
\url{https://www.cis.upenn.edu/~mkearns/teaching/COLT/haussler.pdf}.
\bibitem{Hoeffding}
W. Hoeffding. Probability inequalities for sums of bounded random variables.
\emph{Journal of the American Statistical Association} 58(301):13--30, 1963.
\href{https://www.cs.rpi.edu/academics/courses/spring06/random/hoefding.pdf}{Free original}.

\end{thebibliography}
\end{document}
